\chapter{概率论作为信息的实分析}\label{ch:05}

概率论并没有另造一套与分析无关的语言。事件是可测集合，分布是推前测度，期望是积分，
而条件期望则把一个随机变量压缩到较粗的信息中。真正新增的困难来自这些对象的组合方式：
单变量分布不记录两个变量怎样耦合；逐个条件事件上的比值无法处理零概率条件，也无法协调不可数
多个版本；几种随机收敛对尾部的要求各不相同；渐近收敛又不能回答有限样本究竟偏离多远。

本章依次修补这些缺口。第一节从概率空间和推前积分出发，以独立性、Borel--Cantelli 引理和
零一律说明无穷重复试验怎样把集合论估计变成路径结论。第二节从“只知道部分信息时如何预测”这一
问题得到条件期望，再在标准 Borel 状态空间上把条件平均协调成概率核，并由此建立测度分解、
$L^2$ 最佳预测和贝叶斯公式。第三节比较几乎必然、依概率、$L^p$ 与依分布收敛；尖峰和质量逃逸
迫使我们分别建立一致可积性与紧性，最终得到 Portmanteau 与 Prokhorov 定理。第四节把
独立和送入指数测试和复指数测试：前者产生集中不等式与大数律，后者通过实变量积分与高斯
平滑建立特征函数唯一性、\Levy{} 连续性定理和中心极限定理。

第 1--4 章建立的推前测度、Radon--Nikodym 定理、$L^p$ 完备性、Polish 空间紧性及近似恒等
构成本章的分析基础。这里的振荡指数 $e^{itx}$ 始终作为有界的正弦、余弦测试函数处理；
特征函数的唯一性由高斯平滑证明，而非由一般 Fourier 反演推出。

\section{概率模型、独立性与尾事件}\label{sec:05-01}

\subsection{概率空间、随机变量、分布与期望}\label{subsec:5-1-1}
\index{概率空间}\index{随机变量}\index{分布}\index{期望}

一枚骰子落定以后，点数并不再“随机”；不确定的是观察之前我们能区分哪些情形，以及给这些情形
分配多少质量。在有限模型中，把每个子集都看作事件通常不会出问题。到了无限样本空间，这个做法既
未必与测度相容，也没有反映观测能力的限制。概率论因此不从“随机的数”出发，而从一族允许提问的
事件出发。

\begin{definition}[概率空间]\label{def:5-1-1-1}
概率空间是三元组 $(\Omega,\mathcal F,\Prob)$，其中 $\mathcal F$ 是 $\Omega$ 上的
$\sigma$-代数，$\Prob$ 是 $(\Omega,\mathcal F)$ 上满足 $\Prob(\Omega)=1$ 的测度。
$\mathcal F$ 的元素称为事件；事件 $A$ 的概率是 $\Prob(A)$。
\end{definition}

概率为一的事件不必等于 $\Omega$，概率为零的事件也不必为空。测度论中的 ``a.e.'' 在概率空间上
称为\emph{几乎必然}（almost surely，简写 a.s.）。这只是语言改换：异常点仍须包含在某个可测零集
中，所有积分定理的集合论含义都没有改变。

Kolmogorov 的公理化把概率纳入测度论，却没有替建模者决定 $\Omega$ 和 $\mathcal F$。频率模型把
概率与重复试验的长期比例联系，几何模型把概率与归一化体积联系，信息模型则把 $\mathcal F$ 看成
能够分辨的提问集合；三种解释可以导向同一个概率空间，但都不是测度公理的逻辑后果。

掷骰子的结果是样本点，点数是样本点的函数。这个区别在更复杂的模型里尤其重要：同一个样本点
可以同时产生位置、速度、首次到达时间等许多观测量。

\begin{definition}[随机元、随机变量与分布]\label{def:5-1-1-2}
设 $(\Omega,\mathcal F,\Prob)$ 为概率空间，$(S,\mathcal S)$ 为可测空间。可测映射
\[
  X:(\Omega,\mathcal F)\longrightarrow(S,\mathcal S)
\]
称为取值于 $S$ 的\emph{随机元}；当 $S=\R$ 且 $\mathcal S=\Borel(\R)$ 时，称为
\emph{随机变量}。随机元 $X$ 的分布是推前概率测度
\[
  \Law(X)=X_*\Prob,
  \qquad \Law(X)(B)=\Prob(X\in B)\quad(B\in\mathcal S).
\]
若 $\Law(X)=\Law(Y)$，则称 $X,Y$ 同分布，记作 $X\overset{d}=Y$。
\end{definition}

分布把样本点的名字全部忘掉，只保存状态空间中每个可测问题的答案。这种遗忘非常有用，但也有
代价：它不会记录两个变量在同一样本点上如何相互配合。这个缺口将在下一小节成为独立性的入口。

\begin{remark}[一般可测空间与标准 Borel 空间的分工]
\label{rem:5-1-1-standard-borel}
分布、推前积分和独立性的定义只需要可测空间。正则条件分布、测度分解以及可数个版本的同时选择则会
要求状态空间为标准 Borel 空间，即与某个 Polish 空间的 Borel 可测空间可测同构。
$\R^d$、可数离散空间以及带上确界范数的 $C([0,1])$ 都给出标准 Borel 空间。这个条件不是
随机元定义的一部分；它是后文存在性定理所支付的正则性成本，\S\ref{subsec:5-2-2} 将解释其作用。
\end{remark}

\begin{definition}[期望与矩]\label{def:5-1-1-3}
若 $X$ 是非负扩展实值随机变量，定义
\[
  \E X=\int_\Omega X\dd\Prob\in[0,\infty].
\]
若 $X$ 为实值或复值且可积，仍定义 $\E X=\int_\Omega X\dd\Prob$。更一般地，若
$X$ 取值于 $(S,\mathcal S)$，$f:S\to\R$ 或 $\C$ 可测，并且 $f\circ X$ 非负或可积，则
\[
  \E f(X)=\int_\Omega f(X(\omega))\dd\Prob(\omega)
  =\int_S f(x)\dd\Law(X)(x).
\]
对实值 $X$，$\E|X|^p$ 称为 $p$ 阶绝对矩。若 $X\in L^2(\Prob)$，定义
\[
  \Var(X)=\E|X-\E X|^2.
\]
\end{definition}

最后一个等号是非负推前积分公式（命题~\ref{proposition:3-1-1-5}）及其可积版本
（推论~\ref{corollary:3-1-3-pushforward}）。它给出同一个期望的两种算法：在样本空间上积分，
或先求分布，再在状态空间上积分。

\begin{figure}[htbp]
\centering
\begin{tikzpicture}[node distance=13mm and 19mm]
  \node[book strong node] (omega) {样本空间\\$(\Omega,\mathcal F,\Prob)$};
  \node[book strong node,right=of omega] (state) {状态空间\\$(S,\mathcal S,\Law(X))$};
  \node[book warm node,below=of omega] (sampleint)
    {$\displaystyle\int_\Omega f\!\circ X\,\dd\Prob$};
  \node[book warm node,below=of state] (lawint)
    {$\displaystyle\int_S f\,\dd\Law(X)$};
  \draw[book accent arrow] (omega)--node[above,book label]{$X$}(state);
  \draw[book arrow] (omega)--node[left,book label]{观测后积分}(sampleint);
  \draw[book arrow] (state)--node[right,book label]{按分布积分}(lawint);
  \draw[book dashed arrow] (sampleint.south)--++(0,-9mm)-|(lawint.south);
  \node[book label,above=1mm] at ($(sampleint.south)!0.5!(lawint.south)+(0,-9mm)$)
    {推前积分公式};
\end{tikzpicture}
\caption{随机变量推前概率测度；期望可在箭头的任一端计算。}
\label{fig:5-1-1-pushforward-law}
\end{figure}

\begin{proposition}[期望的分布不变性]\label{proposition:5-1-1-1}
设 $X,Y$ 取值于同一可测空间 $(S,\mathcal S)$，并且 $X\overset{d}=Y$。若
$f:S\to\R$ 或 $\C$ 可测，且 $f(X),f(Y)$ 非负或可积，则
\[
  \E f(X)=\E f(Y).
\]
\end{proposition}

\begin{proof}
由推前积分公式和 $\Law(X)=\Law(Y)$，
\[
 \E f(X)=\int_S f\dd\Law(X)
 =\int_S f\dd\Law(Y)=\E f(Y).
\]
非负情形的等式允许两端同为 $+\infty$；可积情形中各积分是有限标量，因而没有
$\infty-\infty$ 的问题。
\end{proof}

\begin{example}[有限样本空间]\label{ex:5-1-1-1}
令 $\Omega=\{1,\ldots,6\}$、$\mathcal F=2^\Omega$，并令每个单点的概率为 $1/6$。
有限样本空间上的每个实值或复值函数都可测。取点数变量 $X(\omega)=\omega$。对任意函数
$g:\Omega\to\C$，
\[
 \E g(X)=\frac16\sum_{k=1}^6g(k).
\]
特别地，
\[
 \E X=\frac{1+\cdots+6}{6}=\frac72,
 \qquad
 \E X^2=\frac{1^2+\cdots+6^2}{6}=\frac{91}{6},
\]
所以
\[
 \Var(X)=\E X^2-(\E X)^2
 =\frac{91}{6}-\frac{49}{4}=\frac{35}{12}.
\]
其中方差恒等式由逐项展开得到：
\begin{align*}
 \E|X-\E X|^2
 &=\E\bigl[X^2-2X\E X+(\E X)^2\bigr]\\
 &=\E X^2-2(\E X)^2+(\E X)^2
 =\E X^2-(\E X)^2.
\end{align*}
这里期望确实退化为有限加权和；定义的价值在于同一公式不必在连续模型中另起炉灶。
\end{example}

\begin{example}[均匀变量生成指数变量]\label{ex:5-1-1-2}
设 $U$ 在 $(0,1)$ 上具有 Lebesgue 均匀分布，$\lambda>0$，并令
\[
 X=-\lambda^{-1}\log U.
\]
当 $x\ge0$ 时，
\[
 \Prob(X>x)=\Prob(U<e^{-\lambda x})=e^{-\lambda x},
\]
而 $x<0$ 时该概率为一。因此 $X$ 的分布函数和密度分别为
\[
 F_X(x)=
 \begin{cases}0,&x<0,\\1-e^{-\lambda x},&x\ge0,\end{cases}
 \qquad
 p_X(x)=\lambda e^{-\lambda x}\1_{(0,\infty)}(x).
\]
利用分布积分公式，对 $R>0$ 在 $[0,R]$ 上分部积分，有
\begin{align*}
 \int_0^R x\lambda e^{-\lambda x}\dd x
 &=\bigl[-xe^{-\lambda x}\bigr]_{0}^{R}
   +\int_0^R e^{-\lambda x}\dd x\\
 &=-Re^{-\lambda R}+\frac{1-e^{-\lambda R}}{\lambda}.
\end{align*}
由 $Re^{-\lambda R}\to0$ 和 $e^{-\lambda R}\to0$，非负积分的单调收敛定理给出
\[
 \E X=\lim_{R\to\infty}\int_0^R x\lambda e^{-\lambda x}\dd x=\frac1\lambda.
\]
这项计算只使用推前测度、分布函数和一元换元；样本空间在得到 $F_X$ 后便退出了舞台。
\end{example}

\begin{example}[一维高斯分布]\label{ex:5-1-1-gaussian-law}
记 $G\sim N(m,\sigma^2)$（$m\in\R,\sigma>0$）表示 $G$ 的分布具有密度
\[
 q_{m,\sigma}(x)=\frac1{\sqrt{2\pi\sigma^2}}
 \exp\!\left(-\frac{(x-m)^2}{2\sigma^2}\right).
\]
Analysis II 中的高斯积分
$\int_\R e^{-x^2/2}\dd x=\sqrt{2\pi}$ 保证该密度的积分为一。若
$Z\sim N(0,1)$，则对每个 Borel 集 $A$，以 $x=m+\sigma z$ 换元得到
\begin{align*}
 \Prob(m+\sigma Z\in A)
 &=\int_\R\1_A(m+\sigma z)\frac{e^{-z^2/2}}{\sqrt{2\pi}}\dd z\\
 &=\int_A q_{m,\sigma}(x)\dd x.
\end{align*}
故 $m+\sigma Z\sim N(m,\sigma^2)$。由对称区间上奇函数积分为零以及
$|z|e^{-z^2/2}$ 可积，有 $\E Z=0$。此外，对 $R>0$，利用
$(e^{-z^2/2})'=-ze^{-z^2/2}$ 分部积分得
\begin{align*}
 \int_{-R}^{R}z^2e^{-z^2/2}\dd z
 &=-\int_{-R}^{R}z\,\dd\bigl(e^{-z^2/2}\bigr)\\
 &=-\bigl[ze^{-z^2/2}\bigr]_{-R}^{R}
   +\int_{-R}^{R}e^{-z^2/2}\dd z\\
 &=-2Re^{-R^2/2}+\int_{-R}^{R}e^{-z^2/2}\dd z.
\end{align*}
令 $R\to\infty$，由 $Re^{-R^2/2}\to0$ 和单调收敛定理，
\[
 \E Z^2=\frac1{\sqrt{2\pi}}\int_\R z^2e^{-z^2/2}\dd z
 =\frac1{\sqrt{2\pi}}\int_\R e^{-z^2/2}\dd z=1.
\]
因此，若 $G=m+\sigma Z$，则
\begin{align*}
 \E G&=m+\sigma\E Z=m,\\
 \Var(G)&=\E|G-\E G|^2
 =\E|\sigma Z|^2=\sigma^2\E Z^2=\sigma^2.
\end{align*}
这里固定了一维高斯（Gaussian）分布的记号；联合高斯分布、矩母函数和特征函数将在本章后续各节中分别讨论。
\end{example}

\begin{example}[同分布不同耦合]\label{ex:5-1-1-3}
在 $\Omega=\{-1,1\}$ 上取均匀概率，并令 $X(\omega)=\omega$。分别令
$Y_1=X$ 与 $Y_2=-X$。由于 $X$ 的分布关于原点对称，
\[
 \Law(Y_1)=\Law(Y_2)=\Law(X)=\tfrac12\delta_{-1}+\tfrac12\delta_1.
\]
然而
\[
 \Prob(X=Y_1)=1,
 \qquad
 \Prob(X=Y_2)=0.
\]
三个单变量分布完全相同，两个联合关系却走向相反极端。单变量分布不负责保存耦合信息。
\end{example}

积分的单调性可以把一个事件的示性函数压在可积量之下。这条一句话的观察产生了概率论中最常用的
尾估计之一。

\begin{theorem}[Markov 不等式]\label{theorem:5-1-1-2}
若 $X\ge0$ a.s. 且 $a>0$，则
\[
 \Prob(X\ge a)\le\frac{\E X}{a},
\]
其中允许 $\E X=+\infty$。更一般地，若 $X$ 实值，
$\phi:\R\to[0,\infty]$ 可测且非降，并且 $0<\phi(a)<\infty$，则
\[
 \Prob(X\ge a)\le\frac{\E\phi(X)}{\phi(a)}.
\]
\end{theorem}

\begin{proof}
在每个样本点上有
\[
 a\1_{\{X\ge a\}}\le X
\]
（先在使 $X\ge0$ 成立的共同满概率集合上比较，零集上的取值不影响积分）。非负积分的单调性给出
\[
 a\Prob(X\ge a)=\E\bigl[a\1_{\{X\ge a\}}\bigr]\le\E X.
\]
两端除以 $a$ 即得第一式。对一般的 $\phi$，当 $X(\omega)\ge a$ 时，单调性给出
$\phi(X(\omega))\ge\phi(a)$，故
\[
 \phi(a)\1_{\{X\ge a\}}\le\phi(X).
\]
再取非负积分并除以有限正数 $\phi(a)$，得到第二式。
\end{proof}

例如令 $\phi(x)=e^{tx}$，其中 $t>0$，便有
\begin{equation}\label{eq:5-1-1-exponential-markov}
 \Prob(X\ge a)\le e^{-ta}\E e^{tX}.
\end{equation}
右端若有限，还可以对 $t$ 优化。这是 Chernoff 方法的起点；本节只使用
\eqref{eq:5-1-1-exponential-markov} 这条由单调比较直接得到的公式。

\begin{corollary}[Chebyshev 不等式]\label{corollary:5-1-1-3}
若 $X\in L^2(\Prob)$，则对每个 $t>0$，
\[
 \Prob\bigl(|X-\E X|\ge t\bigr)
 \le\frac{\Var(X)}{t^2}.
\]
\end{corollary}

\begin{proof}
概率空间总质量为一，Cauchy--Schwarz 不等式给出
$\E|X|\le(\E|X|^2)^{1/2}$，所以 $\E X$ 有定义且
$|X-\E X|^2$ 可积。对非负变量 $|X-\E X|^2$ 以阈值 $t^2$ 应用
定理~\ref{theorem:5-1-1-2}，得到
\[
 \Prob(|X-\E X|\ge t)
 =\Prob(|X-\E X|^2\ge t^2)
 \le\frac{\E|X-\E X|^2}{t^2}.
\]
分子就是 $\Var(X)$。
\end{proof}

概率空间仍是测度空间，所以极限定理不需要换一套证明。需要换的只是读法：依测度收敛在概率空间上
称为\emph{依概率收敛}。

\begin{proposition}[期望继承积分的极限定理]\label{proposition:5-1-1-4}
在同一概率空间上有以下结论。
\begin{enumerate}[label=\textup{(\roman*)}]
  \item 若 $0\le X_n\uparrow X$ a.s.，则 $\E X_n\uparrow\E X$。
  \item 若 $X_n\ge0$，则
  $\E(\liminf_nX_n)\le\liminf_n\E X_n$。
  \item 若 $X_n\to X$ a.s. 且 $|X_n|\le Y$ a.s. 对所有 $n$ 成立，其中
  $Y\in L^1$，则 $\E|X_n-X|\to0$，从而 $\E X_n\to\E X$。
  \item 若每个 $X_n\in L^1$ 且 $X_n\to X$ 依概率，则
  $\E|X_n-X|\to0$ 当且仅当族 $\{X_n:n\ge1\}$ 一致可积；在这一情形中
  $X\in L^1$。
\end{enumerate}
\end{proposition}

\begin{proof}
从每个 a.s. 陈述中删去容纳全部异常的共同零集，不改变任一积分。
此后 \textup{(i)} 是单调收敛定理~\ref{theorem:3-1-2-1}，
\textup{(ii)} 是 Fatou 引理~\ref{theorem:3-1-2-2}，
\textup{(iii)} 是实／复控制收敛定理~\ref{theorem:3-1-3-5}。
概率空间满足 $\Prob(\Omega)=1<\infty$，且这里的依概率收敛正是
定义~\ref{def:3-4-1-1} 的依测度收敛。因此 \textup{(iv)} 是
Vitali 定理~\ref{theorem:3-4-1-2} 在总质量一情形的直接应用；该定理的正向部分同时给出
$X\in L^1$。每一步的假设都与相应积分定理逐项一致。
\end{proof}

\begin{center}
\begin{tabular}{@{}ll@{}}
\toprule
分析语言 & 概率语言\\
\midrule
$\mu$-a.e. & 几乎必然（a.s.）\\
依测度收敛 & 依概率收敛\\
$\int f\dd\mu$ & $\E f$\\
MCT／Fatou／DCT & 单调、下极限、支配条件下交换期望与极限\\
Vitali 定理 & 依概率收敛加一致可积推出 $L^1$ 收敛\\
\bottomrule
\end{tabular}
\end{center}

只有 a.s. 收敛并不足以交换期望。令 $\Omega=(0,1)$，取 Lebesgue 概率测度并定义
\[
 X_n=n\1_{(0,1/n)}.
\]
对每个 $\omega\in(0,1)$，当 $n>1/\omega$ 时 $X_n(\omega)=0$，故 $X_n\to0$ a.s.；但
\[
 \E X_n=n\,m(0,1/n)=1
\]
对所有 $n$ 成立。这个族也不一致可积：给定 $K>0$，取整数 $n>K$，则
\[
 \E\bigl[|X_n|\1_{\{|X_n|>K\}}\bigr]=1.
\]
失败的不是极限计算，而是缺少支配或一致尾控制。

期望的线性同样继承自积分。若 $X,Y\in L^1$，则
$\E(aX+bY)=a\E X+b\E Y$。对非负变量也可以作非负线性组合并允许 $+\infty$；但若正部、负部
都具有无穷期望，写下 ``$\infty-\infty$'' 并不会产生一个新的概率定理，只会产生一个未定义的式子。

\begin{exercise}
\begin{enumerate}
  \item 设 $F$ 是实直线上的分布函数，$U\sim\mathrm{Unif}(0,1)$。定义
  $F^{\leftarrow}(u)=\inf\{x:F(x)\ge u\}$，证明
  $F^{\leftarrow}(U)$ 的分布函数为 $F$。
  \item 不引用命题~\ref{proposition:5-1-1-1} 的结论，只用推前积分公式重新证明它。
  \item 从 Markov 不等式推出 Chebyshev 不等式，并在骰子变量上对 $t=3/2$ 比较上界与真实概率。
  \item 在同一个两点概率空间上构造三对边缘分布相同而
  $\Prob(X=Y)$ 分别为 $0,1/2,1$ 的耦合；若两点空间不够，说明最小的有限扩张。
\end{enumerate}
\end{exercise}

单个分布已经足以计算单变量期望，却不能回答两个变量是否同涨同落。要描述这种关系，必须回到
联合分布，或等价地，回到各变量所携带的 $\sigma$-代数。


\subsection{独立性、乘积结构与生成类判别}\label{subsec:5-1-2}
\index{独立性}\index{相互独立}\index{两两独立}

考虑两个边缘都均匀的骰子点数。仅凭这两条边缘信息，第二枚骰子可能复制第一枚，也可能与第一枚
毫无关联。所谓独立，不是“两个变量看起来没有关系”，而是一个精确的乘积测度陈述。

\begin{definition}[独立 $\sigma$-代数]\label{def:5-1-2-1}
设 $(\mathcal G_i)_{i\in I}$ 是 $\mathcal F$ 的子 $\sigma$-代数族。若对任意有限个互异指标
$i_1,\ldots,i_k$ 以及任意 $A_j\in\mathcal G_{i_j}$，都有
\[
 \Prob\!\left(\bigcap_{j=1}^k A_j\right)
 =\prod_{j=1}^k\Prob(A_j),
\]
则称 $(\mathcal G_i)_{i\in I}$ 相互独立。
\end{definition}

事件族 $(A_i)$ 的独立性是指 $\bigl(\sigma(A_i)\bigr)$ 独立。这里
$\sigma(A)=\{\varnothing,A,A^c,\Omega\}$。因此事件独立不仅给出原事件交的乘积公式，也自动包含
任意一部分事件换成补事件后的公式。

\begin{definition}[随机元的独立性]\label{def:5-1-2-2}
随机元族 $(X_i)_{i\in I}$ 称为独立的，如果子 $\sigma$-代数族
$(\sigma(X_i))_{i\in I}$ 相互独立，其中
\[
 \sigma(X_i)=\{X_i^{-1}(B):B\text{ 属于 }X_i\text{ 的状态空间 }\sigma\text{-代数}\}.
\]
\end{definition}

\begin{definition}[两两独立与相互独立]\label{def:5-1-2-3}
若族中每两个不同成员独立，则称该族\emph{两两独立}。若它满足
定义~\ref{def:5-1-2-1} 中每个有限子族的乘积公式，则称为\emph{相互独立}。
相互独立蕴含两两独立，反向一般不成立。
\end{definition}

检查独立性的定义似乎要求遍历多个 $\sigma$-代数中的所有事件。前文建立的
$\pi$--$\lambda$ 方法可以把检验缩减到生成类，但“二元成立”不能未经论证便升级为任意有限元成立。

\begin{theorem}[生成类独立判别]\label{theorem:5-1-2-1}
对每个 $i\in I$，设 $\mathcal P_i$ 是生成 $\mathcal G_i$ 的 $\pi$-系统。假设任取有限个互异
指标 $i_1,\ldots,i_k$ 和 $A_j\in\mathcal P_{i_j}$，均有
\[
 \Prob\!\left(\bigcap_{j=1}^kA_j\right)=\prod_{j=1}^k\Prob(A_j).
\]
则 $(\mathcal G_i)_{i\in I}$ 相互独立。
\end{theorem}

\begin{proof}
若某个 $\mathcal P_i$ 不含 $\Omega$，把 $\Omega$ 添入其中。新族仍是 $\pi$-系统，生成同一
$\sigma$-代数，而含有 $\Omega$ 的乘积公式由原假设中删去该坐标得到。因此以下设
$\Omega\in\mathcal P_i$。

先看两个指标。固定 $A\in\mathcal P_{i_1}$，令
\[
 \mathcal D_A=
 \{B\in\mathcal G_{i_2}:\Prob(A\cap B)=\Prob(A)\Prob(B)\}.
\]
$\Omega\in\mathcal D_A$。若 $B\in\mathcal D_A$，则
\[
 \Prob(A\cap B^c)
 =\Prob(A)-\Prob(A\cap B)
 =\Prob(A)\Prob(B^c),
\]
故 $B^c\in\mathcal D_A$；若 $(B_n)$ 两两不交且都属于 $\mathcal D_A$，则测度的可数可加性给出
\[
 \Prob\!\left(A\cap\bigcup_nB_n\right)
 =\sum_n\Prob(A\cap B_n)
 =\Prob(A)\sum_n\Prob(B_n).
\]
所以 $\mathcal D_A$ 是 $\lambda$-系统。它包含 $\mathcal P_{i_2}$，由
$\pi$--$\lambda$ 定理包含 $\mathcal G_{i_2}$。

现在固定 $B\in\mathcal G_{i_2}$，把满足同一乘积公式的 $A\in\mathcal G_{i_1}$ 组成
$\lambda$-系统。上一段说明它包含 $\mathcal P_{i_1}$，再次应用 $\pi$--$\lambda$ 定理，得到
两个完整 $\sigma$-代数独立。

对一般有限个指标，对坐标数 $k$ 归纳。假设少于 $k$ 个坐标时，乘积公式已对相应的完整
$\sigma$-代数成立。固定 $A_j\in\mathcal P_{i_j}$（$1\le j<k$），定义
\[
 \mathcal D_k=\left\{B\in\mathcal G_{i_k}:
 \Prob\!\left(\bigcap_{j=1}^{k-1}A_j\cap B\right)
 =\left(\prod_{j=1}^{k-1}\Prob(A_j)\right)\Prob(B)\right\}.
\]
由归纳假设，$\Omega\in\mathcal D_k$。若 $B\in\mathcal D_k$，令
$C=\bigcap_{j=1}^{k-1}A_j$，则
\begin{align*}
 \Prob(C\cap B^c)
 &=\Prob(C)-\Prob(C\cap B)\\
 &=\prod_{j<k}\Prob(A_j)
   -\left(\prod_{j<k}\Prob(A_j)\right)\Prob(B)\\
 &=\left(\prod_{j<k}\Prob(A_j)\right)\Prob(B^c),
\end{align*}
故 $B^c\in\mathcal D_k$。对两两不交的 $(B_n)\subset\mathcal D_k$，可数可加性给出
\begin{align*}
 \Prob\!\left(C\cap\bigcup_nB_n\right)
 &=\sum_n\Prob(C\cap B_n)
 =\left(\prod_{j<k}\Prob(A_j)\right)\sum_n\Prob(B_n)\\
 &=\left(\prod_{j<k}\Prob(A_j)\right)\Prob\!\left(\bigcup_nB_n\right).
\end{align*}
因此 $\mathcal D_k$ 是包含 $\mathcal P_{i_k}$ 的 $\lambda$-系统，从而
$\mathcal D_k=\mathcal G_{i_k}$。

现在固定一般的 $A_k\in\mathcal G_{i_k}$ 和
$A_j\in\mathcal P_{i_j}$（$1\le j<k-1$），把满足同一 $k$ 元乘积公式的
$B\in\mathcal G_{i_{k-1}}$ 组成 $\mathcal D_{k-1}$。上述三个计算逐字适用：
$\Omega\in\mathcal D_{k-1}$ 由少一个坐标的归纳假设保证，而
$\mathcal P_{i_{k-1}}\subset\mathcal D_{k-1}$ 由上一段已经把第 $k$ 坐标推广到
$\mathcal G_{i_k}$ 保证。故 $\mathcal D_{k-1}=\mathcal G_{i_{k-1}}$。按
$k-2,k-3,\ldots,1$ 重复此步，最终得到任意
$A_j\in\mathcal G_{i_j}$（$1\le j\le k$）的乘积公式。归纳完成。
\end{proof}

\begin{theorem}[联合分布判别]\label{theorem:5-1-2-2}
设随机元 $X_i$ 取值于 $(S_i,\mathcal S_i)$，$1\le i\le n$。则
$X_1,\ldots,X_n$ 独立，当且仅当
\[
 \Law(X_1,\ldots,X_n)=\bigotimes_{i=1}^n\Law(X_i)
\]
作为 $(\prod_iS_i,\bigotimes_i\mathcal S_i)$ 上的概率测度成立。
\end{theorem}

\begin{proof}
对可测矩形 $R=A_1\times\cdots\times A_n$，联合分布的定义给出
\[
 \Law(X_1,\ldots,X_n)(R)
 =\Prob\!\left(\bigcap_{i=1}^n\{X_i\in A_i\}\right).
\]
若 $X_i$ 独立，右端等于 $\prod_i\Prob(X_i\in A_i)$，也就是
$(\bigotimes_i\Law(X_i))(R)$。两个概率测度在生成乘积 $\sigma$-代数的矩形
$\pi$-系统上一致；有限测度唯一性给出它们在整个乘积 $\sigma$-代数上一致。

反之，若两测度相等，则上述矩形计算对所有 $A_i\in\mathcal S_i$ 给出交概率因子化。
而 $\sigma(X_i)$ 的每个事件都形如 $X_i^{-1}(A_i)$，所以
$\sigma(X_1),\ldots,\sigma(X_n)$ 独立。
\end{proof}

\begin{example}[乘积坐标]\label{ex:5-1-2-1}
设每个 $\mu_i$ 都是 $(S_i,\mathcal S_i)$ 上的概率测度。在
\[
 \left(\prod_{i=1}^nS_i,\ \bigotimes_{i=1}^n\mathcal S_i,\
 \bigotimes_{i=1}^n\mu_i\right)
\]
上令 $\pi_i(x_1,\ldots,x_n)=x_i$。对 $A_i\in\mathcal S_i$，
\[
 \Prob(\pi_1\in A_1,\ldots,\pi_n\in A_n)
 =\prod_{i=1}^n\mu_i(A_i).
\]
因此坐标投影独立且 $\Law(\pi_i)=\mu_i$。乘积测度不是独立性的一个方便比喻；它正是独立联合
分布的标准实现。
\end{example}

\begin{figure}[htbp]
\centering
\begin{tikzpicture}[x=1.0cm,y=0.82cm]
  \draw[book line] (0,0) rectangle (7,4.5);
  \foreach \x in {1.1,2.8,4.2,6.1}{\draw[BookInk!45] (\x,0)--(\x,4.5);}
  \foreach \y in {0.9,2.0,3.4}{\draw[BookInk!45] (0,\y)--(7,\y);}
  \fill[BookAccent!12] (2.8,2.0) rectangle (6.1,3.4);
  \draw[BookAccent,semithick] (2.8,2.0) rectangle (6.1,3.4);
  \draw[decorate,decoration={brace,amplitude=4pt},BookWarm]
    (2.8,4.75)--(6.1,4.75) node[midway,above=5pt,book label]{$\Law(X)(A)$};
  \draw[decorate,decoration={brace,amplitude=4pt},BookWarm]
    (7.25,2.0)--(7.25,3.4) node[midway,right=5pt,book label]{$\Law(Y)(B)$};
  \node[book label] at (4.45,2.7) {$\Law(X,Y)(A\times B)$};
  \node[book label,below] at (3.5,-0.15) {$X$ 坐标};
  \node[book label,rotate=90] at (-0.35,2.25) {$Y$ 坐标};
\end{tikzpicture}
\caption{独立时，每个可测矩形的联合质量等于两条边缘质量的乘积；矩形生成整个联合分布。}
\label{fig:5-1-2-rectangle-law}
\end{figure}

联合分布的乘积结构不仅控制事件，也控制可测函数。

\begin{theorem}[期望因子化]\label{theorem:5-1-2-3}
设 $X_1,\ldots,X_n$ 独立，$X_i$ 取值于 $(S_i,\mathcal S_i)$，且
$f_i:S_i\to\R$ 或 $\C$ 可测。
\begin{enumerate}[label=\textup{(\roman*)}]
  \item 若每个 $f_i\ge0$，则
  \[
    \E\prod_{i=1}^nf_i(X_i)=\prod_{i=1}^n\E f_i(X_i),
  \]
  其中若某因子的期望为零，则右端约定为零；否则按通常的扩展非负数乘法解释。
  \item 若每个 $f_i(X_i)\in L^1$，则 $\prod_i f_i(X_i)\in L^1$，并且同一等式在
  $\R$ 或 $\C$ 中成立。
\end{enumerate}
\end{theorem}

\begin{proof}
先设 $f_i=\1_{A_i}$。此时所需等式正是
定义~\ref{def:5-1-2-1} 的有限交公式。对非负简单函数
$f_i=\sum_{j=1}^{m_i}a_{ij}\1_{A_{ij}}$，展开有限乘积，并对每一项应用示性函数情形，得到
\begin{align*}
 \E\prod_{i=1}^nf_i(X_i)
 &=\sum_{j_1,\ldots,j_n}
   \left(\prod_i a_{ij_i}\right)
   \Prob\!\left(\bigcap_i\{X_i\in A_{ij_i}\}\right)\\
 &=\sum_{j_1,\ldots,j_n}
   \prod_i\bigl(a_{ij_i}\Prob(X_i\in A_{ij_i})\bigr)
 =\prod_{i=1}^n\E f_i(X_i).
\end{align*}

对一般非负 $f_i$，取第三章的规范非负简单函数 $s_{i,k}\uparrow f_i$。因为因子数有限，
$\prod_i s_{i,k}(X_i)\uparrow\prod_i f_i(X_i)$。单调收敛定理给出
\[
 \E\prod_i f_i(X_i)
 =\lim_{k\to\infty}\prod_i\E s_{i,k}(X_i)
 =\prod_i\E f_i(X_i).
\]
若某个极限因子为零，则相应非负积分为零，命题~\ref{proposition:3-1-1-3} 说明左端也为零；
若没有零因子，有限个递增扩展实数乘积的极限就是右端。这个简单函数—单调极限过程是函数单调类
推广在当前乘积函数上的具体实现。

若每个 $f_i(X_i)$ 可积，把已证的非负版本用于 $|f_i|$，得到
\[
 \E\prod_i|f_i(X_i)|=\prod_i\E|f_i(X_i)|<\infty.
\]
所以乘积可积。若 $f_i$ 为复值，则逐点有
\[
 f_i=(\Re f_i)^+-(\Re f_i)^-
   +\mathrm i\bigl((\Im f_i)^+-(\Im f_i)^-\bigr),
\]
并且右端四个非负函数在 $X_i$ 下均可积；实值情形只保留前两项。把
$\prod_i f_i(X_i)$ 按此表示式展开为有限和，对每一项使用非负因子化，再重新合并，便得到所述
实值或复值等式；所有项都绝对可积，重排不产生条件收敛问题。
\end{proof}

第四章的测度卷积在这里得到直接的概率读法。若独立随机向量 $X,Y$ 的分布分别为概率测度
$\mu,\nu$，则对每个非负 Borel 函数 $h$，定理~\ref{theorem:5-1-2-2}、Tonelli 定理和
测度卷积的定义给出
\[
 \E h(X+Y)
 =\iint h(x+y)\dd\mu(x)\dd\nu(y)
 =\int h\dd(\mu*\nu).
\]
非负测试函数积分唯一确定测度（取 $h=\1_B$ 即可比较任意 Borel 集 $B$），所以
\begin{equation}\label{eq:5-1-2-independent-sum-convolution}
 \Law(X+Y)=\Law(X)*\Law(Y).
\end{equation}
若 $\mu=fm$、$\nu=gm$，第四章的卷积计算进一步给出
$\Law(X+Y)=(f*g)m$。独立性在此承担的准确工作，是把联合分布变成乘积分布；加法映射随后把
这个乘积推前为卷积。

\begin{proposition}[独立性的可测变换稳定性]\label{proposition:5-1-2-4}
若 $(X_i)_{i\in I}$ 独立，且每个 $f_i$ 是定义在 $X_i$ 状态空间上的可测映射，则
$(f_i(X_i))_{i\in I}$ 独立。
\end{proposition}

\begin{proof}
对每个 $i$，复合映射的可测性给出
\[
 \sigma(f_i(X_i))\subset\sigma(X_i).
\]
从一族独立 $\sigma$-代数中各取子 $\sigma$-代数，定义~\ref{def:5-1-2-1} 的每个有限交
乘积公式仍成立。因此这些复合随机元独立。
\end{proof}

两两独立为何不够，可以在四个样本点上算清。

\begin{example}[两两独立但不相互独立]\label{ex:5-1-2-2}
令 $\Omega=\{-1,1\}^2$ 并给四点各概率 $1/4$。定义
\[
 X(\omega_1,\omega_2)=\omega_1,
 \qquad Y(\omega_1,\omega_2)=\omega_2,
 \qquad Z=XY.
\]
三者都具有 Rademacher 分布
$\tfrac12\delta_{-1}+\tfrac12\delta_1$。映射
$(\omega_1,\omega_2)\mapsto(X,Z)=(\omega_1,\omega_1\omega_2)$ 是
$\{-1,1\}^2$ 上的双射，因此 $(X,Z)$ 在四个有序对上均匀；同理 $(Y,Z)$ 均匀，而
$(X,Y)$ 按构造均匀。故三对都独立。

另一方面，$XYZ=1$ 在每个样本点上成立，并且
\[
 \Prob(X=1,Y=1,Z=1)=\frac14
 \ne\frac12\cdot\frac12\cdot\frac12.
\]
因此 $X,Y,Z$ 不相互独立。缺失的信息恰好是三元交，而不是某一对被漏算。
\end{example}

独立性也比“协方差为零”强。对实值二阶可积变量，暂记
\[
 \Cov(X,Y)=\E[(X-\E X)(Y-\E Y)].
\]

\begin{example}[不相关不等于独立]\label{ex:5-1-2-3}
设 $U\sim\mathrm{Unif}(-1,1)$，令 $X=U$、$Y=U^2$。均匀分布的密度为
$\frac12\1_{(-1,1)}$，因而
\begin{align*}
 \E U&=\frac12\int_{-1}^{1}u\dd u=0,\\
 \E U^2&=\frac12\int_{-1}^{1}u^2\dd u
 =\frac12\left[\frac{u^3}{3}\right]_{-1}^{1}=\frac13,\\
 \E U^3&=\frac12\int_{-1}^{1}u^3\dd u=0.
\end{align*}
所以
\[
 \Cov(X,Y)=\E U^3-(\E U)(\E U^2)=0.
\]
但取事件 $A=\{|X|\le1/2\}$ 与 $B=\{Y>1/4\}$，有
\[
 \Prob(A)=\frac12,\qquad \Prob(B)=\frac12,\qquad \Prob(A\cap B)=0.
\]
若 $X,Y$ 独立，最后一个数应为 $1/4$，矛盾。这里 $Y=X^2$ 所携带的确定关系没有被一个二阶
混合矩检测出来。
\end{example}

反过来，若 $X,Y\in L^2$ 且独立，定理~\ref{theorem:5-1-2-3} 应用于 $X,Y$，并结合
Cauchy--Schwarz 保证的可积性，给出 $\E(XY)=\E X\E Y$，所以 $\Cov(X,Y)=0$。
只有在联合高斯等额外结构下，零协方差才反过来推出独立；联合高斯的定义和证明均留到
后文，在此不作为独立性判据。

\begin{remark}[独立和的振荡指数]\label{rem:5-1-2-cf-preview}
指数测试函数给出独立和的一个基本乘法公式。若 $X,Y$ 独立且实值，则
\[
 \E e^{it(X+Y)}=(\E e^{itX})(\E e^{itY}),\qquad t\in\R,
\]
因为两个因子有模一，定理~\ref{theorem:5-1-2-3} 可直接应用。若 $X,Y$ 是独立 Rademacher
变量，则 $\E e^{itX}=\cos t$，故左端为 $\cos^2t$。另一方面，$X+Y$ 在
$-2,0,2$ 上的概率分别是 $1/4,1/2,1/4$，直接求和得到
\[
 \E e^{it(X+Y)}=\frac14e^{-2it}+\frac12+\frac14e^{2it}
 =\frac{1+\cos(2t)}2=\cos^2t.
\]
这个计算具体展示了期望因子化怎样把和转化为乘积。第~\ref{subsec:5-4-3}~节将把
$t\mapsto\E e^{itX}$ 系统地定义为特征函数，并证明它对分布的识别与极限性质。
\end{remark}

\begin{exercise}
\begin{enumerate}
  \item 设 $X_1,\ldots,X_n$ 独立，$A_i\in\sigma(X_i)$。从定义证明
  $\1_{A_1},\ldots,\1_{A_n}$ 独立，再用命题~\ref{proposition:5-1-2-4} 给出第二种证明。
  \item 为例~\ref{ex:5-1-2-2} 写出 $(X,Y,Z)$ 的完整取值表，并逐项验证三对联合分布。
  \item 设 $X,Y$ 独立且分别有密度 $f,g$。只用 Tonelli 定理和换元证明
  $X+Y$ 的分布具有密度 $f*g$，并说明第四章的卷积定义怎样进入该密度公式。
  \item 从矩形生成类出发重证定理~\ref{theorem:5-1-2-2}，说明有限测度唯一性为何足够。
\end{enumerate}
\end{exercise}

有限乘积结构回答了固定多个事件如何同时发生。无穷序列提出了两个新问题：稀有事件会不会发生
无限多次，以及删掉任意有限初段后还留下多少信息。


\subsection{Borel--Cantelli、尾 \texorpdfstring{$\sigma$}{sigma}-代数与零一律}\label{subsec:5-1-3}
\index{Borel--Cantelli 引理}\index{尾 $\sigma$-代数}\index{Kolmogorov 零一律}

假设第 $n$ 次试验有一个坏事件 $A_n$。逐次让 $\Prob(A_n)$ 变小，并不能单独排除坏事一再发生；
真正可计算的量是整条尾概率级数。收敛时，可数次可加性已经足够。发散时，若事件可以始终一起发生，
概率之和便会严重重复计数；这正是第二个方向需要独立性的原因。

\begin{definition}[limsup 事件]\label{def:5-1-3-1}
对事件列 $(A_n)$，定义
\[
 \{A_n\ \mathrm{i.o.}\}
 :=\limsup_{n\to\infty}A_n
 :=\bigcap_{N=1}^\infty\bigcup_{n\ge N}A_n.
\]
缩写 ``i.o.'' 表示 infinitely often。一个样本点属于该事件，当且仅当每个有限时刻之后仍有某个
$A_n$ 发生，也就是属于无穷多个 $A_n$。
\end{definition}

\begin{theorem}[Borel--Cantelli I]\label{theorem:5-1-3-1}
若 $\sum_{n=1}^\infty\Prob(A_n)<\infty$，则
\[
 \Prob(A_n\ \mathrm{i.o.})=0.
\]
此结论不要求事件独立。
\end{theorem}

\begin{proof}
令 $B_N=\bigcup_{n\ge N}A_n$。由次可加性，
\[
 \Prob(B_N)\le\sum_{n\ge N}\Prob(A_n).
\]
右端是收敛非负级数的尾和，故随 $N\to\infty$ 趋于零。事件列 $(B_N)$ 递减且
$\Prob(B_1)\le1<\infty$；测度从上连续给出
\[
 \Prob(A_n\ \mathrm{i.o.})
 =\Prob\!\left(\bigcap_NB_N\right)
 =\lim_{N\to\infty}\Prob(B_N)=0.
\]
\end{proof}

\begin{theorem}[Borel--Cantelli II]\label{theorem:5-1-3-2}
若事件 $(A_n)$ 相互独立且
$\sum_{n=1}^\infty\Prob(A_n)=\infty$，则
\[
 \Prob(A_n\ \mathrm{i.o.})=1.
\]
\end{theorem}

\begin{proof}
记 $p_n=\Prob(A_n)$。相互独立性对补事件仍成立。固定 $N$，对 $M\ge N$ 有
\begin{align*}
 \Prob\!\left(\bigcap_{n=N}^M A_n^c\right)
 &=\prod_{n=N}^M(1-p_n)\\
 &\le\exp\!\left(-\sum_{n=N}^Mp_n\right),
\end{align*}
其中使用了 $1-u\le e^{-u}$（$0\le u\le1$）。对固定 $N$，发散级数的尾部仍发散，故右端
随 $M\to\infty$ 趋于零。左侧事件随 $M$ 递减，且概率至多为一；测度从上连续给出
\[
 \Prob\!\left(\bigcap_{n=N}^\infty A_n^c\right)=0.
\]
取补集可得
$\Prob(\bigcup_{n\ge N}A_n)=1$。最后，概率一事件的可数交仍有概率一，因为其补集包含于一列
零概率集的可数并。因此
\[
 \Prob\!\left(\bigcap_{N=1}^\infty\bigcup_{n\ge N}A_n\right)=1.
\]
交内的事件正是 $\{A_n\ \mathrm{i.o.}\}$。
\end{proof}

\begin{proposition}[可和概率推出 a.s. 控制]\label{proposition:5-1-3-4}
若 $a_n\ge0$ 且
\[
 \sum_{n=1}^\infty\Prob(|X_n|>a_n)<\infty,
\]
则几乎必然存在依赖于样本点的 $N<\infty$，使得对所有 $n\ge N$ 都有
$|X_n|\le a_n$。
\end{proposition}

\begin{proof}
对 $A_n=\{|X_n|>a_n\}$ 应用定理~\ref{theorem:5-1-3-1}，得到只有有限多个 $A_n$
发生的事件具有概率一。这句话逐点展开正是所述最终控制。
\end{proof}

\begin{example}[抛硬币无限正面]\label{ex:5-1-3-1}
设 $(X_n)$ 是独立 Rademacher 序列，$H_n=\{X_n=1\}$、$T_n=\{X_n=-1\}$。
两族事件分别相互独立，并且
\[
 \sum_n\Prob(H_n)=\sum_n\Prob(T_n)=\sum_n\frac12=\infty.
\]
分别应用 Borel--Cantelli II，得到正面无限次出现的事件和反面无限次出现的事件都具有概率一；二者
的交仍有概率一。因此几乎每条路径都会永远继续出现两种结果，而不是在某次之后固定下来。
\end{example}

\begin{example}[稀有事件的可和阈值]\label{ex:5-1-3-threshold}
设 $U_n\sim\mathrm{Unif}(0,1)$ 相互独立，$\alpha>0$，并令
\[
 A_n=\{U_n\le n^{-\alpha}\}.
\]
由可测变换保持独立，$(A_n)$ 相互独立，且
$\Prob(A_n)=n^{-\alpha}$。于是
\[
 \sum_n\Prob(A_n)=\sum_n n^{-\alpha}
 \begin{cases}
 <\infty,&\alpha>1,\\
 =\infty,&0<\alpha\le1.
 \end{cases}
\]
第一种情形由 Borel--Cantelli I 得 $A_n$ 只发生有限次；第二种情形由 Borel--Cantelli II 得
$A_n$ 几乎必然发生无限次。$p$-级数在 $\alpha=1$ 的收敛阈值，在这里变成路径行为的相变。
\end{example}

\begin{example}[从可和尾界读出收敛速度]\label{ex:5-1-3-2}
假设对每个正有理数 $\varepsilon$，都有
\[
 \sum_{n=1}^\infty\Prob(|X_n-X|>\varepsilon)<\infty.
\]
对每个 $k\ge1$，Borel--Cantelli I 给出概率一事件 $E_k$，使得在 $E_k$ 上
$|X_n-X|\le1/k$ 最终成立。事件 $E=\bigcap_kE_k$ 仍有概率一。若 $\omega\in E$ 且
$\varepsilon>0$，取 $k$ 使 $1/k<\varepsilon$，再取 $E_k$ 对应的最终指标，便有
$|X_n(\omega)-X(\omega)|<\varepsilon$ 最终成立。因此 $X_n\to X$ a.s.
\end{example}

无穷独立系统还有另一种刚性。一个性质若删去任意有限初段后都不改变，那么它属于系统的尾部信息。

\begin{definition}[尾 $\sigma$-代数]\label{def:5-1-3-2}
对随机元序列 $(X_n)$，定义其尾 $\sigma$-代数
\[
 \mathcal T=\bigcap_{N=1}^\infty\sigma(X_N,X_{N+1},\ldots).
\]
$\mathcal T$ 中的事件称为尾事件。
\end{definition}

在规范乘积样本空间上，尾事件可以直观地理解为改变有限多个坐标后不变的事件。然而在一般样本空间
上，“修改坐标”未必定义了另一个样本点；正式定义始终是上述 $\sigma$-代数之交。以下证明也只使用
这个定义。

\begin{lemma}[独立初段与尾段]\label{lemma:5-1-3-3}
若 $(X_n)_{n\ge1}$ 相互独立，则对每个 $N\ge1$，
\[
 \sigma(X_1,\ldots,X_N)
 \quad\text{与}\quad
 \sigma(X_{N+1},X_{N+2},\ldots)
\]
独立。
\end{lemma}

\begin{proof}
令 $\mathcal P_-$ 为所有形如
\[
 \bigcap_{j=1}^N\{X_j\in B_j\}
\]
的事件，其中 $B_j$ 在 $X_j$ 的状态空间中可测；允许 $B_j$ 为全空间。
$\mathcal P_-$ 是生成 $\sigma(X_1,\ldots,X_N)$ 的 $\pi$-系统。
令 $\mathcal P_+$ 为所有有限尾柱事件
\[
 \bigcap_{j=N+1}^{M}\{X_j\in B_j\},\qquad M\ge N+1,
\]
并包含 $\Omega$。它是生成 $\sigma(X_{N+1},X_{N+2},\ldots)$ 的 $\pi$-系统。

若 $A\in\mathcal P_-$、$B\in\mathcal P_+$，相互独立性把构成 $A\cap B$ 的有限个坐标事件
全部因子化；分别把初段、尾段内部的因子重新相乘，得到
\[
 \Prob(A\cap B)=\Prob(A)\Prob(B).
\]
对这两个生成 $\pi$-系统应用定理~\ref{theorem:5-1-2-1} 的二元情形，便得到两个完整
$\sigma$-代数独立。
\end{proof}

\begin{theorem}[Kolmogorov 零一律]\label{theorem:5-1-3-3}
若 $(X_n)_{n\ge1}$ 独立，则每个尾事件 $A\in\mathcal T$ 满足
\[
 \Prob(A)\in\{0,1\}.
\]
\end{theorem}

\begin{proof}
固定 $A\in\mathcal T$。对每个 $N$，有
$A\in\sigma(X_{N+1},X_{N+2},\ldots)$，所以由
引理~\ref{lemma:5-1-3-3}，$A$ 与 $\sigma(X_1,\ldots,X_N)$ 独立。

令 $\mathcal P$ 为所有有限坐标柱事件；它是一个 $\pi$-系统，并生成
\[
 \mathcal H=\sigma(X_1,X_2,\ldots).
\]
上一步说明 $A$ 与每个 $B\in\mathcal P$ 独立。定义
\[
 \mathcal D=
 \{B\in\mathcal H:\Prob(A\cap B)=\Prob(A)\Prob(B)\}.
\]
首先 $\Omega\in\mathcal D$。若 $B\in\mathcal D$，则
\begin{align*}
 \Prob(A\cap B^c)
 &=\Prob(A)-\Prob(A\cap B)\\
 &=\Prob(A)-\Prob(A)\Prob(B)
 =\Prob(A)\Prob(B^c),
\end{align*}
故 $B^c\in\mathcal D$。若 $(B_n)$ 两两不交且 $B_n\in\mathcal D$，则
\begin{align*}
 \Prob\!\left(A\cap\bigcup_{n=1}^{\infty}B_n\right)
 &=\sum_{n=1}^{\infty}\Prob(A\cap B_n)\\
 &=\Prob(A)\sum_{n=1}^{\infty}\Prob(B_n)
 =\Prob(A)\Prob\!\left(\bigcup_{n=1}^{\infty}B_n\right).
\end{align*}
故 $\mathcal D$ 是 $\lambda$-系统。它包含 $\mathcal P$，所以 $\pi$--$\lambda$ 定理给出
$\mathcal D=\mathcal H$。

尾事件 $A$ 本身属于 $\mathcal H$，故可以在定义 $\mathcal D$ 的等式中取 $B=A$。若
$p=\Prob(A)$，便有
\[
 p=\Prob(A\cap A)=p^2.
\]
由于 $0\le p\le1$，等式 $p(1-p)=0$ 蕴含 $p=0$ 或 $p=1$。
\end{proof}

\begin{figure}[htbp]
\centering
\begin{tikzpicture}[node distance=8mm and 9mm]
  \node[book node] (a1) {$A_1$};
  \node[book node,right=of a1] (a2) {$A_2$};
  \node[book node,right=of a2] (dots) {$\cdots$};
  \node[book strong node,right=of dots] (an) {$A_N,A_{N+1},\ldots$};
  \draw[decorate,decoration={brace,mirror,amplitude=4pt},BookAccent]
    ([yshift=-3pt]an.south west)--([yshift=-3pt]an.south east)
    node[midway,below=6pt,book label]{$B_N=\bigcup_{n\ge N}A_n$};
  \node[book warm node,below=19mm of a2] (limsup)
    {$\displaystyle\limsup_nA_n=\bigcap_N B_N$};
  \draw[book accent arrow] (an.south)--(limsup.east);

  \node[book node,below=of limsup] (finite) {有限坐标柱集};
  \node[book strong node,right=of finite] (all) {$\sigma(X_1,X_2,\ldots)$};
  \node[book warning node,right=of all] (self) {$A$ 与自身独立};
  \node[book warm node,below=of self] (p) {$p=p^2$};
  \draw[book arrow] (finite)--node[above,book label]{$\pi$--$\lambda$}(all);
  \draw[book arrow] (all)--(self);
  \draw[book accent arrow] (self)--(p);
\end{tikzpicture}
\caption{上极限事件不受删除有限个初项影响；独立序列的尾事件经生成类推广后与自身独立。}
\label{fig:5-1-3-tail-zero-one}
\end{figure}

\begin{example}[三个尾事件]\label{ex:5-1-3-3}
设 $(X_n)$ 为实值随机变量序列。
\begin{enumerate}[label=\textup{(\roman*)}]
  \item 事件 $\{\limsup_nX_n=+\infty\}$ 可写成
  \[
   \bigcap_{m=1}^\infty\bigcap_{N=1}^\infty
   \bigcup_{n\ge N}\{X_n>m\}.
  \]
  对每个固定 $N_0$，删去式中 $n<N_0$ 的有限项不改变事件，且余式只用
  $X_{N_0},X_{N_0+1},\ldots$，故它属于 $\mathcal T$。
  \item 级数 $\sum_nX_n$ 收敛的事件由 Cauchy 判据写成
  \[
   \bigcap_{k=1}^\infty\bigcup_{N=1}^\infty
   \bigcap_{q\ge p\ge N}
   \left\{\left|\sum_{n=p}^qX_n\right|<\frac1k\right\}.
  \]
  对任意 $N_0$，可把外层所选 $N$ 限制为 $N\ge N_0$ 而不改变该事件，所以它属于每个
  $\sigma(X_{N_0},X_{N_0+1},\ldots)$。
  \item 经验平均 $n^{-1}\sum_{j=1}^nX_j$ 有有限极限的事件也是尾事件。固定 $N_0$ 后，原平均与
  $n^{-1}\sum_{j=N_0}^nX_j$ 的差是固定有限和
  $n^{-1}\sum_{j<N_0}X_j$，逐点趋于零。因此两列同时有或同时没有有限极限。对后一列记
  $Y_n^{(N_0)}=n^{-1}\sum_{j=N_0}^nX_j$，其有有限极限的事件由 Cauchy 判据表示为
  \[
   \bigcap_{k=1}^{\infty}\bigcup_{M=N_0}^{\infty}
   \bigcap_{p,q\ge M}
   \left\{|Y_p^{(N_0)}-Y_q^{(N_0)}|<\frac1k\right\}.
  \]
  每个括号内的事件属于
  $\sigma(X_{N_0},X_{N_0+1},\ldots)$，因而整个事件属于该 $\sigma$-代数。
\end{enumerate}
若序列独立，零一律说明这些事件的概率不能严格落在 $0$ 与 $1$ 之间。只有在另行证明其中一个事件
具有正概率后，才可以把它的概率提升为一；零一律本身不负责证明正概率。
\end{example}

\begin{example}[为何第二部分需要相关性假设]\label{ex:5-1-3-4}
具体取 $\Omega=\{0,1\}$ 上的均匀概率、$A=\{1\}$，并令 $A_n=A$ 对所有 $n$ 成立。则
\[
 \sum_n\Prob(A_n)=\infty,
 \qquad
 \{A_n\ \mathrm{i.o.}\}=A,
\]
所以无限发生事件的概率仍是 $1/2$。这里每个事件都与下一个完全重合；发散的概率和把同一份质量
重复计算了无穷次。这个反例只撤去独立性，精确标出了 Borel--Cantelli II 中该假设的位置。
\end{example}

概率一不表示逐点无例外。例如在无限抛硬币的规范乘积空间中，“从某次以后永远正面”对应的路径
确实存在。若 $H_n$ 是第 $n$ 次正面事件，则对每个 $N$，测度从上连续给出
\[
 \Prob\!\left(\bigcap_{n\ge N}H_n\right)
 =\lim_{M\to\infty}2^{-(M-N+1)}=0.
\]
因此这些事件对 $N$ 的可数并仍为零概率。概率论允许这样的非空异常集；它只保证异常集能够被一个
可测零集容纳。

\begin{exercise}
\begin{enumerate}
  \item 在独立公平硬币模型中证明任意给定的有限正反面图样几乎必然出现无限多次。提示：取互不重叠
  的等长分块，以恢复事件独立性。
  \item 若 $\Prob(|X_n-X|>2^{-n})\le2^{-n}$，用
  命题~\ref{proposition:5-1-3-4} 证明 $X_n\to X$ a.s.，并给出逐点最终误差界。
  \item 固定 $c\in\R$，证明“$\sup_nX_n\le c$”未必是尾事件，而
  “$\limsup_nX_n\le c$”是尾事件，并给出前一断言的两项确定性序列反例。
  \item 补全定理~\ref{theorem:5-1-3-3} 中 $\mathcal D$ 对两两不交可数并封闭的计算。
\end{enumerate}
\end{exercise}

独立性把信息块隔开；条件期望则研究一个信息块已经揭示之后，对其余未知量应当如何作可测预测。
下一节将从 Radon--Nikodym 表示出发，把这种预测构造成一个 $L^1$ 算子。

\section{条件信息、核与贝叶斯公式}\label{sec:05-02}

\subsection{条件期望：Radon--Nikodym 的信息版本}\label{subsec:5-2-1}
\index{条件期望}\index{版本}\index{Doob--Dynkin 因子化}

设 $(\Omega,\mathcal F,\Prob)$ 是概率空间。给定 $X\in L^1(\Prob)$，若观察者能看见
$\mathcal F$ 中的全部事件，那么 $X$ 本身就是可用资料；若观察者只能分辨一个子
$\sigma$-代数 $\mathcal G\subset\mathcal F$ 中的事件，我们就必须把 $X$ 压缩成一个
$\mathcal G$-可测变量。压缩不能随意进行：在每个可观察事件上，它应当保留 $X$ 的总平均。

先看一个没有零概率事件搅局的模型。

\begin{example}[只报奇偶时如何预测赔付]\label{ex:5-2-1-prediction-pressure}
公平骰子的点数记为 $D$，赔付为 $X=D^2$。理赔员只收到“奇数”或“偶数”这一位资料，故可用的
预测必为
\[
 Y=a\1_{\{D\text{ 为奇数}\}}+b\1_{\{D\text{ 为偶数}\}}.
\]
其均方误差为
\[
 \frac16\sum_{d\in\{1,3,5\}}(d^2-a)^2
 +\frac16\sum_{d\in\{2,4,6\}}(d^2-b)^2.
\]
两组平方和可逐项展开为
\begin{align*}
 \sum_{d\in\{1,3,5\}}(d^2-a)^2
 &=3a^2-70a+707
 =3\left(a-\frac{35}{3}\right)^2+\frac{896}{3},\\
 \sum_{d\in\{2,4,6\}}(d^2-b)^2
 &=3b^2-112b+1568
 =3\left(b-\frac{56}{3}\right)^2+\frac{1568}{3}.
\end{align*}
两个平方项的系数都严格为正，因此唯一极小点为
\[
 a=\frac{1+9+25}{3}=\frac{35}{3},
 \qquad
 b=\frac{4+16+36}{3}=\frac{56}{3}.
\]
这个答案还有一个不涉及微分的刻画。若 $A_o=\{D\text{ 为奇数}\}$、
$A_e=\{D\text{ 为偶数}\}$，则
\[
 \E[Y\1_{A_o}]=\frac{35}{6}=\E[X\1_{A_o}],
 \qquad
 \E[Y\1_{A_e}]=\frac{28}{3}=\E[X\1_{A_e}].
\]
具体地，在奇数块上
\[
 \E[Y\1_{A_o}]=\frac12\cdot\frac{35}{3}=\frac{35}{6},
 \qquad
 \E[X\1_{A_o}]=\frac{1^2+3^2+5^2}{6}=\frac{35}{6},
\]
偶数块上两边同为
$(4+16+36)/6=28/3$。也就是说，预测在每个已知信息块上保留了原变量的积分。后一种刻画不需要把信息分割列成有限张
清单，因而可以推广到连续观测。
\end{example}

初等概率中常见的比值
$\E[X\mid Z=y]$ 在 $\Prob(Z=y)>0$ 时可以计算，但连续分布通常满足
$\Prob(Z=y)=0$。在零分母前继续写比值不会产生定义。正确的次序是先给定可观察的
$\sigma$-代数并保留其上全部积分，再在状态空间足够良好时把所得对象组织成关于 $y$ 的核。

\begin{definition}[条件期望]\label{def:5-2-1-1}
设 $X\in L^1(\Prob)$ 且 $\mathcal G\subset\mathcal F$ 是子 $\sigma$-代数。若随机变量
$Y$ 满足
\begin{enumerate}
  \item $Y$ 是 $\mathcal G$-可测且 $Y\in L^1(\Prob)$；
  \item 对每个 $G\in\mathcal G$，
  \[
    \int_GY\,\dd\Prob=\int_GX\,\dd\Prob,
  \]
\end{enumerate}
则称 $Y$ 是 $X$ 关于 $\mathcal G$ 的一个条件期望，记作
$Y=\E[X\mid\mathcal G]$。
\end{definition}

\begin{definition}[关于随机变量的条件期望]\label{def:5-2-1-2}
若 $Z$ 是随机元，则
\[
 \E[X\mid Z]:=\E[X\mid\sigma(Z)].
\]
这是关于 $\sigma(Z)$ 的等价类记号，而不是先对每个 $z$ 定义一个比值。若 $Z$ 取值于标准
Borel 空间，\cref{proposition:5-2-1-doob-dynkin} 将给出一个 Borel 函数 $g$，使可选择
$\E[X\mid Z]=g(Z)$。
\end{definition}

\begin{definition}[版本]\label{def:5-2-1-3}
两个几乎处处相等的 $\mathcal G$-可测函数代表同一个条件期望等价类；其中任一具体代表称为一个
\emph{版本}。版本只在 $\Prob$-a.s. 意义下唯一。若等式带有不可数多个时间或参数，逐参数成立的
零测例外集未必能合并成一个零测集，因而还需另行选择相容版本。
\end{definition}

第~3~章已经从 Radon--Nikodym 定理得到过相同的积分表示。现在把它正式识别为条件期望。

\begin{theorem}[条件期望的存在唯一性]\label{theorem:5-2-1-1}
对每个 $X\in L^1(\Prob)$ 与每个子 $\sigma$-代数
$\mathcal G\subset\mathcal F$，条件期望 $\E[X\mid\mathcal G]$ 存在，并在
$\Prob$-a.s. 意义下唯一。
\end{theorem}

\begin{proof}
先设 $X\geq0$。在 $(\Omega,\mathcal G)$ 上定义
\[
 \nu(G)=\int_GX\,\dd\Prob.
\]
因为 $X\in L^1$，$\nu$ 是有限正测度；若 $\Prob(G)=0$，则 $\nu(G)=0$，故
$\nu\ll\Prob|_{\mathcal G}$。由 Radon--Nikodym 定理
\ref{theorem:3-3-1-1}，存在 $\mathcal G$-可测 $Y\geq0$ 使
\[
 \nu(G)=\int_GY\,\dd\Prob\qquad(G\in\mathcal G).
\]
取 $G=\Omega$ 得 $\int Y=\int X<\infty$，所以 $Y\in L^1$。

若 $X$ 为实值，分别对 $X^+$ 与 $X^-$ 作上述构造，记所得密度为 $Y^+,Y^-$。二者均可积，
故 $Y=Y^+-Y^-$ 可积，并在每个 $G\in\mathcal G$ 上满足所需积分等式。若 $X$ 为复值，分别
处理 $\Re X$ 与 $\Im X$，再组成 $Y=Y_1+iY_2$。

最后证明唯一性。设实值 $Y_1,Y_2$ 都满足定义，并令
$A=\{Y_1>Y_2\}\in\mathcal G$。两条积分保持式给出
\[
 \int_A(Y_1-Y_2)\,\dd\Prob=0.
\]
被积函数在 $A$ 上严格为正；非负函数积分为零当且仅当它几乎处处为零，故
$\Prob(A)=0$。交换 $Y_1,Y_2$ 得 $\Prob(Y_2>Y_1)=0$，所以二者 a.s. 相等。复值情形对
实部和虚部分别使用这个结论。
\end{proof}

“关于 $Z$ 可测”还没有自动变成“是 $Z$ 的函数”。下面的证明把这一点落实到一个可数水平集族，
从而避免为每个水平集分别选像后直接跳到因子化。

\begin{proposition}[Doob--Dynkin 因子化]\label{proposition:5-2-1-doob-dynkin}
设 $Z:\Omega\to T$ 取值于标准 Borel 空间，实值随机变量 $Y$ 对 $\sigma(Z)$ 可测。则存在
Borel 函数 $g:T\to\R$ 使 $Y=g(Z)$。特别地，对 $X\in L^1$ 可选择一个 Borel 函数 $g$，
使 $\E[X\mid Z]=g(Z)$ a.s.
\end{proposition}

\begin{proof}
对每个 $q\in\Q$，集合 $\{Y<q\}$ 属于 $\sigma(Z)$，故可取 Borel 集
$A_q\subset T$ 使
\[
 Z^{-1}(A_q)=\{Y<q\}.
\]
这些 $A_q$ 未必彼此相容。令
\[
 B_q=\bigcup_{r<q,\ r\in\Q}A_r.
\]
于是 $B_q$ 随 $q$ 非降，并且
\[
 B_q=\bigcup_{s<q,\ s\in\Q}B_s,
 \qquad
 Z^{-1}(B_q)=\bigcup_{r<q}\{Y<r\}=\{Y<q\}.
\]
两个集合
\[
 N_+=T\setminus\bigcup_{q\in\Q}B_q,
 \qquad
 N_- =\bigcap_{q\in\Q}B_q
\]
都是 Borel 集，而且都不与 $Z(\Omega)$ 相交：给定 $\omega$，总能取有理数
$q>Y(\omega)$ 与 $r<Y(\omega)$。先定义扩展实值函数
\[
 h(t)=\inf\{q\in\Q:t\in B_q\},
\]
其中空集的下确界取 $+\infty$。对每个 $a\in\R$，
\[
 \{h<a\}=\bigcup_{q<a,\ q\in\Q}B_q,
\]
故 $h$ 可测。令 $g=h$ 于 $T\setminus(N_+\cup N_-)$，并在
$N_+\cup N_-$ 上令 $g=0$，则 $g$ 是处处实值的 Borel 函数。对每个 $\omega$，
\[
 \{q\in\Q:Z(\omega)\in B_q\}=\{q\in\Q:Y(\omega)<q\},
\]
这个集合的下确界等于 $Y(\omega)$，所以 $g(Z)=Y$。

将第一部分用于条件期望的任一实值版本即可。复值条件期望对实部、虚部分别因子化。
\end{proof}

条件期望还需要容纳非负但不可积的变量；稍后核表示中取任意非负测试函数时会用到这一扩展。

\begin{proposition}[非负变量的扩展条件期望]\label{proposition:5-2-1-nonnegative}
设 $X:\Omega\to[0,\infty]$ 可测，不要求 $X$ 可积。可定义一个
$[0,\infty]$ 值 $\mathcal G$-可测变量 $\E[X\mid\mathcal G]$，使
\begin{equation}\label{eq:5-2-extended-ce}
 \int_G\E[X\mid\mathcal G]\,\dd\Prob
 =\int_GX\,\dd\Prob\qquad(G\in\mathcal G),
\end{equation}
两边允许为 $+\infty$。它在 a.s. 意义下唯一。若 $0\leq X_n\uparrow X$，则可相容地选择版本使
\begin{equation}\label{eq:5-2-conditional-mct}
 \E[X_n\mid\mathcal G]\uparrow\E[X\mid\mathcal G]
 \quad\text{a.s.}
\end{equation}
\end{proposition}

\begin{proof}
先记录可积变量的单调性。若 $U,V\in L^1$ 且 $U\leq V$ a.s.，取条件期望版本
$U_0,V_0$，并令 $A=\{U_0>V_0\}\in\mathcal G$。则
\[
 \int_A(U_0-V_0)\,\dd\Prob
 =\int_A(U-V)\,\dd\Prob\leq0.
\]
左端是非负函数的积分，故 $\Prob(A)=0$。所以 $U_0\leq V_0$ a.s.

令 $X_n=X\wedge n$。逐个选取 $Y_n=\E[X_n\mid\mathcal G]$ 的非负版本。由刚证的单调性，
$Y_n\leq Y_{n+1}$ a.s.；删除这些例外集的可数并，并在该并集上把所有 $Y_n$ 改为零，便可假设
$Y_n(\omega)$ 对每个 $\omega$ 都非降。定义 $Y=\lim_nY_n$。它非负且 $\mathcal G$-可测。
对每个 $G\in\mathcal G$，两次使用单调收敛定理得到
\[
 \int_GY\,\dd\Prob
 =\lim_n\int_GY_n\,\dd\Prob
 =\lim_n\int_G(X\wedge n)\,\dd\Prob
 =\int_GX\,\dd\Prob.
\]
这证明 \eqref{eq:5-2-extended-ce}。

若 $Y,Z$ 都满足 \eqref{eq:5-2-extended-ce}，对 $m,n\geq1$ 令
\[
 A_{m,n}=\{Y\geq Z+1/m\}\cap\{Z\leq n\}\in\mathcal G.
\]
在 $A_{m,n}$ 上 $Z$ 的积分至多为 $n$，故两条保持式说明 $Y,Z$ 在该集合上的积分都是同一个
有限数。于是
\[
 0=\int_{A_{m,n}}(Y-Z)\,\dd\Prob
 \geq \frac1m\Prob(A_{m,n}),
\]
所以 $A_{m,n}$ 为零测集。对 $m,n$ 取可数并得 $\Prob(Y>Z)=0$；交换二者便得唯一性。

最后，若 $0\leq X_n\uparrow X$，则对每个 $n,m$ 都有
$X_n\wedge m\leq X_{n+1}\wedge m$。对这些可积截断使用已证的单调性，再合并关于
$(n,m)$ 的可数个例外集并令 $m\to\infty$，得到
\[
 \E[X_n\mid\mathcal G]\leq\E[X_{n+1}\mid\mathcal G]\quad\text{a.s.}
\]
在一个共同零集上修改版本后，令 $W=\lim_n\E[X_n\mid\mathcal G]$。对每个
$G\in\mathcal G$，单调收敛给
$\int_GW=\lim_n\int_GX_n=\int_GX$。扩展条件期望的唯一性说明
$W=\E[X\mid\mathcal G]$ a.s.，即 \eqref{eq:5-2-conditional-mct}。
\end{proof}

\begin{example}[不可积变量的条件平均可以为无穷]\label{ex:5-2-1-nonintegrable}
在 $((0,1),\Borel((0,1)),m)$ 上令 $X(\omega)=\omega^{-1}$，并令
$\mathcal G=\{\varnothing,(0,1)\}$。每个 $\mathcal G$-可测非负函数都是常数，且
\[
 \E[X\wedge n]
 =\int_0^{1/n}n\,\dd\omega+\int_{1/n}^1\frac{\dd\omega}{\omega}
 =1+\log n\longrightarrow+\infty.
\]
所以 $\E[X\mid\mathcal G]$ 是常数 $+\infty$。扩展条件期望并不承诺有限值；它承诺的是
非负积分恒等式。
\end{example}

下面几条运算都从积分刻画推出，而不是从点态条件概率推出。

\begin{proposition}[基本运算与 $L^1$ 收缩]\label{proposition:5-2-1-2}
对 $X,X_1,X_2\in L^1$、标量 $a,b$，条件期望满足
\[
 \E[aX_1+bX_2\mid\mathcal G]
 =a\E[X_1\mid\mathcal G]+b\E[X_2\mid\mathcal G].
\]
若 $X\geq0$，则 $\E[X\mid\mathcal G]\geq0$ a.s.；并且对实值或复值 $X$，
\begin{equation}\label{eq:5-2-l1-contraction}
 \bigl|\E[X\mid\mathcal G]\bigr|
 \leq\E[|X|\mid\mathcal G]\quad\text{a.s.},
 \qquad
 \|\E[X\mid\mathcal G]\|_1\leq\|X\|_1.
\end{equation}
\end{proposition}

\begin{proof}
线性组合是 $\mathcal G$-可测可积函数，并在每个 $G\in\mathcal G$ 上具有正确积分；唯一性给出
线性。正性已包含在存在性证明对非负 Radon--Nikodym 密度的选择中，也可由上一证明的单调性取
$U=0,V=X$ 得到。

若 $X$ 为实值，则 $-|X|\leq X\leq|X|$。单调性和线性给
\[
 -\E[|X|\mid\mathcal G]
 \leq\E[X\mid\mathcal G]
 \leq\E[|X|\mid\mathcal G],
\]
即第一条模长不等式。

若 $X$ 为复值，记 $Y=\E[X\mid\mathcal G]$、
$Z=\E[|X|\mid\mathcal G]$，并取 $[0,2\pi)$ 中一个可数稠密集 $\Theta$。对每个
$\theta\in\Theta$，实值情形给出
\[
 \Re(e^{-i\theta}Y)
 =\E[\Re(e^{-i\theta}X)\mid\mathcal G]
 \leq Z\quad\text{a.s.}
\]
删除这些例外集的可数并后，对同一个 $\omega$ 及全部 $\theta\in\Theta$ 上式成立。由于
$\sup_{\theta\in\Theta}\Re(e^{-i\theta}z)=|z|$，得到 $|Y|\leq Z$。

最后取期望并用定义中的 $G=\Omega$，得到
\[
 \|Y\|_1\leq\E Z=\E|X|.
\]
\end{proof}

\begin{proposition}[取出已知因子]\label{proposition:5-2-1-3}
设 $Z$ 为 $\mathcal G$-可测实值或复值随机变量。若 $Z$ 有界且 $X\in L^1$，则
\begin{equation}\label{eq:5-2-pull-out}
 \E[ZX\mid\mathcal G]=Z\E[X\mid\mathcal G].
\end{equation}
更一般地，只要 $X\in L^1$ 且 $ZX\in L^1$，同一等式仍成立，并且右端自动可积。
\end{proposition}

\begin{proof}
先设 $Z=\sum_{j=1}^mc_j\1_{A_j}$ 是有界 $\mathcal G$-简单函数，其中 $A_j$ 两两不交。
对任意 $G\in\mathcal G$，
\[
 \int_G Z\E[X\mid\mathcal G]\,\dd\Prob
 =\sum_{j=1}^mc_j\int_{G\cap A_j}\E[X\mid\mathcal G]\,\dd\Prob
 =\sum_{j=1}^mc_j\int_{G\cap A_j}X\,\dd\Prob
 =\int_GZX\,\dd\Prob.
\]
故唯一性给出 \eqref{eq:5-2-pull-out}。对有界实值 $Z$，取一致有界且逐点收敛于 $Z$ 的
$\mathcal G$-简单函数 $Z_n$。支配收敛定理分别用于 $Z_nX$ 与
$Z_n\E[X\mid\mathcal G]$，再结合 $L^1$ 收缩
\[
 \|\E[Z_nX-ZX\mid\mathcal G]\|_1\leq\|Z_nX-ZX\|_1,
\]
即可在 $L^1$ 中令 $n\to\infty$。复值 $Z$ 对实部、虚部分别处理。

现在只假设 $ZX\in L^1$。先设 $Z$ 实值，并令
\[
 Z_m=(-m)\vee(Z\wedge m).
\]
则 $|(Z_m-Z)X|\leq|ZX|\1_{\{|Z|>m\}}$，所以 $Z_mX\to ZX$ 于 $L^1$。由
\eqref{eq:5-2-l1-contraction}，
$\E[Z_mX\mid\mathcal G]\to\E[ZX\mid\mathcal G]$ 于 $L^1$。取一个几乎处处收敛的子列；
在该子列上，已证的有界因子公式给左边等于
$Z_m\E[X\mid\mathcal G]$，而后者逐点收敛到
$Z\E[X\mid\mathcal G]$。因此 \eqref{eq:5-2-pull-out} a.s. 成立，且右端因等于一个
$L^1$ 变量而可积。复值情形分解 $Z$ 的实部、虚部；由 $|\Re Z|,|\Im Z|\leq|Z|$，所需积
仍可积。
\end{proof}

\begin{theorem}[塔式性质]\label{theorem:5-2-1-4}
若 $\mathcal H\subset\mathcal G\subset\mathcal F$ 且 $X\in L^1$，则
\begin{equation}\label{eq:5-2-tower}
 \E\bigl[\E[X\mid\mathcal G]\mid\mathcal H\bigr]
 =\E[X\mid\mathcal H]\quad\text{a.s.}
\end{equation}
特别地，$\E\E[X\mid\mathcal G]=\E X$。
\end{theorem}

\begin{proof}
左端是 $\mathcal H$-可测可积变量。对任意 $H\in\mathcal H\subset\mathcal G$，两次使用
条件期望的积分刻画，得到
\[
 \int_H\E[\E[X\mid\mathcal G]\mid\mathcal H]\,\dd\Prob
 =\int_H\E[X\mid\mathcal G]\,\dd\Prob
 =\int_HX\,\dd\Prob.
\]
所以左端是 $\E[X\mid\mathcal H]$ 的版本，唯一性给出 \eqref{eq:5-2-tower}。令
$\mathcal H=\{\varnothing,\Omega\}$ 即得最后一式。
\end{proof}

\begin{example}[有限分割]\label{ex:5-2-1-1}
设 $A_1,\dots,A_m$ 是 $\Omega$ 的有限可测分割，$\mathcal G=\sigma(A_1,\dots,A_m)$。定义
\[
 Y=\sum_{i:\Prob(A_i)>0}
   \frac{\E[X\1_{A_i}]}{\Prob(A_i)}\1_{A_i},
\]
并在零概率块上任取有限常数。每个 $G\in\mathcal G$ 是若干 $A_i$ 的并；在这些块上逐项求和，
便有 $\int_GY=\int_GX$，所以 $Y=\E[X\mid\mathcal G]$。

回到\cref{ex:5-2-1-prediction-pressure}，奇数块上
\[
 \frac{\E[D^2\1_{A_o}]}{\Prob(A_o)}
 =\frac{(1+9+25)/6}{1/2}=\frac{35}{3},
\]
偶数块上同理为 $56/3$。若把被预测量换成 $D$，两块平均分别为
$(1+3+5)/3=3$ 与 $(2+4+6)/3=4$；“知道奇偶”并不会神奇地把平方赔付也变成 $3$ 与 $4$。
\end{example}

\begin{example}[独立信息]\label{ex:5-2-1-2}
若 $X\in L^1$ 且 $\sigma(X)$ 与 $\mathcal G$ 独立，则
\[
 \E[X\mid\mathcal G]=\E X\quad\text{a.s.}
\]
先对 $\sigma(X)$-可测简单函数 $X$，独立性给
$\E[X\1_G]=\E X\,\Prob(G)$，$G\in\mathcal G$。对一般实值或复值 $X\in L^1$，取
相应简单函数作 $L^1$ 逼近，并用
$|\E[(X-X_n)\1_G]|\leq\|X-X_n\|_1$ 令 $n\to\infty$。常数 $\E X$ 因而满足条件期望的
积分刻画。
\end{example}

\begin{example}[已知变量]\label{ex:5-2-1-3}
若 $X\in L^1$ 本身 $\mathcal G$-可测，则 $X$ 已满足定义中的两项条件，故
$\E[X\mid\mathcal G]=X$ a.s.。塔式性质因此可以理解为：先压缩到 $\mathcal G$，再压缩到更粗的
$\mathcal H$，与直接压缩到 $\mathcal H$ 得到同一结果。
\end{example}

\begin{figure}[htbp]
\centering
\begin{tikzpicture}[node distance=13mm and 18mm]
  \node[book strong node,minimum width=31mm] (x) {$X$\\$\mathcal F$-可测};
  \node[book strong node,right=of x,minimum width=38mm] (y)
    {$Y=\E[X\mid\mathcal G]$\\$\mathcal G$-可测};
  \node[book warm node,below=of y,minimum width=38mm] (tests)
    {$G\in\mathcal G$\\$\displaystyle\int_GY=\int_GX$};
  \draw[book accent arrow] (x)--node[above,book label,align=center,font=\scriptsize]{忘掉不可\\观察细节}(y);
  \draw[book arrow] (y)--(tests);
  \draw[book dashed arrow,bend right=28] (x.south)--node[below left,book label]{在全部可观察事件上保持积分}(tests.west);
\end{tikzpicture}
\caption{条件期望把随机变量压到较粗的信息层，但保留每个可观察事件上的积分。这里的箭头不是点态取值公式。}
\label{fig:5-2-1-information-compression}
\end{figure}

因子取出公式要求因子可观察；若 $Z$ 不是 $\mathcal G$-可测，右端甚至未必
$\mathcal G$-可测。一般无界因子还必须保证 $ZX\in L^1$。同样，版本只逐个 a.s. 唯一；未经
可数协调，不能把一族逐参数等式升级为一个同时成立的点态等式。

当信息层随时间递增时，塔式性质给出逐层压缩的一致性；取因子公式则把可观察量移出条件平均。
这两种机制共同保证逐步信息更新中条件平均计算的相容性。

\begin{exercise}[有限信息上的保费]
设 $D$ 为公平骰子，$X=\1_{\{D\geq4\}}$，$\mathcal G=\sigma(\{D\text{ 为奇数}\})$。
计算 $\E[X\mid\mathcal G]$，并在生成分割的每一块上核对积分。
\end{exercise}

\begin{exercise}[独立信息的截断证明]
从有界 $\sigma(X)$-简单函数开始，使用 $X\wedge n$ 及正负部分，完整证明
\cref{ex:5-2-1-2} 的实值情形。
\end{exercise}

\begin{exercise}[塔式性质与已知因子]
设 $\mathcal H\subset\mathcal G$，$Z$ 有界且 $\mathcal H$-可测。证明
\[
 \E[ZX\mid\mathcal H]=Z\E[X\mid\mathcal H],
 \qquad
 \E[Z\E[X\mid\mathcal G]]=\E[ZX].
\]
\end{exercise}

条件期望为每个测试函数给出一个条件平均。若状态空间是标准 Borel，这些平均可以在一个共同的
可数骨架上协调成随机概率测度；下一小节完成这项构造。


\subsection{正则条件分布、测度分解与 Markov 核}\label{subsec:5-2-2}
\index{正则条件分布}\index{测度分解}\index{后验核}

假设 $X$ 与 $Y$ 取离散值。给定 $Y=y$ 时，$X$ 的条件分布是一族概率
$A\mapsto\Prob(X\in A\mid Y=y)$；若 $Y$ 连续，这一条件概率族仍应存在，只是不能再由
$\Prob(Y=y)$ 作分母。我们需要的是一个关于 $y$ 可测的概率测度族。第~3~章已经把这样的族称为
概率核；现在的问题是，条件期望给出的逐事件等价类能否同时选成一个核。

下面两个计算先给出候选，并显示抽象结论必须恢复哪些熟悉公式。

\begin{example}[离散条件分布]\label{ex:5-2-2-1}
设 $Y$ 的值域包含于标准 Borel 空间 $T$ 的可数集 $C=\{y_1,y_2,\dots\}$，而
$X$ 取值于可测空间 $(S,\mathcal S)$。固定一个 $S$ 上的概率测度 $\mu_0$。定义
\begin{equation}\label{eq:5-2-discrete-kernel}
 k(y,A)=
 \begin{cases}
 \displaystyle\frac{\Prob(X\in A,Y=y)}{\Prob(Y=y)},
   &y\in C,\ \Prob(Y=y)>0,\\[7pt]
 \mu_0(A),&\text{其他情形}.
 \end{cases}
\end{equation}
标准 Borel 空间中的单点是 Borel 集，所以对每个 $A$，函数 $y\mapsto k(y,A)$ 是一个只在
可数集上偏离常数的 Borel 函数。固定 $y$ 时，$A\mapsto k(y,A)$ 是概率测度，故 $k$ 是概率核。
对任意 Borel 集 $B\subset T$，
\[
 \begin{aligned}
 \int_{\{Y\in B\}}k(Y,A)\,\dd\Prob
 &=\sum_{j:y_j\in B}\Prob(Y=y_j)k(y_j,A)\\
 &=\sum_{j:y_j\in B}\Prob(X\in A,Y=y_j)\\
 &=\Prob(X\in A,Y\in B).
 \end{aligned}
\]
零概率状态上的 $\mu_0$ 没有进入任何一项，说明这些状态的核值不由联合分布决定。
\end{example}

\begin{example}[联合密度]\label{ex:5-2-2-2}
设 $X\in\R^d$、$Y\in\R^m$ 的联合分布关于 Lebesgue 测度有密度 $p(x,y)$，并令
\[
 p_Y(y)=\int_{\R^d}p(x,y)\,\dd x.
\]
Tonelli 定理说明 $p_Y$ 可测。取任一概率密度 $r$ 于 $\R^d$，定义
\begin{equation}\label{eq:5-2-density-kernel}
 k(y,\dd x)=
 \begin{cases}
 \dfrac{p(x,y)}{p_Y(y)}\,\dd x,&p_Y(y)>0,\\[6pt]
 r(x)\,\dd x,&p_Y(y)=0.
 \end{cases}
\end{equation}
固定 Borel 集 $A$ 后，参数积分的可测性说明 $y\mapsto k(y,A)$ 可测。对 Borel 集
$B\subset\R^m$，Fubini--Tonelli 计算给出
\[
 \begin{aligned}
 \int_{\{Y\in B\}}k(Y,A)\,\dd\Prob
 &=\int_B k(y,A)p_Y(y)\,\dd y\\
 &=\int_B\int_Ap(x,y)\,\dd x\,\dd y
 =\Prob(X\in A,Y\in B).
 \end{aligned}
\]
因此密度比确实给出条件核；它是一般定义的一个后果，而不是一般定义本身。
\end{example}

\begin{remark}[标准 Borel 假设的边界]\label{ex:5-2-2-3}
若目标可测空间不是标准 Borel，正则条件分布可能不存在。本书不在这里构造所需的病态可测空间，
但下面的存在证明会逐处标出标准 Borel 假设的用途：它既提供可数生成的半直线骨架，又提供像为
Borel 集且逆映射可测的编码。只知道某个 $\sigma$-代数可数生成，或只知道一个可测函数族能够分离
点，并不能自动提供这样的双向编码。条件信息 $\mathcal G$ 是否可数生成则不是本定理的假设。
\end{remark}

\begin{example}[确定性条件核]\label{ex:5-2-2-4}
若 $X=h(Y)$，其中 $h:T\to S$ 可测，则
\[
 k(y,\cdot)=\delta_{h(y)}
\]
是从 $T$ 到 $S$ 的概率核。对 $A\in\mathcal S$，
$k(Y,A)=\1_A(h(Y))=\1_{\{X\in A\}}$，故它已经是相应指标函数关于 $Y$ 的条件期望；并且
对非负 Borel 函数 $f$，
\[
 \int_Sf(x)k(y,\dd x)=f(h(y)).
\]
这个公式在 $\Prob(Y=y)=0$ 时仍有确定含义，不过在 $Y$ 的分布零集上改动整个核也不会改变任何
条件积分。
\end{example}

\begin{definition}[正则条件分布]\label{def:5-2-2-1}
设 $X:\Omega\to(S,\mathcal S)$ 是随机元，$\mathcal G\subset\mathcal F$。若从
$(\Omega,\mathcal G)$ 到 $(S,\mathcal S)$ 的概率核 $K$ 满足：对每个 $A\in\mathcal S$，
\begin{equation}\label{eq:5-2-rcd-def}
 K(\cdot,A)=\E[\1_{\{X\in A\}}\mid\mathcal G]
 \quad\text{a.s.},
\end{equation}
则称 $K$ 是 $X$ 给定 $\mathcal G$ 的一个\emph{正则条件分布}，记作
$K=\Law(X\mid\mathcal G)$。
\end{definition}

\begin{definition}[给定随机变量的核]\label{def:5-2-2-2}
设 $S,T$ 都是标准 Borel 空间，$X:\Omega\to S$、$Y:\Omega\to T$。若从 $T$ 到 $S$ 的
概率核 $k$ 满足
\[
 k(Y(\omega),A)
 =\E[\1_{\{X\in A\}}\mid\sigma(Y)](\omega)
 \quad\text{a.s.}
\]
对每个 $A$ 成立，则把 $k(y,\cdot)$ 记作 $\Law(X\mid Y=y)$。这个点值记号包含了版本选择；
它不表示对零概率事件使用比值。
\end{definition}

\begin{definition}[测度分解]\label{def:5-2-2-3}
设 $\pi$ 是 $(S\times T,\mathcal S\otimes\mathcal T)$ 上的概率，第一边缘为 $\mu$。若从
$S$ 到 $T$ 的概率核 $K$ 满足
\begin{equation}\label{eq:5-2-disintegration-rectangles}
 \pi(A\times B)=\int_AK(x,B)\,\mu(\dd x)
 \qquad(A\in\mathcal S,\ B\in\mathcal T),
\end{equation}
就把它写作
\[
 \pi(\dd x,\dd y)=\mu(\dd x)K(x,\dd y),
\]
并称这是 $\pi$ 相对于第一边缘的\emph{测度分解}（disintegration）。
\end{definition}

逐个 $A$ 选择 \eqref{eq:5-2-rcd-def} 的版本不够：$\mathcal S$ 通常不可数，而每个版本各有一个
例外零集。若只在一个可数集合代数上选版本，再声称逐点作 \Caratheodory{} 扩张，也仍有缺口，
因为这个代数拥有不可数多条下降序列，不能用一个可数零集保证逐点可数可加。实线上的分布函数把
可数协调压缩为有理阈值间的单调性与右连续关系，正好绕开这个障碍。

\begin{theorem}[正则条件分布的存在性]\label{theorem:5-2-2-1}
设 $(S,\mathcal S)$ 是标准 Borel 空间，$X:\Omega\to S$。对任意子 $\sigma$-代数
$\mathcal G\subset\mathcal F$，都存在 $\Law(X\mid\mathcal G)$ 的正则版本；这里不要求
$\mathcal G$ 可数生成。

若 $Y:\Omega\to T$ 还取值于标准 Borel 空间，则存在从 $T$ 到 $S$ 的概率核 $k$，使
$k(Y,\cdot)$ 是 $\Law(X\mid\sigma(Y))$ 的正则版本。
\end{theorem}

\begin{proof}
\emph{第一步：把目标编码进实线。}
由标准 Borel 编码定理~\ref{theorem:1-4-2-4}，存在 Borel 集 $B\subset[0,1]$ 与 Borel
同构 $\iota:S\to B$。令 $Z=\iota(X)$。对每个 $q\in\Q$，选择一个版本
\[
 f_q=\E[\1_{\{Z\leq q\}}\mid\mathcal G],
\]
并把它截到 $[0,1]$；截断只改变零测集。

若 $q<r$，条件期望的单调性给 $f_q\leq f_r$ a.s.。若 $q<0$，则 $f_q=0$ a.s.；若
$q\geq1$，则 $f_q=1$ a.s.。还需协调右连续性。对固定 $q$，取有理数
$r_{q,n}\downarrow q$。删去单调性例外集后，$f_{r_{q,n}}$ 逐点非增；记其极限为 $h_q$。对
$G\in\mathcal G$，支配收敛定理和测度从上连续性给出
\[
 \begin{aligned}
 \int_Gh_q\,\dd\Prob
 &=\lim_n\int_Gf_{r_{q,n}}\,\dd\Prob
 =\lim_n\Prob(G\cap\{Z\leq r_{q,n}\})\\
 &=\Prob(G\cap\{Z\leq q\}).
 \end{aligned}
\]
所以 $h_q$ 与 $f_q$ 是同一个条件期望版本。把 $\{f_q\ne\lim_n f_{r_{q,n}}\}$、全部有理对
$q<r$ 的单调性失败集
以及端点关系的失败集取可数并，得到一个 $N\in\mathcal G$、$\Prob(N)=0$。在
$N$ 上把
$f_q$ 全部改成 $\1_{\{q\geq0\}}$；于是对每个 $\omega$ 都有
\begin{equation}\label{eq:5-2-rational-cdf-relations}
 0\leq f_q(\omega)\leq f_r(\omega)\leq1\quad(q<r),
 \qquad
 f_q(\omega)=\inf_{r>q,\ r\in\Q}f_r(\omega),
\end{equation}
并且 $f_q=0$ 于 $q<0$、$f_q=1$ 于 $q\geq1$。
最后一个等式使用了 $(f_r)_{r\in\Q}$ 的单调性：对任意有理数 $r>q$，最终都有
$r_{q,n}<r$，故沿 $(r_{q,n})$ 的极限正是对全部有理数 $r>q$ 所取的下确界。

\emph{第二步：逐点生成概率测度。}
对 $t\in\R$ 定义
\begin{equation}\label{eq:5-2-cdf-extension}
 F_\omega(t)=\inf_{q>t,\ q\in\Q}f_q(\omega).
\end{equation}
它非减。若 $t_n\downarrow t$，则 $F_\omega(t_n)\geq F_\omega(t)$；另一方面，对任意有理
$q>t$，充分大的 $n$ 满足 $t_n<q$，因而
$F_\omega(t_n)\leq f_q(\omega)$。先令 $n\to\infty$，再对 $q>t$ 取下确界，得到
$F_\omega(t_n)\downarrow F_\omega(t)$。所以 $F_\omega$ 右连续。端点关系还给出
\[
 F_\omega(t)=0\quad(t<0),
 \qquad F_\omega(t)=1\quad(t\geq1),
\]
且由 \eqref{eq:5-2-rational-cdf-relations}，$F_\omega(q)=f_q(\omega)$ 对每个有理 $q$ 成立。
Lebesgue--Stieltjes 构造~\ref{theorem:2-3-3-2} 因而给出实线上的概率测度
$L(\omega,\cdot)$，其分布函数为 $F_\omega$，并且每个 $L(\omega,\cdot)$ 集中在 $[0,1]$。

\emph{第三步：验证参数可测性与条件恒等式。}
对每个 $t$，式 \eqref{eq:5-2-cdf-extension} 是可数个 $\mathcal G$-可测函数的下确界，故
$\omega\mapsto L(\omega,(-\infty,t])$ 可测。令
\[
 \mathcal C=\{A\in\Borel(\R):\omega\mapsto L(\omega,A)
 \text{ 是 }\mathcal G\text{-可测的}\}.
\]
$\mathcal C$ 含实线，对补集封闭，并对两两不交的可数并封闭；后一项由逐点测度可加性和可测函数
部分和的极限得到。因此 $\mathcal C$ 是 $\lambda$-系统。它包含所有有理端点左半直线，这些集合
构成生成 $\Borel(\R)$ 的 $\pi$-系统。$\pi$--$\lambda$ 定理给
$\mathcal C=\Borel(\R)$，故 $L$ 是概率核。

再令
\[
 \mathcal D=\{A\in\Borel(\R):L(\cdot,A)
 \text{ 是 }\E[\1_{\{Z\in A\}}\mid\mathcal G]\text{ 的版本}\}.
\]
由 $L(\cdot,(-\infty,q])=f_q$，所有有理端点左半直线属于 $\mathcal D$。常数一处理全空间；
线性处理补集；若 $A_n$ 两两不交且属于 $\mathcal D$，则对 $G\in\mathcal G$，单调收敛给
\[
 \int_G L\!\left(\cdot,\bigcup_nA_n\right)\,\dd\Prob
 =\sum_n\int_GL(\cdot,A_n)\,\dd\Prob
 =\sum_n\Prob(G\cap\{Z\in A_n\}).
\]
所以 $\mathcal D$ 也是 $\lambda$-系统。再次使用 $\pi$--$\lambda$ 定理，得到
$\mathcal D=\Borel(\R)$。

特别地，$L(\cdot,B)$ 是 $\E[\1_{\{Z\in B\}}\mid\mathcal G]=1$ 的版本。故
$N_B=\{L(\cdot,B)\ne1\}\in\mathcal G$ 是零测集。固定 $b_0\in B$，在 $N_B$ 上把
$L(\omega,\cdot)$ 改成 $\delta_{b_0}$。修改后的核在每个 $\omega$ 都集中于 $B$，并且仍满足
全部条件恒等式。最后定义
\[
 K(\omega,A)=L(\omega,\iota(A)),\qquad A\in\mathcal S.
\]
因为 $\iota$ 及其逆映射都可测，$\iota(A)$ 是 $B$ 中、从而是 $\R$ 中的 Borel 集；所以 $K$
是从 $(\Omega,\mathcal G)$ 到 $S$ 的概率核，并满足 \eqref{eq:5-2-rcd-def}。

\emph{第四步：给定 $Y=y$ 的版本。}
现在令 $\mathcal G=\sigma(Y)$。在第一步只出现可数个函数 $f_q$。对每个 $q$ 使用
Doob--Dynkin 因子化，先写成 $f_q=h_q(Y)$；因 $f_q\in[0,1]$ a.s.，再把 $h_q$ 截断到
$[0,1]$，仍保持该等式 a.s. 成立。于是可设 $h_q:T\to[0,1]$ Borel 可测。记
$\nu=\Law(Y)$，并固定 $b_0\in B$。由 $f_q$ 已经满足
\eqref{eq:5-2-rational-cdf-relations}，下列 Borel 集的 $\nu$-测度都为零：
\begin{align*}
 &\{y:h_q(y)>h_r(y)\}\quad(q<r),\\
 &\left\{y:h_q(y)\ne\inf_{r>q,\ r\in\Q}h_r(y)\right\}\quad(q\in\Q),\\
 &\{y:h_q(y)\ne0\}\quad(q<0),
 \qquad
 \{y:h_q(y)\ne1\}\quad(q\ge1).
\end{align*}
把它们的可数并记为 $E$。在 $E$ 上同时改定
\[
 h_q(y)=\1_{\{b_0\le q\}},
\]
则 $h_q(Y)=f_q$ 仍对所有有理数 $q$ 同时 a.s. 成立，且上述单调、端点和右连续关系
对每个 $y\in T$ 成立。

对 $y\in T$、$t\in\R$ 定义
\[
 F_y(t)=\inf_{q>t,\ q\in\Q}h_q(y).
\]
第二步的逐式论证说明 $F_y$ 非降、右连续，且
$F_y(-\infty)=0$、$F_y(+\infty)=1$。因此它唯一决定概率测度
$\ell(y,\cdot)$，并有
\[
 \ell(y,(-\infty,q])=h_q(y)\qquad(q\in\Q).
\]
对每个实数 $t$，$y\mapsto F_y(t)$ 是可数个 Borel 函数的下确界。把使
$y\mapsto\ell(y,A)$ 可测的 Borel 集 $A\subset\R$ 组成 $\lambda$-系统，并从有理端点左半直线出发使用
$\pi$--$\lambda$ 定理，可得 $\ell$ 是从 $T$ 到 $\R$ 的概率核。另一次同样的
$\lambda$-系统计算由等式
\[
 \ell(Y,(-\infty,q])=f_q\quad\text{a.s.}
\]
出发，得到对每个 $A\in\Borel(\R)$，
\[
 \ell(Y,A)=\E[\1_{\{Z\in A\}}\mid\sigma(Y)]
 \quad\text{a.s.}
\]
特别地，$\ell(Y,B)=1$ a.s.，所以
$E_B=\{y:\ell(y,B)\ne1\}$ 满足 $\nu(E_B)=0$。在 $E_B$ 上把整个测度
$\ell(y,\cdot)$ 改为 $\delta_{b_0}$，得到的核对每个 $y$ 都集中于 $B$，且不改变任何一条 a.s. 条件恒等式。
最后定义
\[
 k(y,A)=\ell(y,\iota(A)),\qquad A\in\mathcal S.
\]
则 $k$ 是所需的从 $T$ 到 $S$ 的概率核。这里因子化的是有理阈值的可数族，而不是尚未构造的“整个核”。
\end{proof}

\begin{remark}[正则条件分布构造的三个步骤]\label{rem:5-2-2-existence-checkpoints}
上述构造的三个步骤各有独立作用：有理阈值把版本选择压成可数问题；逐点分布函数保证每个参数处得到的
确是可数可加概率；Borel 同构的可测逆把实线上的核安全地送回原状态空间。

在\cref{ex:5-2-2-1}中，有理阈值条件期望就是每个正概率分割块上的离散累计概率，生成的测度恢复
\eqref{eq:5-2-discrete-kernel}。在\cref{ex:5-2-2-2}中，核在每个左半直线上的值由
\eqref{eq:5-2-density-kernel} 积分得到；测度唯一性因而恢复密度比。标准 Borel 假设作用在目标
空间 $S$，不是条件信息 $\mathcal G$。
\end{remark}

一旦核已经存在，所有 Borel 测试函数可以由指标函数和单调收敛统一处理。

\begin{theorem}[核表示条件期望]\label{theorem:5-2-2-2}
若 $K=\Law(X\mid\mathcal G)$，则对每个非负 Borel 函数 $f:S\to[0,\infty]$，
\begin{equation}\label{eq:5-2-kernel-ce}
 \E[f(X)\mid\mathcal G](\omega)
 =\int_Sf(x)K(\omega,\dd x)\quad\text{a.s.}
\end{equation}
两边允许为 $+\infty$。若 $f$ 为实值或复值且 $f(X)\in L^1$，同一公式仍成立；对 a.e.
$\omega$，函数 $f$ 关于 $K(\omega,\cdot)$ 绝对可积。
\end{theorem}

\begin{proof}
对 $f=\1_A$，公式就是正则条件分布的定义。非负简单函数情形由有限线性性得到。对一般
$f\geq0$，取非负简单函数 $s_n\uparrow f$。固定 $\omega$ 后，对 $K(\omega,\cdot)$ 使用单调
收敛；同时对扩展条件期望使用 \eqref{eq:5-2-conditional-mct}，得到
\[
 \int s_n\,\dd K(\omega,\cdot)\uparrow\int f\,\dd K(\omega,\cdot),
 \qquad
 \E[s_n(X)\mid\mathcal G]\uparrow\E[f(X)\mid\mathcal G]
\]
于同一个可数例外集之外，故 \eqref{eq:5-2-kernel-ce} 成立。

若 $f(X)\in L^1$，先把非负结论用于 $|f|$ 并取期望：
\[
 \E\int_S|f(x)|K(\cdot,\dd x)=\E|f(X)|<\infty.
\]
所以核积分对 a.e. $\omega$ 绝对收敛。实值函数分成正负部，复值函数再分实部、虚部，利用条件
期望线性即可。
\end{proof}

\begin{theorem}[测度分解定理]\label{theorem:5-2-2-3}
设 $S,T$ 是标准 Borel 空间，$\pi$ 是 $S\times T$ 上的概率，第一边缘为 $\pi_S$。则存在从
$S$ 到 $T$ 的概率核 $K$，使
\[
 \pi(\dd x,\dd y)=\pi_S(\dd x)K(x,\dd y).
\]
若 $K'$ 也是这样的核，则
$K(x,\cdot)=K'(x,\cdot)$ 对 $\pi_S$-a.e. $x$ 成立。
\end{theorem}

\begin{proof}
在规范概率空间 $(S\times T,\mathcal S\otimes\mathcal T,\pi)$ 上令
$U(x,y)=x$、$V(x,y)=y$。由\cref{theorem:5-2-2-1}及其中的因子化部分，存在从 $S$ 到
$T$ 的概率核 $K$，使
$K(U,B)=\E[\1_{\{V\in B\}}\mid\sigma(U)]$ a.s.。对 $A\in\mathcal S$、
$B\in\mathcal T$，已知因子公式与塔式性质给出
\[
 \begin{aligned}
 \pi(A\times B)
 &=\E_\pi[\1_A(U)\1_B(V)]\\
 &=\E_\pi[\1_A(U)K(U,B)]
 =\int_AK(x,B)\,\pi_S(\dd x),
 \end{aligned}
\]
故矩形公式成立。

为证唯一性，由标准 Borel 编码取 $T$ 的一个可数生成 $\pi$-系统 $\mathcal P$。若 $K,K'$ 都满足
矩形公式，则对每个 $B\in\mathcal P$ 与每个 $A\in\mathcal S$，
\[
 \int_AK(x,B)\,\pi_S(\dd x)=\int_AK'(x,B)\,\pi_S(\dd x).
\]
Radon--Nikodym 密度的唯一性说明 $K(\cdot,B)=K'(\cdot,B)$ a.e.。对可数个
$B\in\mathcal P$ 合并零集；在其补集上，两个概率测度在生成 $\pi$-系统上相等，有限测度唯一性
定理~\ref{theorem:2-2-1-2} 便给出它们在 $\mathcal T$ 上处处相等。
\end{proof}

给定一个边缘和一个概率核，反向确实可以构造联合测度，而不仅是形式地写两个积分号。

\begin{proposition}[核重构联合分布]\label{proposition:5-2-2-4}
设 $X:\Omega\to(S,\mathcal S)$、$Y:\Omega\to(T,\mathcal T)$ 是随机元，
$\mu=\Law(X)$，并设从 $S$ 到 $T$ 的概率核 $K$ 满足：对每个
$B\in\mathcal T$，
\[
 K(X,B)=\E[\1_{\{Y\in B\}}\mid\sigma(X)]\quad\text{a.s.}
\]
则
$(X,Y)$ 的联合分布为 $\mu K$，即
\begin{equation}\label{eq:5-2-joint-from-kernel}
 \Prob(X\in A,Y\in B)=\int_AK(x,B)\,\mu(\dd x)
 \qquad(A\in\mathcal S,\ B\in\mathcal T).
\end{equation}
对每个非负或关于联合分布可积的可测函数 $f$，
\begin{equation}\label{eq:5-2-kernel-fubini}
 \E f(X,Y)=\int_S\int_T f(x,y)K(x,\dd y)\mu(\dd x).
\end{equation}
\end{proposition}

\begin{proof}
先说明 $\mu K$ 的含义。对 $E\in\mathcal S\otimes\mathcal T$，令 $E_x=\{y:(x,y)\in E\}$。
满足 $x\mapsto K(x,E_x)$ 可测的集合 $E$ 构成一个 $\lambda$-系统：全空间对应常数一，补集对应
$1-K(x,E_x)$，两两不交可数并对应非负可测函数部分和的极限。它包含全部矩形，因为
\[
 K(x,(A\times B)_x)=\1_A(x)K(x,B).
\]
矩形是生成积 $\sigma$-代数的 $\pi$-系统，故 $x\mapsto K(x,E_x)$ 对每个积可测 $E$ 都可测。
定义
\[
 (\mu K)(E)=\int_SK(x,E_x)\,\mu(\dd x).
\]
若 $E_n$ 两两不交，则其每个 $x$-截面也两两不交；核的可数可加性和单调收敛定理说明
$\mu K$ 可数可加。又 $(\mu K)(S\times T)=1$，故这是联合概率。指标函数、非负简单函数和单调
收敛依次给出
\[
 \int f\,\dd(\mu K)=\int_S\int_Tf(x,y)K(x,\dd y)\mu(\dd x)
\]
对全部 $f\geq0$ 成立。

现在利用条件分布。对矩形 $A\times B$，
\[
 \begin{aligned}
 \Prob(X\in A,Y\in B)
 &=\E\bigl[\1_A(X)\E[\1_B(Y)\mid X]\bigr]\\
 &=\E[\1_A(X)K(X,B)]
 =\int_AK(x,B)\,\mu(\dd x).
 \end{aligned}
\]
两个概率测度在生成矩形 $\pi$-系统上相等，故联合分布等于 $\mu K$。非负
\eqref{eq:5-2-kernel-fubini} 已由上段得到；有符号或复值情形先把公式用于 $|f|$，确认两次积分
绝对有限，再分正负部或实虚部。
\end{proof}

\begin{proposition}[核版本的共同零集唯一性]\label{proposition:5-2-2-5}
设 $S$ 为标准 Borel 空间。若 $K,K'$ 都是 $\Law(X\mid\mathcal G)$ 的正则版本，则存在
$N\in\mathcal G$、$\Prob(N)=0$，使对每个 $\omega\notin N$，
\[
 K(\omega,\cdot)=K'(\omega,\cdot)
\]
作为 $S$ 上的概率测度成立。
\end{proposition}

\begin{proof}
标准 Borel 编码给出 $\mathcal S$ 的一个可数生成 $\pi$-系统 $\mathcal P$。对每个
$A\in\mathcal P$，$K(\cdot,A)$ 与 $K'(\cdot,A)$ 都是同一个条件期望的版本，故二者在某个零集
外相等。合并这些可数个零集得到 $N\in\mathcal G$、$\Prob(N)=0$。固定
$\omega\notin N$，两个概率测度 $K(\omega,\cdot)$、$K'(\omega,\cdot)$ 在
$\mathcal P$ 上相等；有限测度唯一性把等式推广到全部 $A\in\mathcal S$。

这个共同零集不能删除。\cref{ex:5-2-2-4}中，在 $Y$ 的分布零集上把
$\delta_{h(y)}$ 换成任一概率测度，仍得到同一个条件分布版本。
\end{proof}

\begin{figure}[htbp]
\centering
\begin{tikzpicture}[node distance=13mm and 18mm]
  \node[book strong node,minimum width=27mm] (mu) {第一边缘\\$\mu(\dd x)$};
  \node[book strong node,right=of mu,minimum width=34mm] (cloud)
    {给定 $x$ 的条件概率\\$K(x,\dd y)$};
  \node[book warm node,right=of cloud,minimum width=34mm] (joint)
    {联合概率\\$\mu(\dd x)K(x,\dd y)$};
  \draw[book accent arrow] (mu)--node[above,book label]{指定条件核}(cloud);
  \draw[book accent arrow] (cloud)--node[above,book label]{对 $x$ 积分}(joint);
  \foreach \p in {-.34,0,.34}{
    \fill[BookAccent!72] ($(cloud.center)+(0.72,\p)$) circle (1.25pt);
  }
  \draw[book dashed arrow,bend right=28] (joint.south)--node[below,book label]{取第一边缘}(mu.south);
\end{tikzpicture}
\caption{测度分解先保留第一边缘，再在每个 $x$ 上附加条件概率；把这些概率测度积分起来便重构联合分布。}
\label{fig:5-2-2-disintegration-clouds}
\end{figure}

正则条件分布只在边缘 a.e. 意义下唯一，符号 $\Law(X\mid Y=y)$ 因而必须连同版本理解。
本小节的存在定理依赖目标状态空间标准 Borel；它没有断言任意可测空间都存在条件核。连续时间路径
分布的逐步构造还需要后面的扩张定理，此处不使用也不预设 Ionescu--Tulcea 定理。

\begin{exercise}[离散核的零概率行]
在\cref{ex:5-2-2-1}中把零概率状态上的 $\mu_0$ 换成另一概率 $\mu_1$。证明两个核给出同一个
关于 $Y$ 的条件分布版本，并找出它们可能不同的全部状态。
\end{exercise}

\begin{exercise}[密度模型的测度分解]
从 \eqref{eq:5-2-density-kernel} 出发，证明
\[
 p(x,y)\,\dd x\,\dd y
 =p_Y(y)\,\dd y\,k(y,\dd x),
\]
并把公式推广到任意非负 Borel 测试函数。
\end{exercise}

\begin{exercise}[核的共同零集]
完整写出\cref{proposition:5-2-2-5}中可数生成 $\pi$-系统：先把 $S$ Borel 编码为
$B\subset[0,1]$，再使用有理端点左半直线与 $B$ 的交。
\end{exercise}

条件核给出每个状态下的完整分布，条件期望则选取其中的积分。下一小节把这种积分预测同凸性、
平方误差几何和贝叶斯换测度公式联系起来。


\subsection{条件 Jensen、\texorpdfstring{$L^2$}{L2} 投影与贝叶斯公式}\label{subsec:5-2-3}
\index{条件 Jensen 不等式}\index{最小均方预测}\index{贝叶斯公式}

在\cref{ex:5-2-1-prediction-pressure}中，两块上的条件平均恰好最小化平方误差。这个现象不是有限
分割的巧合。条件期望一方面受凸性控制，另一方面在 $L^2$ 中与所有可观察方向正交。贝叶斯公式则
回答另一个表面上不同的问题：更换概率测度以后，同一信息层上的条件平均怎样改变？三者最终都只
使用条件期望的积分刻画。

\begin{definition}[信息子空间]\label{def:5-2-3-1}
定义
\[
 L^2(\mathcal G)=\{Y\in L^2(\Prob):Y\text{ 有一个 }\mathcal G\text{-可测版本}\}.
\]
\end{definition}

这是 $L^2(\mathcal F)$ 的闭线性子空间。线性性直接成立。为证闭性，设
$Y_n\in L^2(\mathcal G)$ 且
$Y_n\to Y$ 于 $L^2$，为每个 $Y_n$ 选取一个 $\mathcal G$-可测版本，再取子列使
\[
 \|Y_{n_k}-Y\|_2^2\leq 2^{-3k}\qquad(k\geq1).
\]
Markov 不等式给出
\[
 \sum_{k=1}^\infty
 \Prob\bigl(|Y_{n_k}-Y|>2^{-k}\bigr)
 \leq\sum_{k=1}^\infty2^{-k}<\infty.
\]
由 Borel--Cantelli I，$Y_{n_k}\to Y$ a.s.。令
\[
 C=\{\omega:(Y_{n_k}(\omega))_{k\geq1}\text{ 在 }\C\text{ 中为 Cauchy 列}\}.
\]
$C\in\mathcal G$，因为 Cauchy 条件可写成关于 $|Y_{n_j}-Y_{n_k}|$ 的可数次交并；函数
\[
 \widetilde Y(\omega)=
 \begin{cases}
  \lim_{k\to\infty}Y_{n_k}(\omega),&\omega\in C,\\
  0,&\omega\notin C
 \end{cases}
\]
是 $\mathcal G$-可测，并且 $\widetilde Y=Y$ a.s.，故 $Y\in L^2(\mathcal G)$。这个论证同时
适用于实值与复值 $L^2$。

要把 $\E[X\mid\mathcal G]$ 解释为这个子空间上的投影，还须验证它保持平方可积。
\cref{theorem:5-2-3-1} 中取 $\varphi(u)=u^2$ 将给出所需的 $L^2$ 收缩估计。

\begin{definition}[后验核]\label{def:5-2-3-2}
设参数空间 $\Theta$ 与观测空间 $T$ 为标准 Borel 空间，$\pi(\dd\theta)$ 是先验概率，
$L(\theta,\dd y)$ 是似然核。联合分布
\[
 \pi(\dd\theta)L(\theta,\dd y)
\]
相对于观测边缘的测度分解核 $\pi(\dd\theta\mid y)$ 称为\emph{后验核}。它在观测边缘
a.e. 意义下唯一；零边缘质量的观测值上仍需选择版本。
\end{definition}

\begin{definition}[换测度]\label{def:5-2-3-3}
设 $\mathbb Q$ 是 $(\Omega,\mathcal F)$ 上的概率且 $\mathbb Q\ll\Prob$，记
\[
 Z=\frac{\dd\mathbb Q}{\dd\Prob}.
\]
称这是从 $\Prob$ 到 $\mathbb Q$ 的密度或似然比。对信息层 $\mathcal G$，记
\[
 Z_{\mathcal G}=\E_\Prob[Z\mid\mathcal G].
\]
\end{definition}

\begin{proposition}[限制信息上的密度]
在定义~\ref{def:5-2-3-3} 的条件下，$Z_{\mathcal G}$ 是
$\mathbb Q|_{\mathcal G}$ 关于 $\Prob|_{\mathcal G}$ 的 Radon--Nikodym 导数。
\end{proposition}

\begin{proof}
对每个 $G\in\mathcal G$，
\[
 \mathbb Q(G)=\int_GZ\,\dd\Prob=\int_GZ_{\mathcal G}\,\dd\Prob,
\]
而 $Z_{\mathcal G}$ 非负且 $\mathcal G$-可测。这正是所述 Radon--Nikodym 导数的定义。
\end{proof}

先建立凸性估计。证明中只用一维凸函数的支撑直线，并把不可数条直线缩成可数族。

\begin{theorem}[条件 Jensen]\label{theorem:5-2-3-1}
设 $\varphi:\R\to\R$ 为有限凸函数。若实值 $X$ 满足
$X,\varphi(X)\in L^1$，则
\begin{equation}\label{eq:5-2-conditional-jensen}
 \varphi\bigl(\E[X\mid\mathcal G]\bigr)
 \leq\E[\varphi(X)\mid\mathcal G]
 \quad\text{a.s.}
\end{equation}
左端也属于 $L^1$。
\end{theorem}

\begin{proof}
有限凸函数在实线上连续。对每个 $r\in\Q$，凸割线斜率的单调性给出一个实数 $a_r$，使
\begin{equation}\label{eq:5-2-convex-support}
 \ell_r(u):=\varphi(r)+a_r(u-r)\leq\varphi(u)
 \qquad(u\in\R).
\end{equation}
例如可在左、右导数形成的闭区间
$[\varphi'_-(r),\varphi'_+(r)]$ 中选取 $a_r$；
\eqref{eq:5-2-convex-support} 直接来自左右割线斜率的次序。

这个可数族恢复 $\varphi$：
\begin{equation}\label{eq:5-2-countable-affine-envelope}
 \varphi(u)=\sup_{r\in\Q}\ell_r(u).
\end{equation}
为证反向不等式，固定 $u$ 并取有理数 $r_n\to u$。当 $r_n\in(u-1,u+1)$ 时，凸割线斜率的
单调性给出明确的界
\[
 \varphi(u-1)-\varphi(u-2)
 \leq a_{r_n}
 \leq \varphi(u+2)-\varphi(u+1).
\]
故 $(a_{r_n})$ 有界。由 $\varphi(r_n)\to\varphi(u)$，
\[
 \ell_{r_n}(u)=\varphi(r_n)+a_{r_n}(u-r_n)\longrightarrow\varphi(u).
\]
结合每条支撑线均不超过 $\varphi$，便得 \eqref{eq:5-2-countable-affine-envelope}。

令 $Y=\E[X\mid\mathcal G]$。对每个 $r\in\Q$，线性与正性给出
\[
 \ell_r(Y)
 =\E[\ell_r(X)\mid\mathcal G]
 \leq\E[\varphi(X)\mid\mathcal G]
 \quad\text{a.s.}
\]
合并可数个例外集，再对 $r$ 取上确界并使用
\eqref{eq:5-2-countable-affine-envelope}，得到
\eqref{eq:5-2-conditional-jensen}。最后固定任一支撑线 $\ell_{r_0}$。我们已经得到
\[
 \ell_{r_0}(Y)\leq\varphi(Y)\leq\E[\varphi(X)\mid\mathcal G].
\]
左右两端都可积。更明确地，
$\varphi(Y)^+\leq\E[\varphi(X)\mid\mathcal G]^+$ 且
$\varphi(Y)^-\leq\ell_{r_0}(Y)^-$，所以 $\varphi(Y)\in L^1$。
\end{proof}

\begin{remark}[从凸性进入平方几何]\label{rem:5-2-3-l2-entry}
若实值 $X\in L^2$，在\cref{theorem:5-2-3-1}中取 $\varphi(u)=u^2$，得到
\begin{equation}\label{eq:5-2-l2-contraction}
 \bigl|\E[X\mid\mathcal G]\bigr|^2
 \leq\E[|X|^2\mid\mathcal G],
 \qquad
 \|\E[X\mid\mathcal G]\|_2\leq\|X\|_2.
\end{equation}
所以 $P_{\mathcal G}X:=\E[X\mid\mathcal G]$ 确实属于
$L^2(\mathcal G)$。对\cref{ex:5-2-1-prediction-pressure}，
$P_{\mathcal G}(D^2)$ 正是奇数块上的 $35/3$ 与偶数块上的 $56/3$。残差同每个可观察的
$L^2$ 方向正交后，平方误差的最小化性质将直接由勾股恒等式得到。
\end{remark}

复值 $X\in L^2$ 也满足 \eqref{eq:5-2-l2-contraction}：由
$|\E[X\mid\mathcal G]|\leq\E[|X|\mid\mathcal G]$，再把实值条件 Jensen 用于非负变量
$|X|$ 与函数 $u\mapsto u^2$ 即得。

\begin{proposition}[条件 \Holder{}]\label{proposition:5-2-3-holder}
设 $1<p,q<\infty$、$p^{-1}+q^{-1}=1$，$U\in L^p$、$V\in L^q$。则
\begin{equation}\label{eq:5-2-conditional-holder}
 \bigl|\E[UV\mid\mathcal G]\bigr|
 \leq
 \E[|U|^p\mid\mathcal G]^{1/p}
 \E[|V|^q\mid\mathcal G]^{1/q}
 \quad\text{a.s.}
\end{equation}
\end{proposition}

\begin{proof}
普通 \Holder{} 不等式~\ref{theorem:3-3-2-2} 先给 $UV\in L^1$。记
\[
 A=\E[|U|^p\mid\mathcal G],
 \qquad B=\E[|V|^q\mid\mathcal G].
\]
$A,B$ 非负且有限 a.s.。固定 $\varepsilon>0$，对点态 Young 不等式应用随机归一化
\[
 a=\frac{|U|}{(A+\varepsilon)^{1/p}},
 \qquad
 b=\frac{|V|}{(B+\varepsilon)^{1/q}},
\]
得到
\[
 \frac{|UV|}{(A+\varepsilon)^{1/p}(B+\varepsilon)^{1/q}}
 \leq
 \frac{|U|^p}{p(A+\varepsilon)}
 +\frac{|V|^q}{q(B+\varepsilon)}.
\]
所有分母倒数都是有界的 $\mathcal G$-可测变量，可以从条件期望中取出。于是
\[
 \frac{\E[|UV|\mid\mathcal G]}
 {(A+\varepsilon)^{1/p}(B+\varepsilon)^{1/q}}
 \leq\frac{A}{p(A+\varepsilon)}+\frac{B}{q(B+\varepsilon)}\leq1.
\]
依次取 $\varepsilon=1/n$，合并这可数个不等式的例外集，再令 $n\to\infty$。零分母处也由
该极限给出正确的零值，得到
\[
 \E[|UV|\mid\mathcal G]\leq A^{1/p}B^{1/q}.
\]
再用
$|\E[UV\mid\mathcal G]|\leq\E[|UV|\mid\mathcal G]$ 即得
\eqref{eq:5-2-conditional-holder}。整个论证没有在 $A=0$ 或 $B=0$ 处作除法。
\end{proof}

\begin{example}[在有限信息块上核对条件 \Holder{}]
令 $D$ 为公平骰子点数，$\mathcal G$ 只报告奇偶，取
$U=D$、$V=\1_{\{D\geq4\}}$ 以及 $p=q=2$。在奇数块与偶数块上分别有
\[
 \E[UV\mid\mathcal G]=\frac53,\ \frac{10}{3},\qquad
 \E[U^2\mid\mathcal G]=\frac{35}{3},\ \frac{56}{3},\qquad
 \E[V^2\mid\mathcal G]=\frac13,\ \frac23.
\]
所以条件 Cauchy--Schwarz 在两块上分别化为
\[
 \frac53\leq\frac{\sqrt{35}}3,\qquad
 \frac{10}{3}\leq\frac{\sqrt{112}}3;
\]
平方以后就是 $25\leq35$ 与 $100\leq112$。这正是普通 Cauchy--Schwarz 对每个条件分割块的
应用，而不是把两个块的平均混在一起。
\end{example}

端点版本若 $U\in L^1$、$V\in L^\infty$，则
\[
 |\E[UV\mid\mathcal G]|
 \leq\|V\|_\infty\E[|U|\mid\mathcal G].
\]
这是由 $|UV|\leq\|V\|_\infty|U|$ 和条件期望单调性得到的；估计只使用常数界
$\|V\|_\infty$。

\begin{theorem}[$L^2$ 投影性质]\label{theorem:5-2-3-2}
设 $X$ 为实值且 $X\in L^2$。对每个实值 $Y\in L^2(\mathcal G)$，
\begin{equation}\label{eq:5-2-orthogonality}
 \E\bigl[(X-\E[X\mid\mathcal G])Y\bigr]=0.
\end{equation}
因此 $P_{\mathcal G}X=\E[X\mid\mathcal G]$ 是 $L^2(\mathcal G)$ 中唯一最小化
$\E|X-Y|^2$ 的变量。
\end{theorem}

\begin{proof}
记 $P=P_{\mathcal G}X$。由 \eqref{eq:5-2-l2-contraction}，$P\in L^2(\mathcal G)$，故
$X-P\in L^2$。先设 $Y$ 有界。取出已知因子并使用塔式性质，得到
\[
 \begin{aligned}
 \E[(X-P)Y]
 &=\E\bigl[\E[(X-P)Y\mid\mathcal G]\bigr]\\
 &=\E\bigl[Y\E[X-P\mid\mathcal G]\bigr]=0.
 \end{aligned}
\]
对一般 $Y\in L^2(\mathcal G)$，令 $Y_m=(-m)\vee(Y\wedge m)$。\Holder{} 不等式给
\[
 \bigl|\E[(X-P)(Y-Y_m)]\bigr|
 \leq\|X-P\|_2\|Y-Y_m\|_2\longrightarrow0,
\]
所以 \eqref{eq:5-2-orthogonality} 仍成立。

任取 $Y\in L^2(\mathcal G)$，把
$X-Y=(X-P)+(P-Y)$ 展开平方；交叉项由
\eqref{eq:5-2-orthogonality}（取 $P-Y$）为零。因此
\begin{equation}\label{eq:5-2-pythagorean-prediction}
 \|X-Y\|_2^2=\|X-P\|_2^2+\|P-Y\|_2^2.
\end{equation}
右侧第二项非负，且等于零当且仅当 $Y=P$ 于 $L^2$，这同时给出最优性和唯一性。
\end{proof}

\begin{corollary}[总方差公式]\label{corollary:5-2-3-3}
若实值 $X\in L^2$，定义
\[
 \Var(X\mid\mathcal G)
 =\E\bigl[(X-\E[X\mid\mathcal G])^2\mid\mathcal G\bigr].
\]
则
\begin{equation}\label{eq:5-2-total-variance}
 \Var(X)=\E\Var(X\mid\mathcal G)
 +\Var\bigl(\E[X\mid\mathcal G]\bigr).
\end{equation}
\end{corollary}

\begin{proof}
令 $P=\E[X\mid\mathcal G]$。分解
\[
 X-\E X=(X-P)+(P-\E X)
\]
的两项属于 $L^2$，而第二项 $\mathcal G$-可测。由
\eqref{eq:5-2-orthogonality}，它们的内积为零。平方后取期望即得
\[
 \Var(X)=\E(X-P)^2+\E(P-\E X)^2.
\]
塔式性质把第一项写成
$\E\E[(X-P)^2\mid\mathcal G]=\E\Var(X\mid\mathcal G)$，第二项就是
$\Var(P)$。
\end{proof}

\begin{example}[总方差的骰子回算]\label{ex:5-2-3-2}
仍令 $X=D^2$，$\mathcal G$ 只报告奇偶。直接计算
\begin{align*}
 \E X&=\frac{1^2+2^2+3^2+4^2+5^2+6^2}{6}=\frac{91}{6},\\
 \E X^2&=\frac{1^4+2^4+3^4+4^4+5^4+6^4}{6}
 =\frac{2275}{6},\\
 \Var(X)&=\E X^2-(\E X)^2
 =\frac{2275}{6}-\left(\frac{91}{6}\right)^2
 =\frac{5369}{36}.
\end{align*}
奇数块和偶数块内的方差分别为
\[
 \frac13\left[\left(1-\frac{35}{3}\right)^2
 +\left(9-\frac{35}{3}\right)^2
 +\left(25-\frac{35}{3}\right)^2\right]=\frac{896}{9},
\]
\[
 \frac13\left[\left(4-\frac{56}{3}\right)^2
 +\left(16-\frac{56}{3}\right)^2
 +\left(36-\frac{56}{3}\right)^2\right]=\frac{1568}{9}.
\]
两块等概率，故
$\E\Var(X\mid\mathcal G)=1232/9$。条件均值取 $35/3,56/3$，各以概率 $1/2$ 出现，
其均值为 $91/6$，所以
\[
 \Var(\E[X\mid\mathcal G])=\frac12\left(-\frac72\right)^2
 +\frac12\left(\frac72\right)^2=\frac{49}{4}.
\]
最后
\[
 \frac{1232}{9}+\frac{49}{4}=\frac{5369}{36},
\]
与直接方差完全一致。
\end{example}

\begin{example}[联合高斯分布的条件均值]\label{ex:5-2-3-1}
本章称实随机向量 $(X,Y)$ \emph{联合高斯}，若每个实线性组合 $sX+tY$ 或具有一维
高斯分布，或退化为常数。设
\[
 \mu_X=\E X,\quad \mu_Y=\E Y,\quad
 \sigma_X^2=\Var(X),\quad\sigma_Y^2=\Var(Y)>0,
 \quad c=\Cov(X,Y).
\]
先考虑非退化密度模型。令 $\Delta=\sigma_X^2\sigma_Y^2-c^2>0$，并假设联合密度为
\[
 \frac1{2\pi\sqrt\Delta}
 \exp\!\left[-\frac{
 \sigma_Y^2(x-\mu_X)^2-2c(x-\mu_X)(y-\mu_Y)
 +\sigma_X^2(y-\mu_Y)^2}{2\Delta}\right].
\]
关于 $x$ 完成平方：分子中的二次式等于
\[
 \sigma_Y^2\left(x-\mu_X-\frac{c}{\sigma_Y^2}(y-\mu_Y)\right)^2
 +\frac{\Delta}{\sigma_Y^2}(y-\mu_Y)^2.
\]
记
\[
 m(y)=\mu_X+\frac{c}{\sigma_Y^2}(y-\mu_Y)
\]
。把完全平方式代回联合密度，再对 $x$ 作高斯积分，得到
\begin{align*}
 p_Y(y)
 &=\frac{e^{-(y-\mu_Y)^2/(2\sigma_Y^2)}}{2\pi\sqrt\Delta}
   \int_\R
   e^{-\sigma_Y^2(x-m(y))^2/(2\Delta)}\dd x\\
 &=\frac1{\sqrt{2\pi\sigma_Y^2}}
   e^{-(y-\mu_Y)^2/(2\sigma_Y^2)}.
\end{align*}
该密度对每个 $y$ 都严格为正，因而密度比给出条件密度
\[
 p_{X\mid Y=y}(x)
 =\sqrt{\frac{\sigma_Y^2}{2\pi\Delta}}
  \exp\!\left[-\frac{\sigma_Y^2(x-m(y))^2}{2\Delta}\right].
\]
它是均值 $m(y)$、方差 $\Delta/\sigma_Y^2$ 的高斯密度。确实，令
$u=x-m(y)$，则
\begin{align*}
 \int_\R x p_{X\mid Y=y}(x)\dd x
 &=m(y)\int_\R p_{X\mid Y=y}(x)\dd x\\
 &\quad+
 \sqrt{\frac{\sigma_Y^2}{2\pi\Delta}}
 \int_\R u e^{-\sigma_Y^2u^2/(2\Delta)}\dd u
 =m(y),
\end{align*}
因为最后一项是可积奇函数的积分。于是
\begin{equation}\label{eq:5-2-gaussian-regression}
 \E[X\mid Y]
 =\mu_X+\frac{\Cov(X,Y)}{\Var(Y)}(Y-\mu_Y).
\end{equation}
若 $\Delta=0$，则
\[
 \Var\!\left(X-\mu_X-\frac{c}{\sigma_Y^2}(Y-\mu_Y)\right)
 =\sigma_X^2-\frac{c^2}{\sigma_Y^2}=0,
\]
所以括号内变量 a.s. 为零，同一公式仍成立。若 $\Var(Y)=0$，则 $Y$ a.s. 为常数，条件均值为
$\E X$，而 \eqref{eq:5-2-gaussian-regression} 中的逆方差写法不适用。

第~\ref{subsec:5-4-3}~节将证明：协方差正定时，上述密度描述与“所有线性组合均为高斯分布”
的定义一致。因此该计算也给出一般非退化联合高斯变量对的条件均值。一般 $L^2$ 投影只说明
条件均值是所有 $Y$ 的可测函数中的最佳预测，并不能推出它必为仿射函数；仿射性来自高斯
结构。
\end{example}

现在回到换测度。分母可能为零，故贝叶斯公式必须在积分定义中验证，而不能把随机变量形式地约去。

\begin{theorem}[贝叶斯--Radon--Nikodym 公式]\label{theorem:5-2-3-4}
设 $\mathbb Q\ll\Prob$，$Z=\dd\mathbb Q/\dd\Prob$，且
$X\in L^1(\mathbb Q)$。则在
$Z_{\mathcal G}=\E_\Prob[Z\mid\mathcal G]>0$ 处，
\begin{equation}\label{eq:5-2-bayes-rn}
 \E_\mathbb Q[X\mid\mathcal G]
 =\frac{\E_\Prob[ZX\mid\mathcal G]}
 {\E_\Prob[Z\mid\mathcal G]}.
\end{equation}
在 $\{Z_{\mathcal G}=0\}$ 上可任取有限的 $\mathcal G$-可测值；这个集合是
$\mathbb Q$-零集。
\end{theorem}

\begin{proof}
记 $A=Z_{\mathcal G}$、$N=\E_\Prob[ZX\mid\mathcal G]$，并定义
\[
 W=\begin{cases}N/A,&A>0,\\0,&A=0.\end{cases}
\]
因为 $Z\geq0$ 且 $\E_\Prob Z=1$，$A$ 非负、可积并有限 a.s.。在
$E=\{A=0\}\in\mathcal G$ 上，
\[
 \int_EZ\,\dd\Prob=\int_EA\,\dd\Prob=0,
\]
故 $Z=0$ a.s. 于 $E$，并且 $\mathbb Q(E)=0$。基本模长估计给
\[
 |N|\leq\E_\Prob[Z|X|\mid\mathcal G].
\]
右端在 $E$ 上积分为 $\int_EZ|X|\,\dd\Prob=0$，所以 $N=0$ a.s. 于 $E$。

对任意非负 $\mathcal G$-可测函数 $H$，先对简单 $H$ 使用条件期望定义，再由单调收敛推广，得到
\begin{equation}\label{eq:5-2-density-on-information}
 \int HZ\,\dd\Prob=\int HA\,\dd\Prob,
\end{equation}
两边允许为 $+\infty$。因此
\[
 \begin{aligned}
 \int|W|\,\dd\mathbb Q
 &=\int|W|Z\,\dd\Prob
 =\int|W|A\,\dd\Prob\\
 &=\int_{\{A>0\}}|N|\,\dd\Prob
 \leq\int\E_\Prob[Z|X|\mid\mathcal G]\,\dd\Prob
 =\E_\Prob[Z|X|]<\infty.
 \end{aligned}
\]
所以 $W\in L^1(\mathbb Q)$。对 $G\in\mathcal G$，把
\eqref{eq:5-2-density-on-information} 分别用于 $W$ 的正负部或实虚部，得到
\[
 \begin{aligned}
 \int_GW\,\dd\mathbb Q
 &=\int_GWZ\,\dd\Prob
 =\int_GWA\,\dd\Prob\\
 &=\int_GN\,\dd\Prob
 =\int_GZX\,\dd\Prob
 =\int_GX\,\dd\mathbb Q.
 \end{aligned}
\]
这里 $WA=N$ a.s.，包括 $A=0$ 的集合。故 $W$ 满足关于 $\mathbb Q$ 的条件期望定义，
即为 \eqref{eq:5-2-bayes-rn} 的版本。
\end{proof}

\begin{example}[二项--Beta 共轭更新]\label{ex:5-2-3-3}
设 $\alpha,\beta>0$，参数 $\theta\in(0,1)$ 的先验分布为
$\mathrm{Beta}(\alpha,\beta)$，其密度为
\[
 \pi(\theta)=\frac{\theta^{\alpha-1}(1-\theta)^{\beta-1}}
 {B(\alpha,\beta)},\qquad0<\theta<1,
\]
其中
$B(a,b)=\int_0^1t^{a-1}(1-t)^{b-1}\,\dd t$。给定 $\theta$ 后作 $n$ 次独立 Bernoulli
试验，记成功次数为 $K$。似然为
\[
 \Prob(K=k\mid\theta)=\binom nk\theta^k(1-\theta)^{n-k}.
\]
对 $k\in\{0,\ldots,n\}$，参数与观测值的联合密度（对 $\theta$ 使用 Lebesgue 测度，
对 $k$ 使用计数测度）为
\[
 \binom nk\frac{\theta^{\alpha+k-1}(1-\theta)^{\beta+n-k-1}}
 {B(\alpha,\beta)}\1_{(0,1)}(\theta).
\]
因此 $K=k$ 的边缘概率是
\begin{align*}
 \Prob(K=k)
 &=\frac{\binom nk}{B(\alpha,\beta)}
   \int_0^1\theta^{\alpha+k-1}(1-\theta)^{\beta+n-k-1}\dd\theta\\
 &=\binom nk\frac{B(\alpha+k,\beta+n-k)}{B(\alpha,\beta)}.
\end{align*}
由 $\alpha,\beta>0$ 可知这个数严格为正。联合密度除以该边缘概率后，组合数和原归一化常数相消，得到后验密度
\[
 \begin{aligned}
 \pi(\theta\mid K=k)
 &=\frac{\theta^{\alpha+k-1}(1-\theta)^{\beta+n-k-1}}
 {B(\alpha+k,\beta+n-k)},\qquad0<\theta<1.
 \end{aligned}
\]
所以后验为 $\mathrm{Beta}(\alpha+k,\beta+n-k)$。“共轭”在这里表示先验密度乘以似然后，经一次
明确归一化仍落在同一函数族，而不是一条需要背诵的参数口诀。
\end{example}

\begin{figure}[htbp]
\centering
\resizebox{0.98\textwidth}{!}{%
\begin{tikzpicture}[node distance=11mm and 14mm]
  \node[book strong node,minimum width=24mm] (x) {$X\in L^2$};
  \node[book strong node,right=20mm of x,minimum width=31mm] (p) {$P_{\mathcal G}X$\\最佳可观察预测};
  \node[book node,below=of p,minimum width=31mm] (sub) {$L^2(\mathcal G)$};
  \draw[book accent arrow] (x)--node[above,book label,align=center,font=\scriptsize]{条件\\期望}(p);
  \draw[book arrow] (p)--(sub);
  \draw[book dashed arrow] (x.south)--node[below left,book label]{$X-P_{\mathcal G}X\perp L^2(\mathcal G)$}(sub.west);

  \node[book warm node,right=27mm of p,minimum width=25mm] (prior) {先验\\$\pi(\dd\theta)$};
  \node[book node,below=of prior,minimum width=28mm] (like) {似然\\$L(\theta,\dd y)$};
  \node[book strong node,right=24mm of prior,minimum width=29mm] (post) {后验\\$\pi(\dd\theta\mid y)$};
  \draw[book arrow] (prior)--(like);
  \draw[book accent arrow] (prior)--node[above,book label,align=center,font=\scriptsize]{乘似然\\并归一化}(post);
  \draw[book arrow] (like.east)--(post.south);
\end{tikzpicture}
}
\caption{左：条件期望的平方误差几何。右：贝叶斯更新是联合分布相对于观测边缘的测度分解，也可由 Radon--Nikodym 密度比计算。}
\label{fig:5-2-3-projection-bayes}
\end{figure}

条件 Jensen 需要凸性和相应可积性；若 $\varphi(X)$ 不可积，有限值条件期望公式不能直接套用。
$L^2$ 投影解释依赖平方可积，不能无条件移到一般 $L^1$。贝叶斯公式的分母为零时，联合分布不
决定该状态上的点值；定理只在 $\mathbb Q$-a.s. 意义下给出唯一版本。

平方误差分解还可用来比较不同信息层上的预测误差；这类比较所需的结论已由
投影恒等式~\eqref{eq:5-2-pythagorean-prediction} 给出。

\begin{exercise}[条件 Jensen 的等号]
设 $\varphi$ 严格凸，$X,\varphi(X)\in L^1$。证明若
$X=\E[X\mid\mathcal G]$ a.s.，则条件 Jensen 取等号；并在有限分割情形证明反向：若每块有正
概率且取等号，则 $X$ 在每块上 a.s. 为常数。
\end{exercise}

\begin{exercise}[最小均方预测]
令 $X\in L^2$，$Z$ 为标准 Borel 值随机元。证明在所有形如 $g(Z)\in L^2$ 的预测中，
$\E[X\mid Z]$ 唯一最小化均方误差。
\end{exercise}

\begin{exercise}[换测度的塔式检查]
在\cref{theorem:5-2-3-4}的假设下，证明
$\E_\Prob Z_{\mathcal G}=1$，并对 $X=1$ 直接验证贝叶斯--Radon--Nikodym 公式。再说明为什么
$Z_{\mathcal G}=0$ 的集合必为 $\mathbb Q$-零集，却未必为 $\Prob$-零集。
\end{exercise}

条件信息固定以后，随机变量序列仍可用 a.s.、依概率、$L^p$ 与依分布等不同方式趋近。下一节将
画出这些收敛概念之间的蕴含图，并用明确反例标出每条不能反向的箭头。

\section{概率收敛、弱收敛与紧性}\label{sec:05-03}

一个随机变量序列可以在几乎每条样本路径上安定下来，也可以只让偏离极限的概率趋于零；它还可以在平均误差中收敛，或者仅仅让每一个有界连续观测量的平均值收敛。这四种说法保留的信息不同。把它们都写成“$X_n\to X$”固然省墨，却会把证明中真正需要的结构一起省掉。

本节先把可靠的蕴含和失败的逆向放在同一张图里。随后我们处理两个升级问题：何时依概率收敛可以升级为矩收敛；何时一列分布至少有弱收敛子列。前一个问题的答案是统一控制大值尾部，后一个问题的答案是统一阻止质量逃向空间的远处。


\subsection{随机变量的四种收敛与反例图谱}\label{subsec:5-3-1}
\index{几乎必然收敛}\index{依概率收敛}\index{$L^p$ 收敛}\index{依分布收敛}

先说明依分布收敛的一项局限。比较 $X_n$ 与 $X$ 的分布可以控制
$\E f(X_n)$ 与 $\E f(X)$，却没有比较同一个 $\omega$ 上的 $X_n(\omega)$ 与
$X(\omega)$；分布已经忘掉样本点怎样彼此配对。因此，依分布收敛一般不能推出依概率收敛，
即使所有变量定义在同一个概率空间上。

\begin{definition}[四种收敛]\label{def:5-3-1-1}
设 $X_n,X$ 是同一概率空间上的实值随机变量，$1\le p<\infty$。
\begin{enumerate}
  \item 若存在 $N\in\mathcal F$，$\Prob(N)=0$，使得对每个 $\omega\notin N$ 都有 $X_n(\omega)\to X(\omega)$，则称 $X_n$ \emph{几乎必然收敛}于 $X$，记作 $X_n\to X$ a.s.；
  \item 若对每个 $\varepsilon>0$，
  \[
    \Prob\bigl(|X_n-X|>\varepsilon\bigr)\longrightarrow0,
  \]
  则称 $X_n$ \emph{依概率收敛}于 $X$，记作 $X_n\xrightarrow{\Prob}X$；
  \item 若 $X_n,X\in L^p$ 且
  \[
    \E|X_n-X|^p\longrightarrow0,
  \]
  则称 $X_n$ 在 $L^p$ 中收敛于 $X$；
  \item 若对每个 $f\in C_b(\R)$ 都有
  \[
    \int_{\R}f\,\dd\Law(X_n)\longrightarrow
    \int_{\R}f\,\dd\Law(X),
  \]
  则称 $X_n$ \emph{依分布收敛}于 $X$。
\end{enumerate}
前三项比较同一概率空间上的变量；第四项只比较推前测度。
\end{definition}

\begin{definition}[依分布符号]\label{def:5-3-1-2}
若概率测度 $\mu_n,\mu$ 对每个 $f\in C_b(\R)$ 满足
$\int f\,\dd\mu_n\to\int f\,\dd\mu$，记作 $\mu_n\Rightarrow\mu$。对随机变量写
$X_n\Rightarrow X$，其含义是 $\Law(X_n)\Rightarrow\Law(X)$。因此 $X_n$ 与 $X$ 可以定义在不同的概率空间上。
\end{definition}

\begin{definition}[随机阶记号]\label{def:5-3-1-3}
设数列 $a_n$ 对每个 $n$ 都满足 $a_n>0$。若对每个 $\eta>0$，存在 $M<\infty$ 与 $N$，使得
\[
  \sup_{n\ge N}\Prob\bigl(|X_n|>Ma_n\bigr)<\eta,
\]
则写 $X_n=O_\Prob(a_n)$。若 $X_n/a_n\xrightarrow{\Prob}0$，则写
$X_n=o_\Prob(a_n)$。
\end{definition}

此处的 $O_\Prob$ 直接由尾概率定义。等到\cref{subsec:5-3-3} 定义概率测度族的紧性后，
$X_n=O_\Prob(a_n)$ 就可以改写为归一化分布族 $\{\Law(X_n/a_n)\}$ 的紧性；目前不需要预借这层拓扑语言。

\begin{example}[打字机序列]\label{ex:5-3-1-1}
在 $([0,1),\Borel([0,1)),m)$ 上，对 $r\ge0$ 与 $0\le k<2^r$ 令
\[
 I_{r,k}=[k2^{-r},(k+1)2^{-r}),\qquad
 X_{2^r+k}=\1_{I_{r,k}}.
\]
若第 $n$ 项位于第 $r$ 层，则
\[
 \Prob(|X_n|>1/2)=2^{-r}\longrightarrow0,
 \qquad
 \E|X_n|^p=2^{-r}\longrightarrow0
\]
对每个 $1\le p<\infty$ 成立。因此 $X_n\xrightarrow{\Prob}0$，甚至在每个有限 $L^p$ 中收敛。
可是每个 $x\in[0,1)$ 在每一层恰好落入一个区间 $I_{r,k}$；沿整列看，$X_n(x)$ 在每层取一次 $1$，同时也取许多次 $0$。所以它在任何 $x\in[0,1)$ 都不收敛。这个例子同时否定“依概率 $\Rightarrow$ a.s.”与“$L^p\Rightarrow$ a.s.”。
\end{example}

\begin{example}[尖峰破坏矩收敛]\label{ex:5-3-1-2}
在 $([0,1],\Borel([0,1]),m)$ 上令
\[
  X_n=n\1_{(0,1/n)}.
\]
对每个 $x>0$，当 $n>x^{-1}$ 时 $x\notin(0,1/n)$，故 $X_n(x)\to0$；在 $x=0$ 处各项也为零。因此 $X_n\to0$ a.s.，从而也依概率收敛。另一方面
\[
  \E|X_n|=n\,m(0,1/n)=1.
\]
所以 $X_n$ 不在 $L^1$ 中收敛于零。把高度改为 $n^{1/p}$，同样得到 a.s. 收敛而 $L^p$ 误差恒为一的例子。
\end{example}

\begin{example}[同分布不同配对]\label{ex:5-3-1-3}
在一个承载两枚独立公平符号 $X,Y$ 的概率空间上，令 $X_n=Y$ 对每个 $n$ 成立。$X_n$ 与 $X$ 都服从 Rademacher 分布
$\tfrac12(\delta_{-1}+\delta_1)$，故 $X_n\Rightarrow X$。但是独立性给出
\[
 \Prob(|X_n-X|>1)=\Prob(Y\ne X)=\frac12,
\]
所以 $X_n$ 不依概率收敛于 $X$。这里序列甚至没有随 $n$ 变化；失败完全来自耦合，而不是渐近计算。
\end{example}

\begin{proposition}[基本蕴含]\label{proposition:5-3-1-1}
对 $1\le p<\infty$，有
\[
  X_n\longrightarrow X\text{ 于 }L^p
  \quad\Longrightarrow\quad
  X_n\xrightarrow{\Prob}X,
\]
以及
\[
  X_n\longrightarrow X\text{ a.s.}
  \quad\Longrightarrow\quad
  X_n\xrightarrow{\Prob}X
  \quad\Longrightarrow\quad
  X_n\Rightarrow X.
\]
上述箭头的逆向一般都不成立。
\end{proposition}

\begin{proof}
若 $X_n\to X$ 于 $L^p$，则对任意 $\varepsilon>0$，Markov 不等式
\cref{theorem:5-1-1-2} 给出
\[
 \Prob(|X_n-X|>\varepsilon)
 \le \varepsilon^{-p}\E|X_n-X|^p\longrightarrow0.
\]

再设 $X_n\to X$ a.s.，固定 $\varepsilon>0$，并令
\[
 A_N=\bigcup_{n\ge N}\{|X_n-X|>\varepsilon\}.
\]
集合 $A_N$ 随 $N$ 递减，而
$\bigcap_NA_N$ 是事件 $\{|X_n-X|>\varepsilon\ \text{i.o.}\}$。在 a.s. 收敛的共同满测集上，这一事件不发生，故 $\Prob(\bigcap_NA_N)=0$。概率测度从上连续，于是
\[
 \Prob(|X_n-X|>\varepsilon)
 \le \Prob(A_n)\longrightarrow0
 \qquad(n\to\infty).
\]

最后设 $X_n\xrightarrow{\Prob}X$。先证明任意子列 $(X_{n_j})$ 都有进一步子列 a.s. 收敛于 $X$。递归选取 $j_k$，使
\[
 \Prob\bigl(|X_{n_{j_k}}-X|>2^{-k}\bigr)<2^{-k}.
\]
Borel--Cantelli I（\cref{theorem:5-1-3-1}）表明这些事件只发生有限多次，所以
$X_{n_{j_k}}\to X$ a.s.。对任意
$f\in C_b(\R)$，连续性与 DCT 给
\[
 \E f(X_{n_{j_k}})\longrightarrow\E f(X).
\]
若原数列 $\E f(X_n)$ 不收敛于 $\E f(X)$，就存在 $\delta>0$ 及一个子列，使
$|\E f(X_{n_j})-\E f(X)|\ge\delta$ 对全部 $j$ 成立；刚构造的进一步子列与此矛盾。因此每个 $f\in C_b(\R)$ 的测试积分都收敛，即 $X_n\Rightarrow X$。

逆向失败分别由\cref{ex:5-3-1-1,ex:5-3-1-2,ex:5-3-1-3} 给出。
\end{proof}

\begin{theorem}[依概率子列原理]\label{theorem:5-3-1-2}
$X_n\xrightarrow{\Prob}X$，当且仅当 $(X_n)$ 的每个子列都有一个进一步子列 a.s. 收敛于 $X$。
\end{theorem}

\begin{proof}
正向的进一步子列已经在\cref{proposition:5-3-1-1} 的证明中构造：同时让坏事件的概率和阈值都按 $2^{-k}$ 下降，再用 Borel--Cantelli I。

反之，若 $X_n$ 不依概率收敛于 $X$，则存在 $\varepsilon,\delta>0$ 和子列 $(X_{n_k})$，使
\[
  \Prob(|X_{n_k}-X|>\varepsilon)\ge\delta
  \qquad(k\ge1).
\]
按假设，该子列有进一步子列 a.s. 收敛于 $X$。由\cref{proposition:5-3-1-1}，这个进一步子列又依概率收敛于 $X$，与上述统一下界矛盾。
\end{proof}

子列原理经常比定义更好用：要证明依概率收敛，可以对任意子列寻找 a.s. 收敛的进一步子列；要否定依概率收敛，则应找出一个坏子列，使任何进一步子列都无法摆脱同一概率下界。

\begin{theorem}[连续映射定理：处处连续版]\label{theorem:5-3-1-3}
若 $X_n\Rightarrow X$，且 $g:\R\to\R^m$ 连续，则
$g(X_n)\Rightarrow g(X)$。
\end{theorem}

\begin{proof}
任取 $f\in C_b(\R^m)$。复合函数 $f\circ g$ 属于 $C_b(\R)$，故
\[
 \E f(g(X_n))=\E (f\circ g)(X_n)\longrightarrow
 \E(f\circ g)(X)=\E f(g(X)).
\]
这正是依分布收敛的定义。
\end{proof}

若 $g$ 只在一个 $\Law(X)$-零集之外连续，上述一行证明不再成立，因为 $f\circ g$ 未必连续。
\cref{proposition:5-3-3-3} 将借助 Portmanteau 定理处理这种几乎处处连续的情形。

\begin{remark}[Slutsky 公式所需的联合控制]\label{proposition:5-3-1-4}
若 $X_n\Rightarrow X$、$Y_n\xrightarrow{\Prob}c$，且 $X_n,Y_n$ 定义在同一概率空间上，则
\[
 X_n+Y_n\Rightarrow X+c,
 \qquad
 X_nY_n\Rightarrow cX.
\]
这两个结论不能只由两个边缘分布分别收敛直接推出；关键是先证明
$(X_n,Y_n)\Rightarrow(X,c)$。\cref{proposition:5-3-3-slutsky} 将用联合分布的紧性
建立这一点，并由连续映射定理得到上述公式。
\end{remark}

一个不需要联合极限定理的特例可以直接验证。固定随机变量 $X$ 与有界随机变量 $Z$，令
$Y_n=c+n^{-1}Z$。则 $f(X+Y_n)\to f(X+c)$ a.s.，并由 $|f(X+Y_n)|\le\|f\|_\infty$ 与 DCT 得
$X+Y_n\Rightarrow X+c$。若把固定的 $X$ 换成只知道 $X_n\Rightarrow X$ 的序列，这个论证会在第一步停止：有界连续函数在整条实线上未必一致连续，$|Y_n-c|$ 小不能对所有可能的 $X_n$ 同时给出函数值误差小。真正缺少的是联合分布的紧性，而不是再写一次 DCT。

\begin{figure}[htbp]
\centering
\begin{tikzpicture}[node distance=18mm and 20mm]
  \node[book strong node] (lp) {$L^p$ 收敛};
  \node[book strong node,below=of lp] (as) {a.s. 收敛};
  \node[book strong node,right=of lp,yshift=-9mm] (pr) {依概率};
  \node[book strong node,right=of pr] (di) {依分布};
  \draw[book accent arrow] (lp)--(pr);
  \draw[book accent arrow] (as)--(pr);
  \draw[book accent arrow] (pr)--(di);
  \draw[book dashed arrow] (pr) to[bend left=18]
    node[above,font=\scriptsize,align=center]{逆向失败：\\打字机} (as);
  \draw[book dashed arrow] (di) to[bend left=18]
    node[below,font=\scriptsize,align=center]{逆向失败：\\独立副本} (pr);
  \draw[book dashed arrow] (as) to[bend right=20]
    node[left,font=\scriptsize,align=right]{逆向失败：\\尖峰} (lp);
\end{tikzpicture}
\caption{四种收敛中始终可靠的箭头。打字机序列还说明 $L^p$ 收敛不推出 a.s. 收敛；尖峰说明 a.s. 收敛不推出 $L^p$ 收敛。}
\label{fig:5-3-1-convergence-map}
\end{figure}

\begin{remark}[反例矩阵的读法]\label{rem:5-3-1-prob-to-law}
图\ref{fig:5-3-1-convergence-map} 没有把 a.s. 收敛与 $L^p$ 收敛排成一条链，因为二者互不蕴含。依分布收敛则不比较同一 $\omega$ 上的值。以后证明大数律、中心极限定理或估计量的渐近性质时，应先问结论需要保留哪一种信息，再决定使用哪条箭头；把更弱拓扑写成更强拓扑不会让定理更有力，只会让证明出现缺口。
\end{remark}

\begin{exercise}[四种收敛的比较]\label{exercise:5-3-1-1}
\begin{enumerate}
  \item 对打字机序列逐项验证每个有限 $L^p$ 收敛，并证明它在每一点都不收敛。
  \item 对固定 $X$ 与有界 $Z$，直接由 DCT 证明 $X+n^{-1}Z\Rightarrow X$。
  \item 构造一个 a.s. 收敛到零但在指定 $L^p$ 中不收敛的非负序列。
  \item 证明 $X_n=o_\Prob(a_n)$ 推出 $X_n=O_\Prob(a_n)$，并给出逆向失败的常数序列。
\end{enumerate}
\end{exercise}


\subsection{一致可积、\Scheffe{} 与期望收敛}\label{subsec:5-3-2}
\index{一致可积}\index{Scheffe@\Scheffe{} 引理}\index{全变差距离}

\cref{ex:5-3-1-2} 中，$X_n\to0$ a.s.，但全部期望都等于一。坏集的概率只有 $1/n$，期望却把函数高度也计算进去。若要从概率收敛得到平均收敛，就必须排除“越罕见越巨大”的补偿机制。

\begin{definition}[概率空间上的一致可积]\label{def:5-3-2-1}
一族可积随机变量 $\mathcal X$ 称为\emph{一致可积}（uniformly integrable，简称 UI），若
\[
 \lim_{K\to\infty}\sup_{X\in\mathcal X}
 \E\bigl[|X|\1_{\{|X|>K\}}\bigr]=0.
\]
\end{definition}

由\cref{lemma:3-4-1-ui-equivalence}，在概率空间上这等价于
\[
 \sup_{X\in\mathcal X}\E|X|<\infty
 \tag{5.3.1}\label{eq:5-3-ui-l1-bound}
\]
以及下列小事件条件：对每个 $\varepsilon>0$，存在 $\delta>0$，使
\[
 \Prob(A)<\delta
 \quad\Longrightarrow\quad
 \sup_{X\in\mathcal X}\E(|X|\1_A)<\varepsilon.
 \tag{5.3.2}\label{eq:5-3-ui-small-events}
\]
为了记住两个量词的作用，不妨重做其中的关键估计。由尾部定义，可取 $K_0$ 使
\[
 \sup_{X\in\mathcal X}\E\bigl(|X|\1_{\{|X|>K_0\}}\bigr)<1.
\]
因而对每个 $X\in\mathcal X$，
\[
 \E|X|
 \le K_0+\E\bigl(|X|\1_{\{|X|>K_0\}}\bigr)
 <K_0+1,
\]
这给出\eqref{eq:5-3-ui-l1-bound}。再给定 $\varepsilon>0$，取 $K$ 使统一尾期望小于
$\varepsilon/2$，并令 $\delta=\varepsilon/(2K)$。当 $\Prob(A)<\delta$ 时，
\[
 \E(|X|\1_A)
 \le K\Prob(A)+\E\bigl(|X|\1_{\{|X|>K\}}\bigr)
 <\varepsilon
\]
对全部 $X\in\mathcal X$ 同时成立，即\eqref{eq:5-3-ui-small-events}。

反过来，假设 $C=\sup_X\E|X|<\infty$ 且小事件条件成立。给定 $\varepsilon>0$，取对应的
$\delta>0$，再取 $K>C/\delta$。Markov 不等式给出
\[
 \Prob(|X|>K)\le\frac{\E|X|}{K}\le\frac CK<\delta.
\]
把 $A=\{|X|>K\}$ 代入小事件条件，得
$\sup_X\E(|X|\1_{\{|X|>K\}})<\varepsilon$，这就恢复了统一尾控制。

\begin{definition}[矩一致可积]\label{def:5-3-2-2}
研究 $p$ 阶矩时，称 $(X_n)$ 的 $p$ 阶矩一致可积，是指随机变量族
$\{|X_n|^p:n\ge1\}$ 一致可积。
\end{definition}

\begin{definition}[全变差距离]\label{def:5-3-2-3}
概率测度 $\mu,\nu$ 的\emph{全变差距离}定义为
\[
 d_{\mathrm{TV}}(\mu,\nu)=\sup_A|\mu(A)-\nu(A)|
 =\frac12|\mu-\nu|(S),
\]
其中上确界遍历可测集，$|\mu-\nu|$ 是有符号测度 $\mu-\nu$ 的全变差测度。若二者相对于同一测度 $\lambda$ 有密度 $p,q$，则
\[
 d_{\mathrm{TV}}(\mu,\nu)=\frac12\int|p-q|\,\dd\lambda.
\]
\end{definition}

最后一个等式中的因子 $1/2$ 不能丢掉。令 $r=p-q$；由于 $\int r=0$，有
$\int r^+=\int r^-=\frac12\int|r|$。对任意 $A$，
$|\int_A r|\le\int r^+$，而取 $A=\{r\ge0\}$ 达到等号。

\begin{example}[可计算的 UI 充分条件]\label{ex:5-3-2-1}
若 $p>1$ 且 $C:=\sup_n\E|X_n|^p<\infty$，则在 $\{|X_n|>K\}$ 上
$|X_n|\le K^{1-p}|X_n|^p$，因而
\[
 \sup_n\E\bigl[|X_n|\1_{\{|X_n|>K\}}\bigr]
 \le C K^{1-p}\longrightarrow0.
\]
更一般地，设 $\Phi:[0,\infty)\to[0,\infty)$ 凸、非降，且
$\Phi(t)/t\to\infty$。若 $C_\Phi:=\sup_n\E\Phi(|X_n|)<\infty$，则
当 $K$ 足够大（此时 $\Phi(t)>0$ 对所有 $t>K$ 成立）时，
\[
 \E\bigl[|X_n|\1_{\{|X_n|>K\}}\bigr]
 \le \left(\sup_{t>K}\frac{t}{\Phi(t)}\right)C_\Phi\longrightarrow0.
\]
这是 de la \ValleePoussin{} 判据的充分方向；这里不讨论其必要性。
\end{example}

\begin{example}[一致 $L^1$ 有界并不够]\label{ex:5-3-2-2}
仍取 $X_n=n\1_{(0,1/n)}$。虽然 $\E|X_n|=1$ 对所有 $n$ 成立，但只要 $n>K$，
\[
 \E\bigl[|X_n|\1_{\{|X_n|>K\}}\bigr]=1.
\]
所以该族不 UI。它依概率收敛于零而期望不收敛，失败恰好被 UI 定义检测到。
\end{example}

\begin{theorem}[Vitali 概率版]\label{theorem:5-3-2-1}
若 $X_n\xrightarrow{\Prob}X$，且 $(X_n)$ UI，则 $X\in L^1$，并且
\[
 \E|X_n-X|\longrightarrow0.
\]
特别地，$\E X_n\to\E X$。
\end{theorem}

\begin{proof}
概率空间是有限测度空间，依概率收敛就是\cref{def:3-4-1-1} 的依测度收敛，而
\cref{def:5-3-2-1} 与\cref{def:3-4-1-2} 是同一定义。因此\cref{theorem:3-4-1-2} 直接给出 $X\in L^1$ 与
$\|X_n-X\|_1\to0$。再由积分三角不等式，
\[
 |\E X_n-\E X|\le\E|X_n-X|\longrightarrow0.
\]
\end{proof}

\begin{proposition}[$L^p$ 升级]\label{proposition:5-3-2-2}
设 $1\le p<\infty$。若 $X_n\xrightarrow{\Prob}X$，且
$\{|X_n|^p:n\ge1\}$ UI，则 $X\in L^p$ 且 $X_n\to X$ 于 $L^p$。
\end{proposition}

\begin{proof}
UI 蕴含 $C:=\sup_n\E|X_n|^p<\infty$。由依概率子列原理，从任意子列都能抽出进一步子列
$X_{n_k}\to X$ a.s.；Fatou 引理给
\[
 \E|X|^p\le\liminf_k\E|X_{n_k}|^p\le C.
\]
所以 $X\in L^p$。

令 $Z_n=|X_n-X|^p$。连续映射 $t\mapsto|t|^p$ 给出 $Z_n\xrightarrow{\Prob}0$。
上一步已经证明 $X\in L^p$，现在可以利用点态不等式
\[
 Z_n\le2^{p-1}(|X_n|^p+|X|^p)
\]
先给一致 $L^1$ 界。给定 $\varepsilon>0$，$\{|X_n|^p\}$ 的 UI 小事件条件与固定可积函数 $|X|^p$ 的积分绝对连续性给出 $\delta>0$，使 $\Prob(A)<\delta$ 时
\[
 \sup_n\E(Z_n\1_A)
 \le2^{p-1}\left(
   \sup_n\E(|X_n|^p\1_A)+\E(|X|^p\1_A)
 \right)<\varepsilon.
\]
由\eqref{eq:5-3-ui-l1-bound}--\eqref{eq:5-3-ui-small-events} 的等价形式，$(Z_n)$ UI。Vitali 概率版作用于 $Z_n\to0$，得到
\[
 \E|X_n-X|^p=\E Z_n\longrightarrow0.
\]
\end{proof}

\begin{theorem}[\Scheffe{} 引理]\label{theorem:5-3-2-3}
设 $f_n,f\ge0$ 可测，$f_n\to f$ a.e.，并且
\[
 \int f_n\longrightarrow\int f<\infty.
\]
则
\[
 \int|f_n-f|\longrightarrow0.
\]
\end{theorem}

\begin{proof}
有恒等式
\[
 |f_n-f|=f_n+f-2(f_n\wedge f).
\]
因为 $f_n\wedge f\to f$ a.e.，且 $0\le f_n\wedge f\le f\in L^1$，DCT 给
$\int(f_n\wedge f)\to\int f$。代入恒等式并使用 $\int f_n\to\int f$，得到
\[
 \int|f_n-f|
 =\int f_n+\int f-2\int(f_n\wedge f)\longrightarrow0.
\]
\end{proof}

\begin{example}[密度与全变差的直接计算]\label{ex:5-3-2-3}
在 $[0,1]$ 上令
\[
 p_n(x)=1+\frac1n\sin(2\pi x),\qquad p(x)=1.
\]
这些函数非负，积分都为一，且 $p_n\to p$ 处处。\Scheffe{} 引理给相应概率分布的全变差收敛；这里还可以把距离算完：
\[
 \int_0^1|p_n-p|\,\dd x
 =\frac1n\int_0^1|\sin(2\pi x)|\,\dd x
 =\frac{2}{\pi n},
 \qquad
 d_{\mathrm{TV}}(p_n\dd x,p\dd x)=\frac1{\pi n}.
\]
\end{example}

\Scheffe{} 的“非负且总质量收敛”确实承担工作。若
$g_n=n\1_{(0,1/n)}-n\1_{(1/n,2/n)}$，则 $g_n\to0$ a.e. 且 $\int g_n=0$，但
$\int|g_n|=2$；去掉非负性，结论失败。若
$f_n=1+n\1_{(0,1/n)}$、$f=1$，则 $f_n\to f$ a.e. 而
$\int f_n=2\not\to1$，且 $\int|f_n-f|=1$；去掉总质量收敛，移动尖峰仍能携带固定质量。

\begin{proposition}[$L^1$ 收敛的 UI 必要性]\label{proposition:5-3-2-4}
若 $X_n\to X$ 于 $L^1$，则 $(X_n)$ UI，并且 $X_n\xrightarrow{\Prob}X$。因此，对一个已知依概率收敛于 $X$ 的序列，
\[
 X_n\to X\text{ 于 }L^1
 \quad\Longleftrightarrow\quad
 (X_n)\text{ UI}.
\]
\end{proposition}

\begin{proof}
$L^1$ 收敛蕴含依概率收敛已由 Markov 不等式证明。又有
$\sup_n\|X_n\|_1<\infty$。给定 $\varepsilon>0$，先取 $N$ 使 $n\ge N$ 时
$\|X_n-X\|_1<\varepsilon/2$。可积函数积分的绝对连续性给出 $\delta>0$，使 $\Prob(A)<\delta$ 时
\[
 \E(|X|\1_A)<\frac\varepsilon2,
 \qquad
 \max_{1\le n<N}\E(|X_n|\1_A)<\varepsilon;
\]
第二个条件只需对有限多个可积函数取最小的 $\delta$。于是对 $n\ge N$，
\[
 \E(|X_n|\1_A)
 \le\|X_n-X\|_1+\E(|X|\1_A)<\varepsilon.
\]
故整族满足 UI 的小事件条件，再结合一致 $L^1$ 界即得 UI。最后一个等价的正向由 Vitali 概率版给出，反向就是刚证明的必要性。
\end{proof}

依分布收敛本身仍不足以比较期望。令 $X_n$ 以概率 $1/n$ 取值 $n$，其余时候取零，则
$X_n\xrightarrow{\Prob}0$，因而 $X_n\Rightarrow0$，但 $\E X_n=1$。这不是期望“不连续”的神秘现象，而是测试函数 $x\mapsto x$ 不属于 $C_b(\R)$，同时该族不 UI。

同一截断--尾部分解也适用于鞅的 $L^1$ 极限：依概率或 a.s. 收敛控制截断后的有界部分，UI
统一控制条件期望族的尾部。这里所需的抽象工具正是
\cref{theorem:5-3-2-1,proposition:5-3-2-2,proposition:5-3-2-4}。

\begin{figure}[htbp]
\centering
\begin{tikzpicture}[x=1cm,y=1cm]
  \draw[book line] (0,0)--(11,0);
  \fill[BookAccent!8] (2,0) rectangle (9,2.0);
  \fill[BookWarning!10] (0,0) rectangle (2,2.0);
  \fill[BookWarning!10] (9,0) rectangle (11,2.0);
  \draw[book line] (2,0)--(2,2.2);
  \draw[book line] (9,0)--(9,2.2);
  \node[book label,align=center] at (5.5,1.25) {截断部分：依概率收敛 $+$ 一致有界\\
    $\Longrightarrow$ 平均误差小};
  \node[book label,align=center] at (1,1.15) {左尾\\UI 剪除};
  \node[book label,align=center] at (10,1.15) {右尾\\UI 剪除};
  \draw[decorate,decoration={brace,amplitude=5pt},BookAccent] (2,-0.25)--(9,-0.25)
    node[midway,below=6pt,font=\small]{截断区};
\end{tikzpicture}
\caption{Vitali 定理中的截断--尾部分工。图只记录误差分解；截断部分的估计与尾积分估计仍分别需要证明。}
\label{fig:5-3-2-vitali-body-tail}
\end{figure}

\begin{exercise}[UI 与质量守恒]\label{exercise:5-3-2-1}
\begin{enumerate}
  \item 直接用尾积分定义证明：若 $\sup_n\E|X_n|^p<\infty$ 且 $p>1$，则 $(X_n)$ UI。
  \item 证明 \Scheffe{} 引理，并从有密度的全变差公式推出事件概率的一致收敛。
  \item 对 $X_n=n\1_{(0,1/n)}$ 验证 $X_n\Rightarrow0$，同时定位 UI 定义中失败的量词。
  \item 设 $X_n\xrightarrow{\Prob}X$。证明 $X_n\to X$ 于 $L^p$ 当且仅当 $\{|X_n-X|^p\}$ UI。
\end{enumerate}
\end{exercise}
\subsection{Portmanteau、紧性与 Prokhorov 定理}\label{subsec:5-3-3}
\index{弱收敛}\index{紧性}\index{Portmanteau 定理}\index{Prokhorov 定理}

依分布收敛把概率测度放进一个由 $C_b$ 测试函数产生的拓扑。定义很短，随后却有三个必须分别解决的问题：
\begin{enumerate}
  \item 函数积分的收敛怎样控制开集、闭集的概率；
  \item 什么条件阻止一族概率把质量送到越来越远的地方；
  \item 对角抽取所得的测试积分极限为什么来自可数可加概率，而不是一个仅仅有限可加的“伪积分”。
\end{enumerate}
Portmanteau 回答第一问，紧性刻画第二问；第三问需要把极限泛函表示成可数可加概率测度。

\begin{definition}[概率族的紧性]\label{def:5-3-3-3}
设 $S$ 为 Polish 空间，$\mathcal K$ 是由 $S$ 上 Borel 概率测度组成的族。若对每个
$\varepsilon>0$，存在同一个紧集 $K_\varepsilon\subset S$，使
\[
 \sup_{\nu\in\mathcal K}\nu(K_\varepsilon^c)<\varepsilon,
\]
则称 $\mathcal K$ \emph{具有紧性}，也称它为\emph{紧概率族}。
\end{definition}

这里“同一个”的量词不能省略：紧集可以依赖 $\varepsilon$，但不能依赖
$\nu\in\mathcal K$。

\begin{proposition}[Polish 概率的紧逼近与 Radon 正则性]\label{proposition:5-3-3-polish-radon}
设 $S$ 为 Polish 空间。每个 Borel 概率 $\mu$ 组成的单元素族都具有紧性：对每个 $\varepsilon>0$，存在紧集
$K\subset S$ 使 $\mu(K)>1-\varepsilon$。此外，$\mu$ 是 Radon 概率，即每个 Borel 集都可由紧集从内逼近、由开集从外逼近。
\end{proposition}

\begin{proof}
固定一个有界的兼容完备度量 $d$ 以及稠密列 $(x_j)_{j\ge1}$。这样的度量总能取得：若
$d_0$ 兼容且完备，则 $d=d_0/(1+d_0)$ 诱导同一拓扑。函数
$r\mapsto r/(1+r)$ 非降，且对 $a,b\geq0$ 有
\[
 \frac{a}{1+a}+\frac{b}{1+b}-\frac{a+b}{1+a+b}
 =\frac{ab(2+a+b)}{(1+a)(1+b)(1+a+b)}\geq0,
\]
故 $d$ 仍满足三角不等式，确为度量。更具体地，对 $0<\delta<1$，
\[
 d(x,y)<\delta
 \quad\Longrightarrow\quad
 d_0(x,y)=\frac{d(x,y)}{1-d(x,y)}<\frac{\delta}{1-\delta}.
\]
因此 $d$-Cauchy 列也是 $d_0$-Cauchy 列，从而有 $d_0$-极限；两个度量拓扑相同，所以该极限也是
$d$-极限，即 $d$ 完备。给定 $\varepsilon>0$。对每个
$m\ge1$，开球族
\[
 \{B(x_j,2^{-m-2}):j\ge1\}
\]
覆盖 $S$。这些球的有限前段之并递增到 $S$，故测度从下连续性给出 $N_m$，使
\[
 \mu\left(\bigcup_{j=1}^{N_m}B(x_j,2^{-m-2})\right)
 >1-\varepsilon2^{-m-1}.
\]
令
\[
 F_m=\bigcup_{j=1}^{N_m}\overline B(x_j,2^{-m-1}),
 \qquad K=\bigcap_{m\ge1}F_m.
\]
每个 $F_m$ 闭，因此 $K$ 闭；并且
\[
 \mu(K^c)\le\sum_{m\ge1}\mu(F_m^c)
 <\varepsilon\sum_{m\ge1}2^{-m-1}<\varepsilon.
\]
对每个 $m$，有限集 $\{x_1,\ldots,x_{N_m}\}$ 是 $K$ 的 $2^{-m-1}$-网，所以 $K$ 全有界。
$K$ 是完备空间 $S$ 的闭子集，因而完备；完备且全有界的度量空间紧。紧性得证。

下面证明正则性。若 $F$ 闭，先取上面构造的紧集 $K$ 使 $\mu(K^c)<\varepsilon$，则
$F\cap K$ 是 $F$ 中的紧集，且
\[
 \mu(F\setminus(F\cap K))\le\mu(K^c)<\varepsilon.
\]
又令 $F^{1/n}=\{x:d(x,F)<1/n\}$。当 $F\ne\varnothing$ 时，
$F^{1/n}\downarrow F$，故测度从上连续性给 $\mu(F^{1/n}\setminus F)\to0$；空集情形直接取空开集。
所以闭集同时有紧内逼近与开外逼近。

记 $\mathcal R$ 为同时具有这两种逼近的 Borel 集全体。它包含闭集。若 $A\in\mathcal R$，取开集
$G\supset A$ 使 $\mu(G\setminus A)<\varepsilon/2$，再取全局紧集 $K$ 使
$\mu(K^c)<\varepsilon/2$，则 $K\setminus G$ 是 $A^c$ 中的紧集，并且
\[
 \mu(A^c\setminus(K\setminus G))
 \le\mu(K^c)+\mu(G\setminus A)<\varepsilon.
\]
另一方面，若 $C\subset A$ 紧且 $\mu(A\setminus C)<\varepsilon$，则开集
$C^c\supset A^c$ 满足 $\mu(C^c\setminus A^c)<\varepsilon$。故 $\mathcal R$ 对补集封闭。

最后设 $A_n\in\mathcal R$。分别取开集 $G_n\supset A_n$，使
$\mu(G_n\setminus A_n)<\varepsilon2^{-n-1}$，则
$G=\bigcup_nG_n$ 是 $\bigcup_nA_n$ 的开外逼近，误差小于 $\varepsilon$。
又因 $\bigcup_{n\le N}A_n\uparrow\bigcup_nA_n$，可取 $N$ 使尾部造成的测度差小于
$\varepsilon/2$；再对 $1\le n\le N$ 取紧集 $C_n\subset A_n$，使
$\mu(A_n\setminus C_n)<\varepsilon/(2N)$。有限并 $\bigcup_{n\le N}C_n$ 是所需紧内逼近。
所以 $\mathcal R$ 是包含闭集的 $\sigma$-代数，因而包含全部 Borel 集。
\end{proof}

\begin{example}[无限维中的紧质量盒]\label{ex:5-3-3-hilbert-tight}
设 $(e_n)$ 是 $\ell^2$ 的标准正交基，$\xi_n$ 是同一概率空间上的实值随机变量，且
$|\xi_n|\le1$ a.s.。令
\[
 X=\sum_{n\ge1}2^{-n}\xi_ne_n.
\]
除去一个共同零集后，级数逐样本绝对收敛于 $\ell^2$；其部分和是可测映射，故极限
$X$ 也是 $\ell^2$-值随机元。它的分布集中在
\[
 K=\{x\in\ell^2:|\langle x,e_n\rangle|\le2^{-n}\text{ 对每个 }n\}
\]
中。$K$ 闭，且其尾部满足
\[
 \sup_{x\in K}\sum_{n>N}|\langle x,e_n\rangle|^2
 \le\sum_{n>N}4^{-n}\longrightarrow0.
\]
给定 $\varepsilon>0$，先取 $N$ 使
\[
 \left(\sum_{n>N}4^{-n}\right)^{1/2}<\frac{\varepsilon}{2}.
\]
有限维盒
\[
 Q_N=\left\{(x_1,\ldots,x_N):|x_j|\le2^{-j},\ 1\le j\le N\right\}
 \subset\R^N
\]
紧，故有有限 $\varepsilon/2$-网 $\{q^1,\ldots,q^M\}$。把每个 $q^r$ 在第 $N$ 个坐标后补零，得到
$\widetilde q^{\,r}\in K$。对任意 $x\in K$，选 $r$ 使前 $N$ 个坐标与 $q^r$ 的欧氏距离小于
$\varepsilon/2$，则
\[
 \|x-\widetilde q^{\,r}\|_{\ell^2}^2
 <\frac{\varepsilon^2}{4}+\sum_{n>N}4^{-n}
 <\frac{\varepsilon^2}{2}<\varepsilon^2.
\]
因此 $K$ 对每个 $\varepsilon$ 都有有限 $\varepsilon$-网，即 $K$ 全有界。又 $K$ 是完备空间 $\ell^2$ 的闭子集，因而完备；完备且全有界故紧。这个例子提醒我们：无限维闭球通常不紧，但一个概率仍可把全部质量放在更薄的紧集上。
\end{example}

\begin{definition}[Polish 空间上的弱收敛]\label{def:5-3-3-1}
设 $S$ 为 Polish 空间，记 $\mathsf P(S)$ 为 $S$ 上全体 Borel 概率测度。使每个映射
\[
 \nu\longmapsto\int_S f\,\dd\nu,\qquad f\in C_b(S),
\]
连续的最弱拓扑称为 $\mathsf P(S)$ 上的\emph{弱拓扑}。对
$\mu_n,\mu\in\mathsf P(S)$，若对每个 $f\in C_b(S)$，
\[
 \int_S f\,\dd\mu_n\longrightarrow\int_S f\,\dd\mu,
\]
则称 $\mu_n$ \emph{弱收敛}于 $\mu$，仍记作 $\mu_n\Rightarrow\mu$；这恰是上述弱拓扑中的
序列收敛。由
\cref{proposition:5-3-3-polish-radon}，这里的概率都是 Radon 概率。
测试实值函数已经足够；若采用复值 $C_b(S)$，分别考察实部与虚部得到同一定义。
\end{definition}

在局部紧空间上只用 $C_c$ 测试得到的是\emph{模糊收敛}（vague convergence）；它不自动保留总质量。例如
$\delta_n$ 在 $\R$ 上对每个 $f\in C_c(\R)$ 都满足 $\int f\,\dd\delta_n=f(n)\to0$，但零测度不是概率。这正是弱概率拓扑坚持使用 $C_b$ 的原因。

\begin{definition}[相对弱紧]\label{def:5-3-3-2}
设 $S$ 为 Polish 空间。若 $\mathcal K\subset\mathsf P(S)$ 的弱闭包在弱拓扑下紧，则称 $\mathcal K$
\emph{相对弱紧}。在证明弱拓扑可度量后，这将等价于：$\mathcal K$ 中每个序列都有弱收敛子列。
\end{definition}

由\cref{proposition:5-3-3-polish-radon}，每个单独的 Polish 概率都有自己的紧逼近；这并不说明任意概率族可以共享一个紧集。

\begin{definition}[$\mu$-连续集]\label{def:5-3-3-4}
设 $S$ 为 Polish 空间、$\mu\in\mathsf P(S)$。若 Borel 集 $A\subset S$ 满足
$\mu(\partial A)=0$，则称 $A$ 为
\emph{$\mu$-连续集}。
\end{definition}

\begin{remark}[Prokhorov 定理的证明结构]\label{rem:5-3-3-four-stations}
证明包含四个相互独立的环节。Portmanteau 把 $C_b$ 测试翻译成开闭集不等式；弱紧性给出共同
紧逼近；Daniell 单调连续性排除有限可加的伪极限；最后由弱拓扑的度量化和对角抽取得到
Prokhorov 定理。特别地，子列抽取之后仍须证明极限泛函对应可数可加概率。
\end{remark}

\begin{example}[移动 Dirac 质量]\label{ex:5-3-3-1}
对度量空间 $S$，
\[
 \delta_{x_n}\Rightarrow\delta_x
 \quad\Longleftrightarrow\quad x_n\to x.
\]
正向的逆否命题如下：若某子列满足 $d(x_{n_k},x)\ge\varepsilon$，则
$f(y)=\min\{1,d(y,x)/\varepsilon\}\in C_b(S)$，并有
$\int f\,\dd\delta_{x_{n_k}}=1$、$\int f\,\dd\delta_x=0$。反向则由
$f(x_n)\to f(x)$ 对每个 $f\in C_b(S)$ 得到。

在 $\R$ 上，$\delta_n$ 不是紧族，也没有弱收敛子列。事实上，对任意严格递增子列
$(n_k)$，令 $a_k=0$（$k$ 偶）及 $a_k=1$（$k$ 奇），并在每个区间
$[n_k,n_{k+1}]$ 上定义
\[
 f(x)=a_k+\frac{a_{k+1}-a_k}{n_{k+1}-n_k}(x-n_k),
\]
在 $(-\infty,n_1]$ 上令 $f=a_1$。这给出 $0\le f\le1$ 的连续函数，并且
$\int f\,\dd\delta_{n_k}=f(n_k)=a_k$ 交替取零与一，所以该子列不可能弱收敛。
\end{example}

\begin{example}[模糊收敛丢失质量]\label{ex:5-3-3-3}
对 $f\in C_c(\R)$，有 $f(n)=0$ 当 $n$ 足够大，所以 $\delta_n$ 模糊收敛于零测度。
另一方面，常数函数 $1\in C_b(\R)$ 给
$\int1\,\dd\delta_n=1\not\to0$。这项单行计算精确区分了模糊收敛与概率弱收敛。
\end{example}

\begin{theorem}[Portmanteau 定理]\label{theorem:5-3-3-1}
设 $S$ 为 Polish 空间，$\mu_n,\mu\in\mathsf P(S)$。以下条件等价：
\begin{enumerate}
  \item $\mu_n\Rightarrow\mu$；
  \item 对每个开集 $G$，
  \[
    \mu(G)\le\liminf_n\mu_n(G);
  \]
  \item 对每个闭集 $F$，
  \[
    \limsup_n\mu_n(F)\le\mu(F);
  \]
  \item 对每个 $\mu$-连续集 $A$，有 $\mu_n(A)\to\mu(A)$；
  \item 对每个有界下半连续函数 $h:S\to\R$，
  \[
    \int h\,\dd\mu\le\liminf_n\int h\,\dd\mu_n.
  \]
\end{enumerate}
\end{theorem}

\begin{proof}
先证 (1)$\Rightarrow$(2)。若 $G=S$，结论由概率总质量为一得到。若 $G\ne S$，令
\[
 \phi_k(x)=\min\{1,k\,d(x,G^c)\}.
\]
则 $\phi_k\in C_b(S)$、$0\le\phi_k\le\1_G$，且 $\phi_k\uparrow\1_G$。因此对每个固定 $k$，
\[
 \int\phi_k\,\dd\mu
 =\lim_n\int\phi_k\,\dd\mu_n
 \le\liminf_n\mu_n(G).
\]
令 $k\to\infty$ 并用 MCT，得到 (2)。

(2) 与 (3) 等价，因为 $F^c$ 开且所有测度总质量为一。若 (2)、(3) 成立且
$\mu(\partial A)=0$，则
\[
 \mu(A^\circ)\le\liminf_n\mu_n(A)
 \le\limsup_n\mu_n(A)\le\mu(\overline A).
\]
而 $\overline A\setminus A^\circ\subset\partial A$，所以两端都等于 $\mu(A)$；这证明 (4)。

以下由 (4) 推回 (1)。固定 $f\in C_b(S)$。经仿射变换可设 $0\le f\le1$。
实数 $t$ 满足 $\mu(f^{-1}\{t\})>0$ 的至多可数，因为这些互不相交集合的正质量总和不超过一。
给定 $\eta>0$，选取
\[
 t_0<t_1<\cdots<t_m
\]
使 $t_0<0<t_1$、$t_{m-1}<1<t_m$、相邻距离小于 $\eta$，且每个
$\mu(f^{-1}\{t_j\})=0$。令
$A_j=\{t_{j-1}<f\le t_j\}$。其边界包含于
$f^{-1}\{t_{j-1},t_j\}$，故 $A_j$ 都是 $\mu$-连续集，并且它们分割 $S$。于是
\[
 \sum_j t_{j-1}\mu_n(A_j)
 \le\int f\,\dd\mu_n
 \le\sum_j t_j\mu_n(A_j).
\]
对 $n\to\infty$ 使用 (4)，上下和分别收敛到相应的 $\mu$-和，而两和之差小于
$\eta\sum_j\mu(A_j)=\eta$。令 $\eta\downarrow0$，得到
$\int f\,\dd\mu_n\to\int f\,\dd\mu$，即 (1)。

还需连接 (2) 与 (5)。设 (2) 成立。给有界下半连续 $h$ 加一个常数并缩放，可先设
$0\le h\le M$。下半连续性说明 $\{h>t\}$ 对每个 $t$ 都开。层蛋糕公式与 Fatou 引理给
\[
\begin{aligned}
 \int h\,\dd\mu
 &=\int_0^M\mu(h>t)\,\dd t\\
 &\le\int_0^M\liminf_n\mu_n(h>t)\,\dd t\\
 &\le\liminf_n\int_0^M\mu_n(h>t)\,\dd t
 =\liminf_n\int h\,\dd\mu_n.
\end{aligned}
\]
一般有界 $h$ 通过加常数还原。因此 (2)$\Rightarrow$(5)。反之，开集的指标函数
$\1_G$ 是有界下半连续函数；把它代入 (5) 便得到 (2)。
\end{proof}

\begin{proposition}[弱紧族的一致紧逼近]\label{proposition:5-3-3-compact-tight}
设 $S$ 为 Polish 空间，$\mathcal K\subset\mathsf P(S)$ 在弱拓扑下紧。则对每个
$\varepsilon>0$，存在紧集 $K\subset S$，使
\[
 \inf_{\nu\in\mathcal K}\nu(K)>1-\varepsilon.
\]
\end{proposition}

\begin{proof}
取有界兼容完备度量 $d$。固定 $m\ge1$。对每个 $\mu\in\mathcal K$，
\cref{proposition:5-3-3-polish-radon} 先给出紧集 $C_\mu$，使
$\mu(C_\mu)>1-\varepsilon2^{-m-1}$。由 $C_\mu$ 的紧性，可从以其各点为中心、半径
$2^{-m-1}$ 的开球覆盖中取有限子覆盖；把这些球之并记为 $G_\mu$。于是
\[
 \mu(G_\mu)>1-\varepsilon2^{-m-1}.
\]
对开集 $G$，映射 $\nu\mapsto\nu(G)$ 在弱拓扑下下半连续。的确，当 $G\ne S$ 时
\[
 \nu(G)=\sup_{k\ge1}\int\min\{1,kd(x,G^c)\}\,\dd\nu,
\]
右侧是连续函数族的逐点上确界；$G=S$ 时该映射恒为一。因此
\[
 U_\mu=\{\nu\in\mathcal K:\nu(G_\mu)>1-\varepsilon2^{-m-1}\}
\]
是 $\mu$ 在 $\mathcal K$ 中的开邻域。紧性给出有限子覆盖
$U_{\mu_1},\ldots,U_{\mu_r}$。

把 $G_{\mu_1},\ldots,G_{\mu_r}$ 中出现的全部球心合并为有限集 $E_m$，令
\[
 F_m=\bigcup_{x\in E_m}\overline B(x,2^{-m}).
\]
每个 $\nu\in\mathcal K$ 属于某个 $U_{\mu_i}$，且 $G_{\mu_i}\subset F_m$，故
\[
 \nu(F_m)>1-\varepsilon2^{-m-1}.
\]
令 $K=\bigcap_{m\ge1}F_m$。它闭，并在每个尺度有有限网；因此它在完备空间 $S$ 中紧。最后对每个
$\nu\in\mathcal K$，
\[
 \nu(K^c)\le\sum_{m\ge1}\nu(F_m^c)
 <\varepsilon\sum_{m\ge1}2^{-m-1}<\varepsilon.
\]
\end{proof}

\begin{remark}[弱紧性与一致紧逼近]\label{rem:5-3-3-compact-tight-no-loop}
上一命题的紧性假设位于概率测度空间 $\mathsf P(S)$，结论则是在原空间 $S$ 中找到共同紧集。
有限子覆盖把每个概率各自的紧逼近统一成对整个弱紧族有效的逼近；这个论证只使用弱拓扑的定义
和单个 Polish 概率的紧逼近。
\end{remark}

\begin{proposition}[弱拓扑的有界 Lipschitz 度量]\label{proposition:5-3-3-weak-metric}
设 $S$ 为 Polish 空间，并固定 $S$ 上一个有界兼容完备度量 $d$。定义
\[
 \Lip_d(f)=\sup_{x\ne y}\frac{|f(x)-f(y)|}{d(x,y)},
 \qquad
 \mathrm{BL}_b(S;\R)=\{f:S\to\R:\|f\|_\infty+\Lip_d(f)<\infty\}.
\]
（常数函数的 Lipschitz 常数按零计算；本小节固定 $d$ 后也把 $\Lip_d$ 简写为
$\Lip$。）再定义
\[
 d_{\mathrm{BL}}(\mu,\nu)=
 \sup\left\{
 \left|\int f\,\dd(\mu-\nu)\right|:
 \|f\|_\infty+\Lip_d(f)\le1
 \right\}.
 \tag{5.3.3}\label{eq:5-3-bl-metric}
\]
则 $d_{\mathrm{BL}}$ 是 $\mathsf P(S)$ 上的度量，并且它诱导由 $C_b(S)$ 测试定义的弱拓扑。
\end{proposition}

\begin{proof}
对称性与三角不等式从积分的相应性质得到。若
$d_{\mathrm{BL}}(\mu,\nu)=0$，则所有有界 Lipschitz 函数在两测度下积分相同。对开集
$G\ne S$，距离帽
$\phi_k=\min\{1,kd(\,\cdot\,,G^c)\}$ 是有界 Lipschitz 函数，且
$\phi_k\uparrow\1_G$；MCT 给 $\mu(G)=\nu(G)$。对 $G=S$，两者都取值一。
全体开集对有限交封闭并生成 $\Borel(S)$，有限测度唯一性定理遂给 $\mu=\nu$。所以
$d_{\mathrm{BL}}$ 的确是度量。

先证明弱拓扑中的每个 $\mu$ 都有落在给定 $d_{\mathrm{BL}}$-球内的弱开邻域。
记
\[
 \mathcal F_{\mathrm{BL}}
 =\{f:\|f\|_\infty+\Lip_d(f)\le1\}.
\]
给定 $0<\eta<1/2$ 与 $r>0$，由紧逼近取紧集 $K$ 使 $\mu(K^c)<\eta$，并令
$G=K^r=\{x:d(x,K)<r\}$。族 $\mathcal F_{\mathrm{BL}}|_K$ 一致有界且等度连续；
\ArzelaAscoli{} 给出有限 $\eta$-网 $f_1|_K,\ldots,f_N|_K$，其中
$f_j\in\mathcal F_{\mathrm{BL}}$。考虑弱开邻域
\[
 \mathcal U=\left\{\nu:
 \nu(G)>1-2\eta,\quad
 \left|\int f_j\,\dd(\nu-\mu)\right|<\eta
 \ (1\le j\le N)\right\}.
\]
这里 $\nu\mapsto\nu(G)$ 下半连续，故 $\mathcal U$ 确为弱开。

任取 $f\in\mathcal F_{\mathrm{BL}}$，选 $f_j$ 使
$\sup_K|f-f_j|<\eta$。若 $x\in G$，选 $y\in K$ 使 $d(x,y)<r$，则
\[
 |f(x)-f_j(x)|
 \le |f(x)-f(y)|+|f(y)-f_j(y)|+|f_j(y)-f_j(x)|
 <\eta+2r.
\]
又 $|f-f_j|\le2$。因此对 $\nu\in\mathcal U$，
\[
\begin{aligned}
 \left|\int f\,\dd(\nu-\mu)\right|
 &\le \left|\int f_j\,\dd(\nu-\mu)\right|
    +\int|f-f_j|\,\dd\nu+\int|f-f_j|\,\dd\mu\\
 &<\eta+(\eta+2r)+4\eta+\eta+2\eta
 =9\eta+2r.
\end{aligned}
\]
其中 $\nu(G^c)<2\eta$，而在 $\mu$ 的积分中直接在 $K$ 与 $K^c$ 上分裂。
先选 $\eta,r$ 使 $9\eta+2r$ 小于给定半径，就得到所需邻域。

反过来，证明每个 $C_b$ 测试积分关于 $d_{\mathrm{BL}}$ 连续。固定
$f\in C_b(S)$、$\mu\in\mathsf P(S)$，记 $M=\|f\|_\infty$。给定
$0<\eta<1/2$，取紧集 $K$ 使 $\mu(K^c)<\eta$。有界 Lipschitz 函数在紧度量空间
$K$ 上构成含常数并分离点的实函数代数；由 Stone--Weierstrass，可取
$g_0:K\to\R$ 为 Lipschitz 函数且
\[
 \sup_K|f-g_0|<\eta.
\]
记 $L_0=\Lip(g_0)$，并用 McShane 公式
\[
 \widetilde g(x)=\inf_{y\in K}\bigl(g_0(y)+L_0\,d(x,y)\bigr)
\]
延拓。对 $x\in K$，Lipschitz 不等式给
$g_0(y)+L_0d(x,y)\ge g_0(x)$，而取 $y=x$ 达到等号，所以
$\widetilde g|_K=g_0$。三角不等式还给
$|\widetilde g(x)-\widetilde g(x')|\le L_0d(x,x')$。再在
$[-M-\eta,M+\eta]$ 之间截断，得到
$g\in\mathrm{BL}_b(S)$，其在 $K$ 上仍等于 $g_0$。记 $L=\Lip(g)$。
连续性与 $K$ 的紧性保证存在 $r_0>0$，使得只要 $y\in K$ 且
$d(x,y)<r_0$，就有 $|f(x)-f(y)|<\eta$。否则可取
$y_n\in K$、$x_n\in S$，满足 $d(x_n,y_n)<1/n$ 而
$|f(x_n)-f(y_n)|\ge\eta$；从 $(y_n)$ 抽取收敛子列后，$(x_n)$ 收敛到同一极限，
与 $f$ 的连续性矛盾。因此可取 $0<r\le r_0$，同时满足：若
$y\in K$、$d(x,y)<r$，则 $|f(x)-f(y)|<\eta$，并且 $Lr<\eta$。
于是 $\sup_{K^r}|f-g|<3\eta$。

令 $\psi(x)=(1-d(x,K)/r)^+$。它有界 Lipschitz、在 $K$ 上等于一、在
$(K^r)^c$ 上等于零。由\eqref{eq:5-3-bl-metric}，对每个有界 Lipschitz $u$，
\[
 \left|\int u\,\dd(\nu-\mu)\right|
 \le(\|u\|_\infty+\Lip(u))\,d_{\mathrm{BL}}(\mu,\nu).
 \tag{5.3.4}\label{eq:5-3-bl-scaled}
\]
因此当 $d_{\mathrm{BL}}(\mu,\nu)$ 足够小时，
$\int\psi\,\dd\nu>1-2\eta$，从而 $\nu((K^r)^c)<2\eta$；同时
$|\int g\,\dd(\nu-\mu)|<\eta$。在 $K^r$ 及其补集上分裂后，
\[
\begin{aligned}
 \left|\int f\,\dd(\nu-\mu)\right|
 &\le\left|\int g\,\dd(\nu-\mu)\right|
   +\int|f-g|\,\dd\nu+\int|f-g|\,\dd\mu\\
 &\le \eta+6\eta
   +2(M+\|g\|_\infty)\bigl(\nu((K^r)^c)+\mu(K^c)\bigr),
\end{aligned}
\]
截断还给出 $\|g\|_\infty\le M+\eta$。代入
$\nu((K^r)^c)<2\eta$ 与 $\mu(K^c)<\eta$，上式右端不超过
\[
 7\eta+6\eta(2M+\eta),
\]
因而随 $\eta\downarrow0$ 而趋零。故 $d_{\mathrm{BL}}$ 收敛蕴含全部
$C_b$ 测试积分收敛。两个拓扑相同。
\end{proof}


\begin{proposition}[Daniell--\Caratheodory{} 表示步骤]\label{proposition:5-3-3-dc-step}
设 $S$ 为 Polish 空间，固定一个有界兼容完备度量 $d$，并令
\[
 H=\{f:S\to\R:\|f\|_\infty+\Lip_d(f)<\infty\}
   =\mathrm{BL}_b(S;\R).
\]
设
$L:H\to\R$ 是正线性泛函，$L(1)=1$，并满足下列 Daniell 单调连续性：
\[
 0\le f_k\downarrow0\text{ 点态}
 \quad\Longrightarrow\quad
 L(f_k)\downarrow0.
 \tag{5.3.5}\label{eq:5-3-daniell-continuity}
\]
则存在唯一 Radon 概率 $\mu$，使
\[
 L(f)=\int f\,\dd\mu\qquad(f\in H).
\]
\end{proposition}

\begin{proof}
下面直接作 Daniell--\Caratheodory{} 构造。对任意 $A\subset S$ 定义
\[
 \mu^*(A)=\inf\left\{
 \lim_{n\to\infty}L(g_n):
 0\le g_n\in H,\quad g_n\uparrow g,\quad g\ge\1_A
 \right\},
 \tag{5.3.6}\label{eq:5-3-dc-outer}
\]
其中极限函数 $g$ 允许取 $+\infty$，空下确界约定为 $+\infty$。有界
Lipschitz 函数对有限和、乘积、最大值和正部运算封闭；以下出现的
$\phi g_n$、$(1-g_n)^+$ 因而仍在 $H$ 中。
零序列给 $\mu^*(\varnothing)=0$。常数序列 $1$ 给 $\mu^*(S)\le1$。
反过来，若 $(g_n)$ 是 $S$ 的任一允许覆盖列，则
$(1-g_n)^+\downarrow0$。由\eqref{eq:5-3-daniell-continuity} 与
$1\le g_n+(1-g_n)^+$，
\[
 1=L(1)\le L(g_n)+L((1-g_n)^+)\longrightarrow\lim_nL(g_n).
\]
所以 $\mu^*(S)=1$。定义还直接给出单调性。

验证可列次可加性。设 $A=\bigcup_{k\ge1}A_k$，并先假设
$\sum_k\mu^*(A_k)<\infty$。给定 $\varepsilon>0$，为每个 $k$ 取允许列
$g_{k,n}\uparrow g_k\ge\1_{A_k}$，使
\[
 \lim_nL(g_{k,n})<\mu^*(A_k)+\varepsilon2^{-k}.
\]
这一可数次近优选择按第~1~章的约定使用 $\textsf{CC}$。把每列重复前若干项后可令双指标从
$n\ge1$ 开始，并置
\[
 h_n=\sum_{k=1}^n g_{k,n}.
\]
则 $h_n\uparrow h$，且 $h\ge\1_A$。又因每个
$L(g_{k,n})\le\lim_jL(g_{k,j})$，
\[
 \mu^*(A)\le\lim_nL(h_n)
 \le\sum_{k\ge1}\bigl(\mu^*(A_k)+\varepsilon2^{-k}\bigr).
\]
令 $\varepsilon\downarrow0$ 得次可加性；若右侧无穷，结论自动成立。因此 $\mu^*$ 是概率外测度。

它还是度量外测度。若 $A$ 或 $B$ 为空，正距离可加性已经由
$\mu^*(\varnothing)=0$ 得到。以下设二者非空且 $d(A,B)=\delta>0$，令
\[
 \phi(x)=\min\{1,d(x,B)/\delta\}.
\]
则 $\phi=1$ 于 $A$，$\phi=0$ 于 $B$。对任一覆盖 $A\cup B$ 的允许列
$g_n\uparrow g$，列 $\phi g_n$ 与 $(1-\phi)g_n$ 分别覆盖 $A,B$；线性给
\[
 \lim_nL(g_n)
 =\lim_nL(\phi g_n)+\lim_nL((1-\phi)g_n)
 \ge\mu^*(A)+\mu^*(B).
\]
与次可加性合并，得到
$\mu^*(A\cup B)=\mu^*(A)+\mu^*(B)$。由度量外测度的 Borel 性
命题~\ref{proposition:2-1-2-3}，所有 Borel 集都 $\mu^*$-可测。以下把 $\mu^*$ 在
$\Borel(S)$ 上的限制记为 $\mu$；它是 Borel 概率。

还须恢复泛函 $L$。先处理开集 $G$。若 $G\ne S$，距离帽
\[
 \phi_n(x)=\min\{1,nd(x,G^c)\}
\]
递增到 $\1_G$；$G=S$ 时取常数列一。它们在\eqref{eq:5-3-dc-outer} 中允许，故
\[
 \mu(G)\le
 \sup\{L(f):0\le f\le1,\ f\in H,\ f=0\text{ 于 }G^c\}.
 \tag{5.3.7}\label{eq:5-3-dc-open-lower}
\]
反之，固定右侧集合中的 $f$，并取任一覆盖 $G$ 的允许列 $g_n\uparrow g$。
由于 $f=0$ 于 $G^c$ 且 $g\ge1$ 于 $G$，有 $(f-g_n)^+\downarrow0$。于是
\[
 L(f)\le L(g_n)+L((f-g_n)^+)\longrightarrow\lim_nL(g_n).
\]
对覆盖列取下确界，再对 $f$ 取上确界，得到
\[
 \mu(G)=
 \sup\{L(f):0\le f\le1,\ f\in H,\ f=0\text{ 于 }G^c\}.
 \tag{5.3.8}\label{eq:5-3-dc-open-formula}
\]

现取 $0\le f\in H$，选 $M>\|f\|_\infty$。令
$\delta=2^{-m}$、$N_m=\lceil M/\delta\rceil$，并置开集
$G_k=\{f>k\delta\}$，$1\le k<N_m$。点态有
\[
 \delta\sum_{k=1}^{N_m-1}\1_{G_k}
 \le f
 \le\delta+\delta\sum_{k=1}^{N_m-1}\1_{G_k}.
 \tag{5.3.9}\label{eq:5-3-dc-layer-bounds}
\]
由\eqref{eq:5-3-dc-open-formula}，对每个 $k$ 可取
$0\le q_k\le1$、$q_k=0$ 于 $G_k^c$，使
$L(q_k)>\mu(G_k)-\varepsilon/N_m$。由于一个点上可能非零的 $q_k$
数目不超过满足 $k\delta<f(x)$ 的指标数，
$\delta\sum_kq_k\le f$。正性于是给
\[
 L(f)\ge
 \delta\sum_{k=1}^{N_m-1}\mu(G_k)-\delta\varepsilon.
 \tag{5.3.10}\label{eq:5-3-dc-lower-sum}
\]

为得反向估计，对每个 $G_k$ 取\eqref{eq:5-3-dc-outer} 中的允许列
$g_{k,n}$，其成本小于 $\mu(G_k)+\varepsilon/N_m$。令
$u_n=\delta+\delta\sum_kg_{k,n}$。由
\eqref{eq:5-3-dc-layer-bounds}，$(f-u_n)^+\downarrow0$，从而
\[
 L(f)\le\lim_nL(u_n)
 \le\delta+\delta\sum_{k=1}^{N_m-1}\mu(G_k)+\delta\varepsilon.
 \tag{5.3.11}\label{eq:5-3-dc-upper-sum}
\]
同一层集上下和也把 $\int f\,\dd\mu$ 夹在
\eqref{eq:5-3-dc-lower-sum}--\eqref{eq:5-3-dc-upper-sum} 的主项之间。
先令 $\varepsilon\downarrow0$，再令 $\delta=2^{-m}\downarrow0$，得到
$L(f)=\int f\,\dd\mu$。一般实值 $f\in H$ 由
$f=f^+-f^-$ 得到表示。

由\cref{proposition:5-3-3-polish-radon}，$\mu$ 是 Radon 概率。若另一概率也表示
$L$，则\eqref{eq:5-3-dc-open-formula} 说明二者在所有开集上相同；有限测度唯一性给二者相等。
\end{proof}

\begin{example}[单调连续性排除的伪极限]\label{ex:5-3-3-ultrafilter-limit}
在 $S=\N$ 上取离散度量 $d(m,n)=\1_{\{m\ne n\}}$。每个有界函数都是 Lipschitz 函数，所以
$\mathrm{BL}_b(S)=\ell^\infty$。取第~1~章所述的一个自由超滤子 $\mathcal U$，并定义
\[
 L(x)=\lim_{n\to\mathcal U}x_n,\qquad x\in\ell^\infty.
\]
有界数列的闭值域紧，故超滤极限存在；加法和数乘的连续性说明 $L$ 正、线性且 $L(1)=1$。
它很像概率积分，却不来自可数可加概率。令
\[
 f_k(n)=\1_{\{n\ge k\}}.
\]
对每个固定 $n$，有 $f_k(n)\downarrow0$；但每个余有限尾集都属于自由超滤子，所以
$L(f_k)=1$ 对所有 $k$ 成立。Daniell 单调连续性正好排除了这个伪极限。

这个例子也说明，仅要求“有一个函数在某个紧集上等于一且截断误差小”不够：恒取函数
$1$ 会使该条件没有内容。Prokhorov 证明中必须使用同一个真正的紧集控制原概率列，并在其上用
Dini 定理验证\eqref{eq:5-3-daniell-continuity}。
\end{example}

\begin{theorem}[Prokhorov 定理]\label{theorem:5-3-3-2}
设 $S$ 为 Polish 空间。概率族 $\mathcal K\subset\mathsf P(S)$ 相对弱紧，当且仅当
$\mathcal K$ 具有紧性。
\end{theorem}

\begin{proof}
先设 $\mathcal K$ 相对弱紧。其弱闭包 $\overline{\mathcal K}^{\,w}$ 紧，
\cref{proposition:5-3-3-compact-tight} 给出对整个闭包有效的共同紧集，因而
$\mathcal K$ 具有紧性。

反过来设 $\mathcal K$ 具有紧性，并固定一个有界兼容完备度量 $d$；以下的
$\Lip$ 与 $\mathrm{BL}_b$ 都相对于这个 $d$。取 $\mathcal K$ 中任意序列 $(\mu_n)$。对每个 $m$，选紧集
$K_m'$ 使
\[
 \sup_n\mu_n((K_m')^c)<2^{-m}.
\]
以有限并替换后可设 $K_m\uparrow$，且仍有
$\sup_n\mu_n(K_m^c)<2^{-m}$。这一步没有假设 $S$ 有紧耗尽；一般 Polish 空间未必
$\sigma$-紧，$K_m$ 只追踪这列概率的大部分质量。

我们构造可数测试族。对整数 $A,B,m\ge1$，$K_m$ 上满足
$\|f\|_\infty\le A$、$\Lip(f)\le B$ 的函数族在 $C(K_m)$ 中一致有界、等度连续，
因而由\cref{theorem:1-3-2-2}（\ArzelaAscoli{} 紧域定理）知该族全有界。对每个精度 $1/q$ 取一个有限网；把网中函数用
McShane 公式
\[
 \widetilde g(x)=\inf_{y\in K_m}\bigl(g(y)+B\,d(x,y)\bigr)
\]
延拓到 $S$，再截断到 $[-A,A]$。与
\cref{proposition:5-3-3-weak-metric} 证明中的计算相同，该延拓在
$K_m$ 上等于 $g$，且 Lipschitz 常数不超过 $B$；截断也不增大该常数。所得函数保持
$K_m$ 上的值。遍历 $A,B,m,q$ 得到一个可数族 $\mathcal D\subset\mathrm{BL}_b(S)$。
它的性质是：每个有界 Lipschitz $f$ 在每个 $K_m$ 上都可由
$\mathcal D$ 中具有共同取值界的函数一致逼近。

对 $\mathcal D$ 作对角抽取，得到子列 $(\mu_{n_j})$，使
$\int g\,\dd\mu_{n_j}$ 对每个 $g\in\mathcal D$ 都收敛。固定任意
$f\in\mathrm{BL}_b(S)$，取整数 $A,B$ 使
$\|f\|_\infty\le A$、$\Lip(f)\le B$。给定 $\varepsilon>0$，先取 $m$ 使
$4A2^{-m}<\varepsilon$，再取 $g\in\mathcal D$，使
$\|g\|_\infty\le A$ 且 $\sup_{K_m}|f-g|<\varepsilon$。于是
\[
\begin{aligned}
 \left|\int f\,\dd\mu_{n_j}-\int f\,\dd\mu_{n_\ell}\right|
 &\le
 \left|\int g\,\dd\mu_{n_j}-\int g\,\dd\mu_{n_\ell}\right|\\
 &\quad+2\varepsilon
   +2A\bigl(\mu_{n_j}(K_m^c)+\mu_{n_\ell}(K_m^c)\bigr).
\end{aligned}
\]
先令 $j,\ell\to\infty$，再令 $\varepsilon\downarrow0$，可知每个有界 Lipschitz
测试积分都有极限。定义
\[
 L(f)=\lim_{j\to\infty}\int f\,\dd\mu_{n_j},
 \qquad f\in\mathrm{BL}_b(S;\R).
\]
$L$ 正、线性且 $L(1)=1$。

关键是验证 Daniell 单调连续性。若 $0\le f_k\downarrow0$，记
$M=\|f_1\|_\infty$。给定 $\varepsilon>0$，由原列的紧性取紧集 $K$，使
$\sup_j\mu_{n_j}(K^c)<\varepsilon/(M+1)$。Dini 定理给
$\sup_K f_k\downarrow0$，故
\[
 0\le L(f_k)
 =\lim_j\int f_k\,\dd\mu_{n_j}
 \le \sup_K f_k+M\sup_j\mu_{n_j}(K^c).
\]
先令 $k\to\infty$，再令 $\varepsilon\downarrow0$，得到 $L(f_k)\downarrow0$。
\cref{proposition:5-3-3-dc-step} 因而产生唯一 Radon 概率 $\mu$，满足
$L(f)=\int f\,\dd\mu$ 对每个有界 Lipschitz $f$ 成立。

目前得到的是逐个 Lipschitz 测试函数的积分收敛，还需把它升级为
$d_{\mathrm{BL}}$ 中对整个单位球的一致收敛。给定 $\eta>0$，原列的紧性与
$\mu$ 的紧逼近给出同一个紧集 $K$，使
\[
 \sup_j\mu_{n_j}(K^c)<\eta,\qquad \mu(K^c)<\eta
\]
（取两个紧集的有限并）。有界 Lipschitz 单位球限制到 $K$ 后由
\ArzelaAscoli{} 有有限 $\eta$-网 $f_1,\ldots,f_N$。这些网点的积分都已收敛。
对任一单位球函数 $f$ 选 $f_i$ 使 $\sup_K|f-f_i|<\eta$；在 $K^c$ 上使用
$|f-f_i|\le2$，得到
\[
 \left|\int f\,\dd(\mu_{n_j}-\mu)\right|
 \le
 \left|\int f_i\,\dd(\mu_{n_j}-\mu)\right|+6\eta
\]
当 $j$ 足够大时对所有 $f$ 同时成立。故
$d_{\mathrm{BL}}(\mu_{n_j},\mu)\to0$，由
\cref{proposition:5-3-3-weak-metric}，$\mu_{n_j}\Rightarrow\mu$。
所以 $\mathcal K$ 中每个序列都有弱收敛子列。

最后把序列结论升级为闭包紧性。弱拓扑由 $d_{\mathrm{BL}}$ 度量。若
$(\lambda_n)$ 是 $\overline{\mathcal K}^{\,w}$ 中的序列，对每个 $n$ 选
$\nu_n\in\mathcal K$，使 $d_{\mathrm{BL}}(\lambda_n,\nu_n)<1/n$。
序列 $(\nu_n)$ 有收敛子列，相应的 $(\lambda_n)$ 子列收敛到同一极限。因此闭包序列紧；
由度量空间的序列判别\cref{theorem:1-2-1-4}，序列紧等价于紧，故
$\mathcal K$ 相对弱紧。
\end{proof}

\begin{example}[中心极限定理中的紧性]\label{ex:5-3-3-2}
设 $X_1,X_2,\ldots$ i.i.d.，$\E X_1=0$、$\E X_1^2=1$，并令
$Z_n=n^{-1/2}\sum_{k=1}^nX_k$。独立性给出 $\E Z_n^2=1$，所以对每个 $R>0$，
\[
 \sup_n\Prob(|Z_n|>R)\leq R^{-2}.
\]
因此分布族 $(\Law(Z_n))$ 具有紧性，Prokhorov 定理保证每个子列都有弱收敛的进一步子列。若再证明
$\varphi_{Z_n}(t)\to e^{-t^2/2}$，特征函数唯一性便把每个子列极限都识别为 $N(0,1)$。
\cref{subsec:5-4-3}中的 \Levy{} 连续性定理把这两步合并成一个通用判据。
\end{example}

\begin{proposition}[几乎处处连续的推前]\label{proposition:5-3-3-3}
设 $S,T$ 为 Polish 空间，$g:S\to T$ Borel 可测，并令 $D_g$ 为 $g$ 的不连续点集。
若 $\mu_n\Rightarrow\mu$ 且 $\mu(D_g)=0$，则
\[
 g_*\mu_n\Rightarrow g_*\mu.
\]
\end{proposition}

\begin{proof}
任取闭集 $F\subset T$。若
$x\in\overline{g^{-1}(F)}\setminus D_g$，则存在 $x_k\in g^{-1}(F)$ 收敛到 $x$；
$g$ 在 $x$ 连续，故 $g(x_k)\to g(x)$，而 $F$ 闭，所以 $g(x)\in F$。因此
\[
 \overline{g^{-1}(F)}
 \subset g^{-1}(F)\cup D_g.
\]
Portmanteau 的闭集不等式给
\[
\begin{aligned}
 \limsup_n(g_*\mu_n)(F)
 &=\limsup_n\mu_n(g^{-1}(F))\\
 &\le\limsup_n\mu_n(\overline{g^{-1}(F)})\\
 &\le\mu(\overline{g^{-1}(F)})
 \le\mu(g^{-1}(F))+\mu(D_g)\\
 &=(g_*\mu)(F).
\end{aligned}
\]
再次使用 Portmanteau，即得推前弱收敛。
\end{proof}

处处连续映射定理是这个命题中 $D_g=\varnothing$ 的特例；它在
\cref{theorem:5-3-1-3} 中已经由定义给出更短的证明。新命题真正增加的是：不连续点允许存在，但极限分布不能在那里放质量。

\begin{proposition}[Slutsky 定理]\label{proposition:5-3-3-slutsky}
设 $X_n\Rightarrow X$、$Y_n\xrightarrow{\Prob}c$，且 $X_n,Y_n$ 定义在同一概率空间上。
则
\[
 (X_n,Y_n)\Rightarrow(X,c).
\]
因此
\[
 X_n+Y_n\Rightarrow X+c,\qquad
 X_nY_n\Rightarrow cX.
\]
\end{proposition}

\begin{proof}
记 $\alpha_n=\Law(X_n)$、$\beta_n=\Law(Y_n)$ 与
$\lambda_n=\Law(X_n,Y_n)$。由
\cref{proposition:5-3-1-1}，$\beta_n\Rightarrow\delta_c$。
一个收敛列连同其极限在任意拓扑空间中都是紧的：包含极限的开集最终包含整条序列，剩余只有有限项。
所以
\[
 \{\alpha_n:n\ge1\}\cup\{\Law(X)\},
 \qquad
 \{\beta_n:n\ge1\}\cup\{\delta_c\}
\]
分别是弱紧集。由\cref{proposition:5-3-3-compact-tight}，给定
$\varepsilon>0$，存在紧集 $K,L\subset\R$，使
\[
 \sup_n\alpha_n(K^c)<\varepsilon/2,\qquad
 \sup_n\beta_n(L^c)<\varepsilon/2.
\]
$K\times L$ 紧，且
\[
 \lambda_n((K\times L)^c)
 \le\alpha_n(K^c)+\beta_n(L^c)<\varepsilon.
\]
故 $(\lambda_n)$ 具有紧性。

任取 $(\lambda_n)$ 的一个子列。Prokhorov 定理给出进一步子列弱收敛于某个
$\lambda\in\mathsf P(\R^2)$。坐标投影连续，故该极限的两个边缘分别为
$\Law(X)$ 与 $\delta_c$。第二边缘为 $\delta_c$ 意味着
\[
 \lambda(\R\times(\R\setminus\{c\}))=0,
\]
即 $\lambda$ 集中在 $\R\times\{c\}$。对任意 $h\in C_b(\R^2)$，
$x\mapsto h(x,c)$ 属于 $C_b(\R)$，于是由第一边缘
\[
 \int_{\R^2}h(x,y)\,\lambda(\dd x,\dd y)
 =\int_\R h(x,c)\,\Law(X)(\dd x).
\]
所以 $\lambda$ 唯一等于映射 $x\mapsto(x,c)$ 对 $\Law(X)$ 的推前，即
$\Law(X,c)$。

我们已经证明每个子列都有进一步子列弱收敛到同一个 $\Law(X,c)$。
由于弱拓扑由 $d_{\mathrm{BL}}$ 度量，若原列不收敛，就能选出一条与该极限距离至少为某个
$\eta>0$ 的子列；它的任何进一步收敛子列都不能以 $\Law(X,c)$ 为极限，矛盾。因此
$(X_n,Y_n)\Rightarrow(X,c)$。

最后把连续映射 $(x,y)\mapsto x+y$ 与 $(x,y)\mapsto xy$ 代入
\cref{proposition:5-3-3-3}，得到两个标量结论。
\end{proof}

\begin{proposition}[矩界推出紧性]\label{proposition:5-3-3-4}
设 $\mathcal K\subset\mathsf P(\R^d)$。若存在 $p>0$ 与 $C<\infty$，使
\[
 \sup_{\mu\in\mathcal K}\int_{\R^d}|x|^p\,\dd\mu(x)\le C,
\]
则 $\mathcal K$ 具有紧性。
\end{proposition}

\begin{proof}
给定 $R>0$，对每个 $\mu\in\mathcal K$，Markov 不等式给
\[
 \mu(\overline B(0,R)^c)
 =\mu(|x|^p>R^p)
 \le R^{-p}\int|x|^p\,\dd\mu
 \le CR^{-p}.
\]
给定 $\varepsilon>0$，取 $R>(C/\varepsilon)^{1/p}$。闭球
$\overline B(0,R)$ 在有限维欧氏空间中紧，并且对全族留下的集外质量小于
$\varepsilon$。
\end{proof}

\begin{corollary}[弱收敛列不会逃逸]\label{corollary:5-3-3-tight-sequence}
设 $S$ 为 Polish 空间，且 $\mu_n,\mu\in\mathsf P(S)$。若 $\mu_n\Rightarrow\mu$，则
\[
 \{\mu_n:n\ge1\}\cup\{\mu\}
\]
是紧概率族。
\end{corollary}

\begin{proof}
收敛列连同其极限在弱拓扑中紧；对这个紧集应用
\cref{proposition:5-3-3-compact-tight}。
\end{proof}

特别地，\cref{def:5-3-1-3} 中
\[
 X_n=O_\Prob(a_n)
\]
当且仅当归一化分布的整族
$\{\Law(X_n/a_n):n\ge1\}$ 具有紧性。事实上，$O_\Prob$ 定义对每个误差水平先控制一条尾列；
剩下的有限多个分布各自可由紧集逼近，扩大同一个紧区间即可把它们一并纳入。反之，整族
具有紧性时可在定义中取 $N=1$。在 $\R$ 上，紧区间 $[-M,M]$ 正好把定义中的尾概率写成
$\Prob(|X_n/a_n|>M)$。随机阶记号因而不是另一个极限概念，而是概率分布族的统一紧性估计。

\begin{figure}[htbp]
\centering
\begin{tikzpicture}[x=1cm,y=1cm]
  \begin{scope}
    \draw[book line] (0,0) ellipse (2.6 and 1.6);
    \node[book strong node,minimum width=2.8cm,minimum height=1.3cm] at (0,0)
      {共同紧集 $K$\\质量 $>1-\varepsilon$};
    \node[book label,align=center] at (0,2.05) {紧性：同一个 $K$ 控制整族};
    \draw[book accent arrow] (2.9,0.9)--(1.7,0.45);
    \node[book label,align=left] at (4.0,1.15) {集外质量\\统一小};
  \end{scope}
  \begin{scope}[xshift=7.7cm]
    \draw[book line] (-2.5,0)--(2.5,0);
    \draw[book line] (-0.8,-0.15)--(-0.8,1.7);
    \draw[book line] (0.9,-0.15)--(0.9,1.7);
    \draw[BookAccent,semithick] plot[smooth] coordinates {
      (-2.3,0) (-1.4,0) (-0.8,0.15) (-0.3,0.85) (0.2,1.35)
      (0.9,1.55) (1.5,1.55) (2.3,1.55)
    };
    \node[book label] at (-1.55,1.2) {$G^c$};
    \node[book label] at (1.55,1.85) {$G$};
    \node[book label,align=center] at (0,-0.65) {$\min\{1,kd(x,G^c)\}\uparrow\1_G$};
  \end{scope}
\end{tikzpicture}
\caption{左：Prokhorov 的质量控制。右：Portmanteau 中从下逼近开集指标的距离帽。有限网只在共同紧集上工作，图外质量由紧性控制。}
\label{fig:5-3-3-tight-portmanteau}
\end{figure}

弱收敛与紧性由此形成一条可执行路线。要证明一列分布有极限点，先找统一紧集；在
$\R^d$ 中常用矩界完成这一步。Prokhorov 只给存在性，随后仍需用特征函数、矩、方程或其他决定类识别所有可能的极限。反过来，一旦弱收敛已经成立，
\cref{corollary:5-3-3-tight-sequence} 会自动提供以后估计中所需的统一截断。
经验测度、数值近似的分布与随机过程的路径分布都遵循这条“先证紧性、再识别极限”的路线；路径空间上的紧性还需要第~6~章建立增量模量，不能仅由本节的有限维矩界替代。

\begin{remark}[边界条件不能省略]\label{rem:5-3-3-boundaries}
\begin{enumerate}
  \item 模糊收敛允许概率质量逃逸，不能替代概率弱收敛；
  \item 紧性是对整族的统一量词，不是“每个成员各自紧”；
  \item Prokhorov 的此种形式使用 Polish 结构：可分性提供可数对角测试，完备性把全有界闭集变成紧集；
  \item Daniell 单调连续性承担可数可加性。没有它，超滤极限可以通过全部有限线性检验；
  \item Prokhorov 只产生子列，不识别极限；唯一性仍需新的观测工具。
\end{enumerate}
\end{remark}

\begin{exercise}[Portmanteau 与紧性]\label{exercise:5-3-3-1}
\begin{enumerate}
  \item 证明 $\delta_{x_n}\Rightarrow\delta_x$ 当且仅当 $x_n\to x$，并给出
  $\delta_n$ 在 $\R$ 上不是紧族的量词证明。
  \item 由 Portmanteau 定理证明：若 $A$ 是 $\mu$-连续集且
  $\mu_n\Rightarrow\mu$，则 $\mu_n(A)\to\mu(A)$。
  \item 设 $\sup_n\E|X_n|^p<\infty$，其中 $p>0$。证明
  $\{\Law(X_n)\}$ 具有紧性，并计算使集外质量小于给定 $\varepsilon$ 所需的半径。
  \item 设 $g:S\to T$ 的不连续点集为 $D_g$。补全
  $\overline{g^{-1}(F)}\subset g^{-1}(F)\cup D_g$ 的逐点证明，并由闭集版
  Portmanteau 重证\cref{proposition:5-3-3-3}。
  \item 在离散空间 $\N$ 上验证
  $f_k=\1_{\{n\ge k\}}\downarrow0$，并解释为什么任意可数可加概率都满足
  $\int f_k\downarrow0$，而自由超滤极限不满足。
\end{enumerate}
\end{exercise}

现在我们已经能控制两种“逃逸”：UI 防止期望质量集中到越来越高的尖峰，紧性防止分布的质量逃到越来越远的空间位置。下一节把这些工具用于独立和：先以指数矩得到集中不等式，再证明大数律，最后用特征函数识别中心极限定理的唯一极限。

\section{独立和、强弱大数律与中心极限}\label{sec:05-04}

\subsection{矩母函数、Chernoff 方法与集中不等式}\label{subsec:5-4-1}
\index{矩母函数}\index{Chernoff 方法}\index{集中不等式}

Chebyshev 不等式已经能回答一个基本问题。若 $X_1,\dots,X_n$ 独立同分布，
$\E X_i=\mu$、$\Var(X_i)=\sigma^2<\infty$，则
\[
 \Prob\left(\left|\frac1n\sum_{i=1}^nX_i-\mu\right|\geq\varepsilon\right)
 \leq \frac{\sigma^2}{n\varepsilon^2}.
\]
这个上界随 $n$ 衰减，却只按 $n^{-1}$ 衰减。对有界独立变量，实际尾概率通常按
$e^{-cn}$ 衰减。缺失的指数从哪里来？把和 $S$ 直接放进 Markov 不等式只会使用一阶矩；
把 $S$ 先送进单调函数 $x\mapsto e^{sx}$，则独立性会把和的指数变成可分解的乘积。
这一小步就是 Chernoff 方法的核心机制。

Chernoff 的指数变换、Hoeffding 的有界差估计与 Bernstein 的方差信息修正都遵循同一逻辑。
它们的共同点不是某个需背诵的公式，而是“指数化、因子化、
控制单项、优化参数”这组通用证明步骤。

在正式推导之前，先看三种截然不同的指数矩行为。

\begin{example}[Rademacher]\label{ex:5-4-1-1}
设 $X$ 以相同概率取 $1$ 与 $-1$。则 $\E X=0$，并且对每个 $t\in\R$，
\[
 \E e^{tX}=\frac{e^t+e^{-t}}2=\cosh t.
\]
令 $h(t)=\log\cosh t-t^2/2$。由于
$h'(t)=\tanh t-t$，而当 $t\geq0$ 时
$h''(t)=\operatorname{sech}^2t-1\leq0$，故 $h'(t)\leq h'(0)=0$，从而
$h(t)\leq h(0)=0$。函数 $h$ 为偶函数，所以
\[
 \cosh t\leq e^{t^2/2}\qquad(t\in\R).
\]
这正是下文定义的最基本的次高斯性质。
\end{example}

\begin{example}[高斯分布]\label{ex:5-4-1-2}
设 $G\sim N(0,\sigma^2)$，其中 $\sigma>0$。只在实轴上配方，便有
\begin{align*}
 \E e^{tG}
 &=\frac1{\sqrt{2\pi\sigma^2}}
   \int_{\R}\exp\left(tx-\frac{x^2}{2\sigma^2}\right)\dd x\\
 &=e^{\sigma^2t^2/2}\frac1{\sqrt{2\pi\sigma^2}}
   \int_{\R}\exp\left(-\frac{(x-\sigma^2t)^2}{2\sigma^2}\right)\dd x
 =e^{\sigma^2t^2/2}.
\end{align*}
第二个等号只是实变量平移 $u=x-\sigma^2t$；这里没有复变换元。高斯分布使
下文的次高斯上界取等，是以后比较常数的基准。
\end{example}

\begin{example}[Pareto 重尾变量]\label{ex:5-4-1-3}
令 $X$ 的密度为
\[
 p(x)=\alpha x^{-\alpha-1}\1_{[1,\infty)}(x),\qquad \alpha>0.
\]
首先
\[
 \int_1^\infty\alpha x^{-\alpha-1}\dd x
 =\bigl[-x^{-\alpha}\bigr]_{1}^{\infty}=1,
\]
所以这确实是概率密度。且对 $x\ge1$，
\[
 \Prob(X>x)=\int_x^\infty\alpha u^{-\alpha-1}\dd u=x^{-\alpha}.
\]
对任意 $t>0$，
\[
 \E e^{tX}
 \geq\alpha\sum_{n=1}^{\infty}
       \int_n^{n+1}e^{tx}x^{-\alpha-1}\dd x
 \geq\alpha\sum_{n=1}^{\infty}e^{tn}(n+1)^{-\alpha-1}=\infty,
\]
因为最后一级数的通项甚至不趋于零。因此正指数矩处处发散，Chernoff 方法在这里无从启动。
这不是计算技巧不足，而是重尾分布本身不具有有限的正指数矩；此时应改用截断或幂矩估计。
\end{example}

\begin{definition}[矩母函数与累积量生成函数]\label{def:5-4-1-1}
实值随机变量 $X$ 的矩母函数定义为
\[
 M_X(t)=\E e^{tX}\in(0,\infty],\qquad t\in\R.
\]
在 $M_X(t)<\infty$ 的参数集合上，记
$\Lambda_X(t)=\log M_X(t)$，称为\emph{累积量生成函数}。
在某一个正参数处有限已经足以给相应的单侧 Chernoff 界；要求零点邻域内处处有限是更强的条件，
二者不作混同。
\end{definition}

\begin{definition}[次高斯变量]\label{def:5-4-1-2}
\index{次高斯变量}
若实值随机变量 $X$ 满足 $\E X=0$，并且存在 $\sigma^2\geq0$ 使
\[
 \E e^{tX}\leq e^{\sigma^2t^2/2}\qquad(t\in\R),
\]
则称 $X$ 是参数 $\sigma^2$ 的\emph{次高斯}（sub-Gaussian）变量。这个参数是指数矩的控制常数，
并不按定义等于 $\Var(X)$。
\end{definition}

\begin{definition}[Chernoff 变换]\label{def:5-4-1-3}
设随机变量 $S$ 的对数矩母函数为 $\Lambda_S$。对 $a\in\R$，定义单侧 Chernoff 变换
\[
 I_S^+(a)=\sup_{\substack{s>0\\\Lambda_S(s)<\infty}}
 \{sa-\Lambda_S(s)\}.
\]
若上确界所取集合为空，则约定 $I_S^+(a)=-\infty$。
\end{definition}

\begin{proposition}[Chernoff 上界]
若至少存在一个 $s>0$ 使 $\E e^{sS}<\infty$，则
\[
 \Prob(S\geq a)
 \leq\inf_{\substack{s>0\\\Lambda_S(s)<\infty}}
 e^{-sa}\E e^{sS}
 =e^{-I_S^+(a)}.
\]
\end{proposition}

\begin{proof}
对每个可行的 $s>0$，指数函数的单调性和 Markov 不等式给出
\[
 \Prob(S\geq a)
 =\Prob(e^{sS}\geq e^{sa})
 \leq e^{-sa}\E e^{sS}.
\]
对 $s$ 取下确界，再用 $\Lambda_S(s)=\log\E e^{sS}$，便得到结论。
\end{proof}

真正需要估计的是单个有界中心变量的指数矩。凸性用两个端点的凸组合控制区间内的指数函数，而端点计算仍需一个
不会把最优常数藏起来的微积分步骤。

\begin{lemma}[Hoeffding 引理]\label{lemma:5-4-1-1}
设 $a\leq X\leq b$ a.s. 且 $\E X=0$。则对每个 $t\in\R$，
\[
 \E e^{tX}\leq\exp\left(\frac{t^2(b-a)^2}{8}\right).
\]
\end{lemma}

\begin{proof}
若 $a=b$，则 $X=0$ a.s.，结论成立。以下设 $a<b$。由 $\E X=0$ 可知
$a\leq0\leq b$。凸函数 $x\mapsto e^{tx}$ 在区间端点之间低于弦，因而对
$x\in[a,b]$，
\[
 e^{tx}\leq \frac{b-x}{b-a}e^{ta}+\frac{x-a}{b-a}e^{tb}.
\]
取期望并用 $\E X=0$，得到
\[
 \E e^{tX}\leq\frac{b}{b-a}e^{ta}-\frac{a}{b-a}e^{tb}.
\]
置
\[
 p=\frac{-a}{b-a}\in[0,1],\qquad s=t(b-a).
\]
右端变为
\[
 (1-p)e^{-ps}+pe^{(1-p)s}.
\]
令
\[
 H(s)=\log\bigl((1-p)e^{-ps}+pe^{(1-p)s}\bigr).
\]
把分母中的两项归一化为概率 $1-q_s,q_s$，其中
\[
 q_s=\frac{pe^{(1-p)s}}
 {(1-p)e^{-ps}+pe^{(1-p)s}},
\]
直接求导得到
\[
 H'(s)=q_s-p,\qquad H''(s)=q_s(1-q_s)\leq\frac14.
\]
又 $H(0)=H'(0)=0$。当 $s\geq0$ 时，Taylor 积分余项给
\[
 H(s)=\int_0^s(s-u)H''(u)\dd u\leq\frac{s^2}{8};
\]
当 $s<0$ 时，同一恒等式写成
$H(s)=\int_s^0(u-s)H''(u)\dd u$，仍得到相同上界。因此
$H(s)\leq s^2/8$ 对所有实数 $s$ 成立。代回 $s=t(b-a)$ 即得结论。
\end{proof}

独立性现在把每一项的估计相乘起来。

\begin{theorem}[Hoeffding 不等式]\label{theorem:5-4-1-2}
设 $X_1,\dots,X_n$ 独立，且
$a_i\leq X_i\leq b_i$ a.s.。则对每个 $u>0$，
若 $V:=\sum_{i=1}^n(b_i-a_i)^2>0$，则
\[
 \Prob\left(\sum_{i=1}^n(X_i-\E X_i)\geq u\right)
 \leq
 \exp\left(-\frac{2u^2}{V}\right).
\]
若 $V=0$，则左端为零。
\end{theorem}

\begin{proof}
记 $Y_i=X_i-\E X_i$ 以及
$V=\sum_i(b_i-a_i)^2$。若 $V=0$，则每个 $X_i$ a.s. 为常数，因而
$\sum_iY_i=0$ a.s.；结论成立。以下设 $V>0$。对任意 $s>0$，Chernoff 变换、
独立变量的期望因子化以及 Hoeffding 引理依次给出
\begin{align*}
 \Prob\left(\sum_iY_i\geq u\right)
 &\leq e^{-su}\E\exp\left(s\sum_iY_i\right)\\
 &=e^{-su}\prod_i\E e^{sY_i}\\
 &\leq\exp\left(-su+\frac{s^2}{8}
                  \sum_i(b_i-a_i)^2\right)
 =\exp\left(-su+\frac{s^2V}{8}\right).
\end{align*}
右端指数是关于 $s$ 的二次函数，在 $s=4u/V$ 处达到最小值
$-2u^2/V$。代入即得所述估计。
\end{proof}

\begin{corollary}[样本均值的双侧界]\label{corollary:5-4-1-hoeffding-two-sided}
若 $X_1,\dots,X_n$ 独立同分布，且 $a\leq X_i\leq b$ a.s.，其中 $a<b$，则
对每个 $\varepsilon>0$，
\[
 \Prob\left(
 \left|\frac1n\sum_{i=1}^nX_i-\E X_1\right|\geq\varepsilon
 \right)
 \leq2\exp\left(-\frac{2n\varepsilon^2}{(b-a)^2}\right).
\]
\end{corollary}

\begin{proof}
对上尾把定理中的 $u$ 取为 $n\varepsilon$。对变量 $-X_i$ 再应用一次定理得到下尾估计，
最后使用并集界。
\end{proof}

这个推论立即给出一个强结论。固定 $m\geq1$，右端在 $n$ 上可和，故 Borel--Cantelli
引理说明样本均值与期望相差至少 $1/m$ 的事件只发生有限次。对所有 $m$ 取可数交，得到
有界独立同分布变量的样本均值 a.s. 收敛到期望。这里不能只写“令 $\varepsilon\downarrow0$”：
把不可数多个满测集相交并不保测度一，可数阈值 $1/m$ 才完成这一步。

\begin{proposition}[次高斯变量之和]\label{proposition:5-4-1-3}
若 $X_1,\dots,X_n$ 独立、中心化，且 $X_i$ 是参数 $\sigma_i^2$ 的
次高斯变量，则 $\sum_iX_i$ 是参数 $\sum_i\sigma_i^2$ 的次高斯变量。
特别地，若 $v:=\sum_i\sigma_i^2>0$，则对每个 $u>0$，
\[
 \Prob\left(\sum_iX_i\geq u\right)
 \leq\exp\left(-\frac{u^2}{2v}\right).
\]
若 $v=0$，则 $\sum_iX_i=0$ a.s.
\end{proposition}

\begin{proof}
独立性与期望因子化给出
\[
 \E\exp\left(t\sum_iX_i\right)
 =\prod_i\E e^{tX_i}
 \leq\exp\left(\frac{t^2}{2}\sum_i\sigma_i^2\right).
\]
这就是第一个结论。若 $v=\sum_i\sigma_i^2>0$，Chernoff 上界的指数为
$-tu+t^2v/2$，取 $t=u/v$ 得尾界。若 $v=0$，则每个 $\sigma_i^2=0$，且
$\E e^{X_i}\leq1$。另一方面，$e^x\geq1+x$ 对所有实数 $x$ 成立，并且只在 $x=0$ 处取等，故
\[
 0\leq\E(e^{X_i}-1-X_i)\leq1-1-\E X_i=0.
\]
于是 $e^{X_i}-1-X_i=0$ a.s.，从而 $X_i=0$ a.s.。所以和为零 a.s.，所述上尾概率为零。
\end{proof}

Hoeffding 不等式只利用区间长度，没有利用实际方差。若方差远小于 $M^2$，下一条估计能把这份信息
放回指数中。

\begin{proposition}[Bernstein 型界（选读）]\label{proposition:5-4-1-4}
设 $X_1,\dots,X_n$ 为独立实值随机变量，$\E X_i=0$ 且
$|X_i|\leq M$ a.s.，其中 $M>0$。令
\[
 S=\sum_{i=1}^nX_i,\qquad v=\sum_{i=1}^n\E X_i^2.
\]
则对每个 $u>0$，
\[
 \Prob(S\geq u)
 \leq\exp\left(-\frac{u^2}{2(v+Mu/3)}\right).
\]
\end{proposition}

\begin{proof}
若 $0<s<3/M$，由 $\E X_i=0$、$|X_i|\leq M$ 及幂级数的绝对收敛，
\begin{align*}
 \E e^{sX_i}
 &=1+\sum_{k=2}^{\infty}\frac{s^k\E X_i^k}{k!}\\
 &\leq1+\sum_{k=2}^{\infty}
       \frac{s^kM^{k-2}\E X_i^2}{k!}.
\end{align*}
对 $k\geq2$ 有 $k!\geq2\,3^{k-2}$，所以
\[
 \E e^{sX_i}
 \leq1+\frac{s^2\E X_i^2}{2}
        \sum_{j=0}^{\infty}\left(\frac{sM}{3}\right)^j
 =1+\frac{s^2\E X_i^2}{2(1-sM/3)}.
\]
由 $\log(1+x)\leq x$，
\[
 \log\E e^{sX_i}
 \leq\frac{s^2\E X_i^2}{2(1-sM/3)}.
\]
若 $v=0$，则 $\E X_i^2=0$ 对每个 $i$ 成立，因而 $X_i=0$ a.s.，所述尾概率为零。
以下设 $v>0$。独立性与 Chernoff
变换给
\[
 \Prob(S\geq u)
 \leq\exp\left(-su+\frac{s^2v}{2(1-sM/3)}\right).
\]
取
$s=u/(v+Mu/3)$。此时 $sM/3=Mu/(3v+Mu)<1$。记 $D=v+Mu/3$，则
$s=u/D$ 且 $1-sM/3=v/D$，因而
\begin{align*}
 -su+\frac{s^2v}{2(1-sM/3)}
 &=-\frac{u^2}{D}
   +\frac{(u^2/D^2)v}{2(v/D)}\\
 &=-\frac{u^2}{D}+\frac{u^2}{2D}
 =-\frac{u^2}{2(v+Mu/3)}.
\end{align*}
\end{proof}

三种尾界的尺度现在可以直接比较。对 $[a,b]$ 中的 i.i.d. 变量，Chebyshev 给
$O((n\varepsilon^2)^{-1})$，Hoeffding 给
$2e^{-2n\varepsilon^2/(b-a)^2}$；Bernstein 在小偏差区近似
$e^{-n\varepsilon^2/(2\Var X_1)}$，在大偏差区转成 $e^{-cn\varepsilon/M}$。
例如若要以至少 $1-\delta$ 的置信度保证样本均值误差小于 $\varepsilon$，Hoeffding 给出充分条件
\[
 n\geq\frac{(b-a)^2}{2\varepsilon^2}\log\frac2\delta.
\]
这就是经验频率和有界独立抽样平均的有限样本保证。

\begin{center}
\begin{tikzpicture}[node distance=8mm and 9mm]
 \node[book strong node] (tail) {尾事件\\$\{S\geq u\}$};
 \node[book node,right=of tail] (exp) {指数化\\$e^{sS}\geq e^{su}$};
 \node[book node,right=of exp] (prod) {独立分解\\$\E e^{sS}=\prod_i\E e^{sX_i}$};
 \node[book warm node,below=of prod] (quad) {单项控制\\$\log\E e^{sX_i}\leq c_i(s)$};
 \node[book strong node,left=of quad] (opt) {参数优化\\$\inf_{s>0}$};
 \draw[book accent arrow] (tail)--(exp);
 \draw[book accent arrow] (exp)--(prod);
 \draw[book accent arrow] (prod)--(quad);
 \draw[book accent arrow] (quad)--(opt);
\end{tikzpicture}
\end{center}

图中的四步分别使用单调性、独立性、单变量估计与实变量优化。任何一步的假设失效，都不能靠
“集中现象应当成立”的直觉补齐。Pareto 例使第一项指数矩已经发散；相关变量使乘积公式失效；
次高斯参数也不是通常方差的另一个名字。

若失去独立性，上述因子化必须由针对具体依赖结构的估计取代。本小节的结论只针对独立变量。

\begin{remark}[本章只使用实参数矩母函数]\label{rem:5-4-1-complex-mgf}
本章不定义复参数矩母函数，也不调用其全纯性理论。Chernoff 估计只需要实数 $t>0$；
对 Rademacher 变量，实参数矩母函数是
$M(t)=(e^t+e^{-t})/2=\cosh t$，而 Pareto 变量对每个 $t>0$ 都有
$M(t)=\infty$。与之相反，第~\ref{subsec:5-4-3}~节所讨论的特征函数
$\E e^{itX}$ 因 $|e^{itX}|=1$ 而对每个实数 $t$ 都存在。
\end{remark}

\begin{exercise}[Bernoulli 平均的集中估计]
设 $B_1,\dots,B_n$ 为独立 Bernoulli$(p)$ 变量。
\begin{enumerate}
 \item 由 Hoeffding 不等式证明
 $\Prob(|n^{-1}\sum_iB_i-p|\geq\varepsilon)\leq2e^{-2n\varepsilon^2}$；
 \item 比较同一问题的 Chebyshev 上界，并求出两种方法为达到误差概率 $\delta$ 所给的样本量；
 \item 验证若去掉独立性并令 $B_i=B_1$，Hoeffding 的 $e^{-cn}$ 结论可能失败。
\end{enumerate}
\end{exercise}

\subsection{弱大数律、强大数律与截断方法}\label{subsec:5-4-2}
\index{弱大数律}\index{强大数律}\index{Kolmogorov 最大不等式}

集中不等式同时控制有限样本和渐近行为，但它要求有界性或指数矩。若只有二阶矩，
Chebyshev 不等式仍能给依概率收敛；要把结论升级为 a.s. 收敛，逐个时刻的概率趋零还不够，
必须控制一整段路径中的最大振荡。若进一步只保留一阶绝对矩，连方差都可能不存在，
于是还要先把罕见的大跳截掉。本小节依次完成这三次升级。

Bernoulli 的频率定理给出了大数律的早期形式；Borel 把结论推进到路径上的强收敛，Kolmogorov
又用最大不等式与截断给出自然的 $L^1$ 边界。下面的证明顺序保留这条技术升级：先缩放方差，
再控制二进分块，最后截去稀有大值。

\begin{definition}[样本和与样本平均]\label{def:5-4-2-1}
对随机变量序列 $(X_n)$，记
\[
 S_n=\sum_{k=1}^nX_k,
 \qquad
 \overline X_n=\frac{S_n}{n}.
\]
缩写 i.i.d. 表示这些变量相互独立且同分布。
\end{definition}

先从只需一行方差计算的结论开始。

\begin{theorem}[弱大数律（有限方差）]\label{theorem:5-4-2-1}
若 $(X_n)$ i.i.d.，$\E X_1=\mu$ 且 $\Var(X_1)=\sigma^2<\infty$，则
\[
 \frac{S_n}{n}\xrightarrow{\Prob}\mu.
\]
\end{theorem}

\begin{proof}
独立性和中心化给出
\[
 \Var\left(\frac{S_n}{n}\right)
 =\frac1{n^2}\sum_{k=1}^n\Var(X_k)
 =\frac{\sigma^2}{n}.
\]
因此对每个 $\varepsilon>0$，Chebyshev 不等式给
\[
 \Prob\left(\left|\frac{S_n}{n}-\mu\right|\geq\varepsilon\right)
 \leq\frac{\sigma^2}{n\varepsilon^2}\longrightarrow0.
\]
\end{proof}

这个证明不能直接给强律，因为 $\sum_n n^{-1}$ 发散，Borel--Cantelli 无法读取这组上界。
把时刻稀疏到 $2^k$ 后，上界变成可和；随之产生的新问题是如何控制相邻两个二进时刻之间的所有中间时刻。

\begin{example}[Bernoulli 频率]\label{ex:5-4-2-1}
设 $X_i=\1_{A_i}$ i.i.d.，且 $\Prob(A_i)=p$。由推论
\ref{corollary:5-4-1-hoeffding-two-sided}，
\[
 \Prob\left(\left|\frac{S_n}{n}-p\right|>\varepsilon\right)
 \leq2e^{-2n\varepsilon^2}.
\]
对每个 $\varepsilon=1/m$，右端在 $n$ 上可和。Borel--Cantelli I 和可数交于是给出
$S_n/n\to p$ a.s.。在这个有界模型中，集中估计已经给出了频率的强收敛。
\end{example}

可积性是一般强律的边界条件。标准 Cauchy 分布给出一个可计算的反例；其卷积稳定性将在
\cref{ex:5-4-3-cauchy-stability}中用纯实变量方法证明。

\begin{example}[Cauchy 反例]\label{ex:5-4-2-2}
标准 Cauchy 分布的密度为
\[
 f_1(x)=\frac1{\pi(1+x^2)}.
\]
归一化由
\[
 \int_{\R}f_1(x)\,\dd x
 =\frac1\pi\bigl[\arctan x\bigr]_{-\infty}^{\infty}=1
\]
得到。它不满足 $\E|X|<\infty$，因为
\[
 \int_1^R\frac{x}{\pi(1+x^2)}\,\dd x
 =\frac1{2\pi}\log\frac{1+R^2}{2}\longrightarrow\infty;
\]
由密度的对称性，正部和负部的期望
也都为无穷，所以 $\E X$ 不存在。例~\ref{ex:5-4-3-cauchy-stability} 将直接证明：若
$X_1,\dots,X_n$ 独立且均为标准 Cauchy，则
\[
 \frac1n\sum_{k=1}^nX_k\overset{\mathrm d}=X_1
\]
对每个 $n$ 成立。于是这些平均不可能依概率收敛到某个常数：若它们依概率收敛到 $c$，
便也依分布收敛到 $c$，但其分布恒为非退化 Cauchy 分布。这个例子说明删去可积性后强律可能失败；
它没有声称每个不可积分布都以同一种方式失败。
\end{example}

\begin{example}[有限方差的二进尺度]\label{ex:5-4-2-3}
令 $T_n=S_n-n\mu$。Chebyshev 在 $n=2^k$ 时给
\[
 \Prob(|T_{2^k}|>\varepsilon2^k)
 \leq\frac{\sigma^2}{\varepsilon^2 2^k},
\]
其和有限。故 $T_{2^k}/2^k\to0$ a.s.。然而这只控制二进时刻；
在区间 $[2^k,2^{k+1}]$ 内仍可能出现大幅偏离。下面的最大不等式正是填补这些区间的工具。
\end{example}

\begin{theorem}[Kolmogorov 最大不等式]\label{theorem:5-4-2-2}
设 $X_1,\dots,X_n$ 独立、实值、中心化且属于 $L^2$，并令
$S_k=\sum_{j=1}^kX_j$。则对每个 $\lambda>0$，
\[
 \Prob\left(\max_{1\leq k\leq n}|S_k|\geq\lambda\right)
 \leq\frac{\Var(S_n)}{\lambda^2}.
\]
\end{theorem}

\begin{proof}
对 $1\leq k\leq n$，令
\[
 A_k=\{|S_1|<\lambda,\dots,|S_{k-1}|<\lambda,
            |S_k|\geq\lambda\}.
\]
这些首次越界事件两两不交，其并为
$\{\max_{j\leq n}|S_j|\geq\lambda\}$。事件 $A_k$ 与 $S_k$ 都对
$\sigma(X_1,\dots,X_k)$ 可测，而增量 $S_n-S_k$ 与该 $\sigma$-代数独立并且均值为零。
由 Cauchy--Schwarz，不同因子的乘积可积；再用独立变量的期望因子化，得到
\[
 \E\bigl[S_k(S_n-S_k)\1_{A_k}\bigr]
 =\E[S_k\1_{A_k}]\,\E[S_n-S_k]=0.
\]
展开平方并使用最后一项非负，得到
\begin{align*}
 \E[S_n^2\1_{A_k}]
 &=\E[S_k^2\1_{A_k}]
   +2\E[S_k(S_n-S_k)\1_{A_k}]
   +\E[(S_n-S_k)^2\1_{A_k}]\\
 &\geq\E[S_k^2\1_{A_k}]
 \geq\lambda^2\Prob(A_k).
\end{align*}
对 $k$ 求和，并利用 $A_k$ 两两不交，
\[
 \lambda^2\Prob\left(\max_{j\leq n}|S_j|\geq\lambda\right)
 \leq\sum_{k=1}^n\E[S_n^2\1_{A_k}]
 \leq\E S_n^2=\Var(S_n).
\]
除以 $\lambda^2$ 即得结论。
\end{proof}

\begin{center}
\begin{tikzpicture}[x=1.05cm,y=0.8cm]
 \draw[->,book line] (0,0)--(9.5,0) node[right] {$k$};
 \draw[BookWarning,dashed] (0,2.6)--(9.1,2.6) node[right] {$+\lambda$};
 \draw[BookWarning,dashed] (0,-2.6)--(9.1,-2.6) node[right] {$-\lambda$};
 \draw[book accent arrow] (0,0)--(1,0.6)--(2,0.2)--(3,1.1)--(4,1.7)--(5,2.9)--(6,2.2)--(7,3.1)--(8,2.5)--(9,3.3);
 \fill[BookWarning] (5,2.9) circle (2pt);
 \draw[dashed,BookWarning] (5,0)--(5,2.9);
 \node[book warning node,above=3pt] at (5,2.9) {$A_5$：首次越界};
 \node[book node] at (7.2,0.9) {越界后的增量\\与过去独立且中心化};
\end{tikzpicture}
\end{center}

图只解释事件分割；交叉项为何消失仍由上面的可测性、独立性和条件期望计算保证。

\begin{theorem}[强大数律（有限方差）]\label{theorem:5-4-2-3}
若 $(X_n)$ i.i.d. 且 $X_1\in L^2$，则令 $\mu=\E X_1$，有
\[
 \frac{S_n}{n}\longrightarrow\mu\qquad\text{a.s.}
\]
\end{theorem}

\begin{proof}
置 $W_i=X_i-\mu$、$T_n=\sum_{i=1}^nW_i$，并记
$\sigma^2=\E W_1^2$。固定 $\varepsilon>0$。上一例已经给出
\[
 \sum_{k=0}^{\infty}
 \Prob(|T_{2^k}|>\varepsilon2^k)<\infty.
\]
另一方面，把最大不等式应用于独立块
$W_{2^k+1},\dots,W_{2^{k+1}}$，得到
\begin{align*}
 &\Prob\left(
 \max_{1\leq j\leq2^k}
 |T_{2^k+j}-T_{2^k}|>\varepsilon2^k
 \right)\\
 &\hspace{35mm}\leq
 \frac{2^k\sigma^2}{\varepsilon^2 2^{2k}}
 =\frac{\sigma^2}{\varepsilon^2 2^k}.
\end{align*}
这组概率同样可和。先令 $\varepsilon$ 遍历 $1/m$，再用 Borel--Cantelli I 并取可数交，
得到一个概率一事件；在该事件上，
\[
 \frac{|T_{2^k}|}{2^k}\longrightarrow0,
 \qquad
 \frac1{2^k}\max_{1\leq j\leq2^k}
 |T_{2^k+j}-T_{2^k}|\longrightarrow0.
\]
若 $2^k\leq n<2^{k+1}$，则
\[
 \frac{|T_n|}{n}
 \leq\frac{|T_{2^k}|}{2^k}
 +\frac1{2^k}\max_{1\leq j\leq2^k}
 |T_{2^k+j}-T_{2^k}|.
\]
右端趋于零，因此 $T_n/n\to0$，也就是 $S_n/n\to\mu$ a.s.
\end{proof}

有限方差并不是强律的最终边界。为了把条件降到 $L^1$，需要把一个随机级数的收敛转成
加权平均的收敛。下面的实分析引理承担这个转换。

\begin{lemma}[Kronecker 引理]\label{lemma:5-4-2-3}
设 $0<a_1\leq a_2\leq\cdots$ 且 $a_n\to\infty$。若实数或复数级数
$\sum_{n=1}^{\infty}x_n/a_n$ 收敛，则
\[
 \frac1{a_n}\sum_{k=1}^nx_k\longrightarrow0.
\]
\end{lemma}

\begin{proof}
令 $b_n=\sum_{k=1}^nx_k/a_k$、$b_0=0$，并设 $b_n\to b$。Abel 分部求和给
\[
 \sum_{k=1}^nx_k
 =\sum_{k=1}^na_k(b_k-b_{k-1})
 =a_nb_n-\sum_{k=1}^{n-1}(a_{k+1}-a_k)b_k.
\]
把 $b_k=b+(b_k-b)$ 代入并除以 $a_n$。常数 $b$ 的贡献为
$a_1b/a_n\to0$。再给定 $\varepsilon>0$，取 $K$ 使
$|b_k-b|<\varepsilon$ 对 $k\geq K$ 成立。有限前段的贡献除以 $a_n$ 后趋于零，而尾段满足
\[
 \frac1{a_n}\sum_{k=K}^{n-1}(a_{k+1}-a_k)|b_k-b|
 \leq\varepsilon\frac{a_n-a_K}{a_n}\leq\varepsilon.
\]
同时 $|b_n-b|\to0$。先令 $n\to\infty$，再令 $\varepsilon\downarrow0$，得到结论。
\end{proof}

\begin{theorem}[Kolmogorov 收敛判据]\label{theorem:5-4-2-kolmogorov-series}
设 $(Z_n)$ 独立、实值、中心化且 $Z_n\in L^2$。若
\[
 \sum_{n=1}^{\infty}\Var(Z_n)<\infty,
\]
则级数 $\sum_nZ_n$ 在 $L^2$ 中并且 a.s. 收敛。
\end{theorem}

\begin{proof}
令 $R_N=\sum_{n=1}^NZ_n$。独立性和中心化消去不同指标间的交叉项，所以对 $M>N$，
\[
 \E|R_M-R_N|^2=\sum_{n=N+1}^M\Var(Z_n)\longrightarrow0.
\]
$L^2$ 的完备性给出 $R_N$ 的 $L^2$ 极限。

为证明路径收敛，选取严格递增整数 $N_j$，使
\[
 \sum_{n>N_j}\Var(Z_n)\leq2^{-3j}.
\]
对每个 $j$，把 Kolmogorov 最大不等式用于有限块
$Z_{N_j+1},\dots,Z_{N_{j+1}}$，得到
\begin{align*}
 &\Prob\left(
 \max_{N_j<m\leq N_{j+1}}
 \left|R_m-R_{N_j}\right|>2^{-j}
 \right)\\
 &\hspace{30mm}\leq
 2^{2j}\sum_{n=N_j+1}^{N_{j+1}}\Var(Z_n)
 \leq2^{-j}.
\end{align*}
概率和有限，故 Borel--Cantelli I 表明：对 a.e. $\omega$，存在 $J(\omega)$，使对
$j\geq J(\omega)$，该块中的每个部分和都与块起点相差至多 $2^{-j}$。
若 $m>n\geq N_J$，设 $n,m$ 所在块的指标分别为 $j_0\leq j_1$。具体地，
\[
 R_m-R_n
 =-(R_n-R_{N_{j_0}})
  +\sum_{j=j_0}^{j_1-1}(R_{N_{j+1}}-R_{N_j})
  +(R_m-R_{N_{j_1}}),
\]
其中 $j_0=j_1$ 时中间的和为空和。两个端块增量与每个完整块增量都受相应的
$2^{-j}$ 控制，因而
\[
 |R_m(\omega)-R_n(\omega)|
 \leq2^{-j_0}+\sum_{j=j_0}^{j_1-1}2^{-j}+2^{-j_1}
 \leq4\,2^{-j_0}\leq4\,2^{-J}.
\]
让 $J\to\infty$，右端趋于零，所以 $(R_N(\omega))$ 是 Cauchy 数列并收敛。
记这个逐点极限为 $R_\infty$。前半部分已经给出某个 $R\in L^2$，使
$R_N\to R$ 于 $L^2$，从而依概率收敛；刚得到的 a.s. 收敛也蕴含
$R_N\xrightarrow{\Prob}R_\infty$。对任意 $\varepsilon>0$，
\[
 \Prob(|R-R_\infty|>\varepsilon)
 \leq
 \Prob(|R-R_N|>\varepsilon/2)
 +\Prob(|R_N-R_\infty|>\varepsilon/2)\longrightarrow0.
\]
因此 $R=R_\infty$ a.s.，这里的 $L^2$ 极限与路径极限确为同一个随机变量。
\end{proof}

作为抵消机制的一个例子，令 $\varepsilon_n$ 为独立 Rademacher 变量并取
$Z_n=\varepsilon_n/n$。方差级数 $\sum n^{-2}$ 收敛，故
$\sum\varepsilon_n/n$ a.s. 收敛；但
$\sum|\varepsilon_n|/n=\sum1/n$ 在每个样本点上发散。判据利用的是独立中心项的抵消，
不是绝对收敛的改写。

\begin{definition}[逐项截断]\label{def:5-4-2-2}
对 i.i.d. 序列 $(X_n)$，定义
\[
 Y_n=X_n\1_{\{|X_n|\leq n\}}.
\]
\end{definition}

\begin{lemma}[可积变量的逐项截断最终不变]\label{lemma:5-4-2-eventual-truncation}
若 $\E|X_1|<\infty$，则定义~\ref{def:5-4-2-2} 中的截断满足 $X_n=Y_n$ 最终 a.s.
\end{lemma}

\begin{proof}
若 $\E|X_1|<\infty$，则由 Tonelli 定理
\begin{align*}
 \sum_{n=1}^{\infty}\Prob(X_n\ne Y_n)
 &=\sum_{n=1}^{\infty}\Prob(|X_1|>n)
 =\E\sum_{n=1}^{\infty}\1_{\{|X_1|>n\}}\\
 &=\E\,\#\{n\in\N:n<|X_1|\}
 \leq\E|X_1|<\infty.
\end{align*}
因此 Borel--Cantelli I 保证 $X_n=Y_n$ 最终 a.s.。
\end{proof}

\begin{theorem}[Kolmogorov 强大数律]\label{theorem:5-4-2-4}
若 $(X_n)$ i.i.d. 且 $\E|X_1|<\infty$，则
\[
 \frac{S_n}{n}\longrightarrow\E X_1\qquad\text{a.s.}
\]
\end{theorem}

\begin{proof}
使用定义~\ref{def:5-4-2-2} 中的 $Y_n$。这些变量仍然独立。先验证随机级数判据的方差条件：
\begin{align*}
 \sum_{n=1}^{\infty}\frac{\Var(Y_n)}{n^2}
 &\leq\sum_{n=1}^{\infty}\frac{\E Y_n^2}{n^2}\\
 &=\E\left[|X_1|^2
       \sum_{\substack{n\geq1\\n\geq|X_1|}}\frac1{n^2}\right].
\end{align*}
存在绝对常数 $C$，使对每个 $r\geq0$，
\[
 r^2\sum_{\substack{n\geq1\\n\geq r}}\frac1{n^2}\leq Cr.
\]
事实上，$0\leq r<1$ 时用 $r^2\leq r$ 与 $\sum n^{-2}<\infty$；$r\geq1$ 时令
$N=\lceil r\rceil$，并用
$\sum_{n=N}^{\infty}n^{-2}\leq N^{-2}+\int_N^\infty x^{-2}\dd x\leq2/r$。
所以
\[
 \sum_{n=1}^{\infty}\frac{\Var(Y_n)}{n^2}
 \leq C\E|X_1|<\infty.
\]

令 $Z_n=(Y_n-\E Y_n)/n$。Kolmogorov 收敛判据说明
$\sum_nZ_n$ a.s. 收敛。对每条使该级数收敛的路径，把 Kronecker 引理用于
$a_n=n$、$x_n=Y_n-\E Y_n$，得到
\[
 \frac1n\sum_{k=1}^n(Y_k-\E Y_k)\longrightarrow0
 \qquad\text{a.s.}
\]
另一方面，
\[
 \E Y_n=\E\bigl[X_1\1_{\{|X_1|\leq n\}}\bigr]\longrightarrow\E X_1
\]
由支配收敛定理成立，故 \Cesaro{} 平均给
$n^{-1}\sum_{k\leq n}\E Y_k\to\E X_1$。

最后，\cref{lemma:5-4-2-eventual-truncation} 给出 $X_n=Y_n$ 最终 a.s.。在该概率一事件上，
$\sum_{k=1}^n(X_k-Y_k)$ 从某个时刻起是一个有限随机常数，除以 $n$ 后趋于零。
把三部分相加即得
\[
 \frac1n\sum_{k=1}^nX_k
 =\frac1n\sum_{k=1}^n(Y_k-\E Y_k)
  +\frac1n\sum_{k=1}^n\E Y_k
  +\frac1n\sum_{k=1}^n(X_k-Y_k)
 \longrightarrow\E X_1
\]
a.s.
\end{proof}

弱律与强律之间不是一个可以省略证明的箭头。弱律只控制每个固定时刻的坏事件；强律要求一条路径上
从某时刻起所有时刻都好。最大不等式把逐点控制升级为整块控制，截断再把罕见大跳交给
Borel--Cantelli。三种工具各自承担一段不同的逻辑。

因此经验均值、经验频率与独立抽样平均都由定理~\ref{theorem:5-4-2-4} 得到 a.s.
一致性。相关样本需要另行证明与其依赖结构相适应的极限定理；独立强律可以作为比较基准，却不能
直接用于相关样本。

\begin{exercise}[强律证明的关键步骤]
\begin{enumerate}
 \item 设 $(X_n)$ 独立、中心且 $\sum_n\Var(X_n)/n^2<\infty$。证明
 $n^{-1}\sum_{k\leq n}X_k\to0$ a.s.；
 \item 在有限方差强律的证明中，把二进时刻换成 $r^k$（固定 $r>1$），写出块最大估计；
 \item 构造一列 $X_n\xrightarrow{\Prob}0$ 但不 a.s. 收敛的变量，说明弱律形式本身不能推出强律。
\end{enumerate}
\end{exercise}

\subsection{特征函数、\Levy{} 连续性与中心极限定理}\label{subsec:5-4-3}
\index{特征函数}\index{Levy@\Levy{} 连续性定理}\index{中心极限定理}

大数律把 $S_n/n$ 的随机性压到一个常数。若要观察剩余波动，正确尺度通常是 $\sqrt n$。
此时直接计算归一化和的密度并不现实：每增加一项就要再做一次卷积，变量没有密度时这条路线甚至
无法起步。我们需要一种对每个概率分布都存在、又能把独立和变成乘积的观测。

矩母函数有乘法性质，却可能像 Pareto 例那样根本不存在。把实指数换成模长恒为一的振荡函数，
可积性问题消失了；代价是最后必须证明这种振荡观测足以唯一识别分布，并且其极限仍来自概率测度。
这两个缺口分别由唯一性定理与 \Levy{} 连续性定理填补。

de Moivre--Laplace 的二项近似早于现代测度论；\Levy{} 的贡献之一，是把这种计算组织成
普遍适用的特征函数连续性原理。现代证明必须同时交代两件事：紧性产生极限，唯一性识别极限。
把其中任何一件压缩掉，都会在质量逃逸的例子前失效。

\begin{definition}[特征函数]\label{def:5-4-3-1}
若 $X$ 是 $\R^d$ 值随机向量，其特征函数为
\[
 \varphi_X(t)=\E e^{i\langle t,X\rangle},\qquad t\in\R^d.
\]
若 $\mu$ 是 $\R^d$ 上的概率测度，也记
$\varphi_\mu(t)=\int e^{i\langle t,x\rangle}\mu(\dd x)$。
由于被积函数模长为一，特征函数对每个概率分布都有定义，并满足
$|\varphi_X(t)|\leq1$、$\varphi_X(0)=1$。
\end{definition}

\begin{definition}[正定函数]\label{def:5-4-3-2}
函数 $\varphi:\R^d\to\C$ 称为正定的，若对任意
$m\geq1$、$t_1,\dots,t_m\in\R^d$ 与 $c_1,\dots,c_m\in\C$，
\[
 \sum_{j,k=1}^mc_j\overline{c_k}\,\varphi(t_j-t_k)\geq0.
\]
这里左端按要求是非负实数。稍后证明每个特征函数都正定；正定连续函数何时反过来是特征函数，
属于 Bochner 定理的内容，这里不讨论该逆命题。
\end{definition}

\begin{proposition}[特征函数的基本性质]\label{proposition:5-4-3-1}
设 $X$ 为 $\R^d$ 值随机向量。
\begin{enumerate}
 \item $\varphi_X$ 在 $\R^d$ 上一致连续并且正定；
 \item 若 $X,Y$ 独立且同为 $\R^d$ 值，则
 $\varphi_{X+Y}=\varphi_X\varphi_Y$；
 \item 若 $A$ 是实 $m\times d$ 矩阵、$b\in\R^m$，则
 \[
  \varphi_{AX+b}(t)=e^{i\langle t,b\rangle}\varphi_X(A^Tt),
  \qquad t\in\R^m.
 \]
\end{enumerate}
\end{proposition}

\begin{proof}
对任意 $t,h\in\R^d$，
\[
 |\varphi_X(t+h)-\varphi_X(t)|
 \leq\E|e^{i\langle h,X\rangle}-1|.
\]
右端与 $t$ 无关，并由支配收敛定理在 $h\to0$ 时趋于零，所以连续性是一致的。
此外
\begin{align*}
 \sum_{j,k}c_j\overline{c_k}\varphi_X(t_j-t_k)
 &=\E\sum_{j,k}c_j\overline{c_k}
       e^{i\langle t_j-t_k,X\rangle}\\
 &=\E\left|\sum_jc_je^{i\langle t_j,X\rangle}\right|^2\geq0,
\end{align*}
这证明正定性。

若 $X,Y$ 独立，则期望因子化给
\[
 \varphi_{X+Y}(t)
 =\E\bigl(e^{i\langle t,X\rangle}e^{i\langle t,Y\rangle}\bigr)
 =\varphi_X(t)\varphi_Y(t).
\]
最后，内积恒等式
$\langle t,AX+b\rangle=\langle A^Tt,X\rangle+\langle t,b\rangle$
直接给出线性变换公式。
\end{proof}

第一项重要计算仍是高斯分布，但这次不能把实矩母函数公式中的 $t$ 形式地替换为 $it$；
那样做已经假定了尚未建立的复参数理论。下面从实变量积分重新计算。

\begin{example}[高斯特征函数]\label{ex:5-4-3-1}
令
\[
 F(t)=\frac1{\sqrt{2\pi}}\int_{\R}e^{itx}e^{-x^2/2}\dd x.
\]
对 $t$ 在任意有界区间变化时，导数被 $|x|e^{-x^2/2}$ 支配，故可在积分号下求导。
又 $xe^{-x^2/2}=-(e^{-x^2/2})'$。对 $R>0$ 在 $[-R,R]$ 上分部积分，
\begin{align*}
 \int_{-R}^{R}xe^{itx}e^{-x^2/2}\dd x
 &=-\int_{-R}^{R}e^{itx}\dd\bigl(e^{-x^2/2}\bigr)\\
 &=-\bigl[e^{itx}e^{-x^2/2}\bigr]_{-R}^{R}
   +it\int_{-R}^{R}e^{itx}e^{-x^2/2}\dd x.
\end{align*}
边界项的模不超过 $2e^{-R^2/2}$，因而趋于零；其余积分由
$e^{-x^2/2}$ 控制。令 $R\to\infty$ 得
\begin{align*}
 F'(t)
 &=\frac{i}{\sqrt{2\pi}}\int_{\R}xe^{itx}e^{-x^2/2}\dd x\\
 &=\frac{i}{\sqrt{2\pi}}\int_{\R}it e^{itx}e^{-x^2/2}\dd x
 =-tF(t).
\end{align*}
由 $F(0)=1$，实参数常微分方程的唯一解是
\[
 F(t)=e^{-t^2/2}.
\]

现在令 $Z=(Z_1,\dots,Z_m)$，其中坐标独立且均为 $N(0,1)$；对实
$d\times m$ 矩阵 $A$，暂把 $AZ$ 的分布记为 $N(0,\Sigma)$，其中
$\Sigma=AA^T$，允许 $\Sigma$ 奇异。独立乘法与线性变换公式给
\[
 \varphi_{AZ}(t)
 =\prod_{j=1}^m e^{-(A^Tt)_j^2/2}
 =\exp\left(-\frac12t^TAA^Tt\right)
 =\exp\left(-\frac12t^T\Sigma t\right).
\]
右端只含 $\Sigma$。定理~\ref{theorem:5-4-3-2} 将证明不同因子 $A$ 产生同一分布，
从而这个高斯分布记号不依赖所选分解。
\end{example}

如果两个概率的特征函数相同，我们知道它们对所有指数振荡测试相同；尚不能直接说“Fourier
反演”便结束，因为一般 Fourier 反演要到后章才建立。可用高斯核把两个测度先平滑，
而所需的高斯反演恒等式可以在本节用实变量常微分方程独立证明。

\begin{theorem}[特征函数唯一性定理]\label{theorem:5-4-3-2}
若 $\mu,\nu$ 是 $\R^d$ 上的概率测度，且
$\varphi_\mu(t)=\varphi_\nu(t)$ 对每个 $t\in\R^d$ 成立，则 $\mu=\nu$。
\end{theorem}

\begin{proof}
对 $\varepsilon>0$，记
\[
 g_\varepsilon(z)=(2\pi\varepsilon)^{-d/2}
 \exp\left(-\frac{|z|^2}{2\varepsilon}\right).
\]
我们先证明局部需要的高斯恒等式
\begin{equation}\label{eq:gaussian-local-inversion}
 g_\varepsilon(z)=\frac1{(2\pi)^d}
 \int_{\R^d}e^{-i\langle t,z\rangle}
 e^{-\varepsilon|t|^2/2}\dd t.
\end{equation}
把右端记为 $F_\varepsilon(z)$。高斯权重乘任意一次多项式仍可积，故可对每个
$z_j$ 在积分号下求导。对 $t_j$ 作实变量分部积分，并用
$\partial_{t_j}e^{-\varepsilon|t|^2/2}=-\varepsilon t_j
e^{-\varepsilon|t|^2/2}$，得到
\begin{align*}
 \partial_{z_j}F_\varepsilon(z)
 &=-\frac{i}{(2\pi)^d}\int t_j e^{-i\langle t,z\rangle}
 e^{-\varepsilon|t|^2/2}\dd t\\
 &=-\frac{z_j}{\varepsilon}F_\varepsilon(z).
\end{align*}
分部积分的边界项由高斯衰减趋于零。另一方面，多维实高斯积分的乘积分解给
\[
 F_\varepsilon(0)
 =\frac1{(2\pi)^d}\left(\frac{2\pi}{\varepsilon}\right)^{d/2}
 =(2\pi\varepsilon)^{-d/2}.
\]
沿每个坐标依次解上述实参数常微分方程，得到
$F_\varepsilon(z)=F_\varepsilon(0)e^{-|z|^2/(2\varepsilon)}$，即
\eqref{eq:gaussian-local-inversion}。

卷积 $\mu*g_\varepsilon$ 有密度
\begin{align*}
 p_{\mu,\varepsilon}(x)
 &=\int g_\varepsilon(x-y)\mu(\dd y)\\
 &=\frac1{(2\pi)^d}\int_{\R^d}
 e^{-i\langle t,x\rangle}e^{-\varepsilon|t|^2/2}
 \varphi_\mu(t)\dd t.
\end{align*}
Fubini 定理可用，因为高斯权重可积而 $|\varphi_\mu|\leq1$。
所以 $\varphi_\mu=\varphi_\nu$ 蕴含
$\mu*g_\varepsilon=\nu*g_\varepsilon$ 对每个 $\varepsilon>0$ 成立。

取一个分布为 $\mu$ 的随机向量 $X$，并在乘积空间上取与它独立的标准高斯向量 $Z$。
则 $X+\sqrt\varepsilon Z$ 的分布是 $\mu*g_\varepsilon$。对任意 $f\in C_b(\R^d)$，
$f(X+\sqrt\varepsilon Z)\to f(X)$ 逐点成立，且被 $\|f\|_\infty$ 支配，故
\[
 \int f\dd(\mu*g_\varepsilon)\longrightarrow\int f\dd\mu.
\]
同理，$\nu*g_\varepsilon\Rightarrow\nu$。由于两个平滑测度对每个 $\varepsilon$ 相同，
$\mu$ 与 $\nu$ 对所有 $C_b$ 函数积分相同，特别地对所有 $C_c$ 函数积分相同。
定理~\ref{theorem:3-1-1-4} 用 $C_c$ 测试恢复 Radon 开集质量，所以二者在所有开集、继而在
所有 Borel 集上相同。
\end{proof}

因此例~\ref{ex:5-4-3-1} 中若 $AA^T=BB^T$，则 $AZ$ 与 $BZ'$ 的特征函数相同，
唯一性定理保证它们同分布。奇异协方差也不例外；唯一性证明从未要求密度存在。

\begin{remark}[联合高斯分布的特征函数表示]\label{rem:5-4-3-joint-gaussian-recovery}
回到例~\ref{ex:5-2-3-1}中的联合高斯定义。设
$W=(W_1,\ldots,W_d)$ 满足该例的联合高斯定义，记
$m=\E W$ 且 $\Sigma=(\Cov(W_j,W_k))_{j,k}$。对任意 $t\in\R^d$，
$\langle t,W-m\rangle$ 或为常数零，或为均值零、方差
$t^T\Sigma t$ 的一维高斯变量。上面的一维计算因而给出
\[
 \varphi_W(t)=\exp\!\left(i\langle t,m\rangle-\frac12t^T\Sigma t\right).
\]
$\Sigma$ 半正定，故线性代数给出实矩阵 $A$ 使 $AA^T=\Sigma$。若
$Z$ 的坐标独立且均为 $N(0,1)$，则 $m+AZ$ 具有同一特征函数；由
定理~\ref{theorem:5-4-3-2}，
\[
 W\overset{\mathrm d}=m+AZ.
\]
因此“所有线性组合均为高斯分布”与上面的线性构造给出同一类分布。
若 $\Sigma$ 正定，在标准高斯乘积密度中作可逆线性换元，又得
\[
 p_W(x)=\frac{1}{(2\pi)^{d/2}(\det\Sigma)^{1/2}}
 \exp\!\left[-\frac12(x-m)^T\Sigma^{-1}(x-m)\right].
\]
$d=2$ 时这正是例~\ref{ex:5-2-3-1} 完成平方所用的联合密度。若
$\Sigma$ 奇异，分布集中于一个真仿射子空间，不应强行写成
$d$ 维 Lebesgue 密度；线性构造与特征函数仍然有效。
\end{remark}

特征函数比矩稳健得多。下一例用纯实变量卷积证明 Cauchy 分布的稳定性，同时避免使用留数定理。

\begin{example}[Cauchy 平均的稳定性]\label{ex:5-4-3-cauchy-stability}
对 $a>0$，令
\[
 f_a(x)=\frac{a}{\pi(a^2+x^2)}.
\]
我们证明 $f_a*f_b=f_{a+b}$。固定 $x$，置
\[
 \Delta=(x^2+(a+b)^2)(x^2+(a-b)^2).
\]
当 $\Delta\ne0$ 时，有部分分式分解
\begin{align*}
 \frac1{(y^2+b^2)((x-y)^2+a^2)}
 &=\frac{Ay+B}{y^2+b^2}
   +\frac{-A(y-x)+C}{(y-x)^2+a^2},
\end{align*}
其中
\[
 A=\frac{2x}{\Delta},\qquad
 B=\frac{x^2+a^2-b^2}{\Delta},\qquad
 C=\frac{x^2-a^2+b^2}{\Delta}.
\]
将上式两边乘以两个二次分母后，所需验证的多项式恒等式为
\[
 (Ay+B)\bigl((x-y)^2+a^2\bigr)
 +\bigl(-A(y-x)+C\bigr)(y^2+b^2)=1;
\]
左边的 $y^3,y^2,y$ 和常数项系数分别是
\begin{align*}
 &A-A=0,\\
 &-Ax+B+C,\\
 &A(x^2+a^2-b^2)-2Bx,\\
 &B(x^2+a^2)+(Ax+C)b^2.
\end{align*}
代入 $A,B,C$ 后，中间两项为
\[
 \frac{-2x^2+(x^2+a^2-b^2)+(x^2-a^2+b^2)}{\Delta}=0,
\]
以及
\[
 \frac{2x(x^2+a^2-b^2)-2x(x^2+a^2-b^2)}{\Delta}=0.
\]
常数项的分子则是
\begin{align*}
 &(x^2+a^2-b^2)(x^2+a^2)+(3x^2-a^2+b^2)b^2\\
 &=x^4+2(a^2+b^2)x^2+(a^2-b^2)^2
 =\Delta,
\end{align*}
故常数项等于一，多项式恒等式得证。
在线段 $[-R,R]$ 上积分时，两个线性分子项必须先合并。它们的原函数之和含有
\[
 \frac A2\log\frac{y^2+b^2}{(y-x)^2+a^2};
\]
先在两个端点作差、再令 $R\to\infty$，该对数贡献趋于零。不能把两个并不各自
Lebesgue 可积的线性分子项分开积分。常数项则给
\begin{align*}
 \int_{\R}\frac{\dd y}{(y^2+b^2)((x-y)^2+a^2)}
 &=\pi\left(\frac B b+\frac C a\right)\\
 &=\frac{\pi(a+b)}{ab\,[x^2+(a+b)^2]}.
\end{align*}
最后一步的代数化简是
\begin{align*}
 \frac Bb+\frac Ca
 &=\frac{a(x^2+a^2-b^2)+b(x^2-a^2+b^2)}{ab\Delta}\\
 &=\frac{(a+b)\bigl(x^2+(a-b)^2\bigr)}
 {ab\bigl(x^2+(a+b)^2\bigr)\bigl(x^2+(a-b)^2\bigr)}\\
 &=\frac{a+b}{ab\bigl(x^2+(a+b)^2\bigr)}.
\end{align*}
乘以 $ab/\pi^2$ 即得 $(f_a*f_b)(x)=f_{a+b}(x)$。唯一未被上述代数覆盖的情形是
$x=0$ 且 $a=b$；两边关于 $x$ 连续，故由 $x\to0$ 得到同一等式。
这里左边的连续性也可直接由 $f_a$ 的一致连续性估计：
\[
 |(f_a*f_b)(x+h)-(f_a*f_b)(x)|
 \leq\|\tau_hf_a-f_a\|_\infty\|f_b\|_1\longrightarrow0.
\]

所以独立的 Cauchy 变量相加时尺度参数相加。若 $X_1,\dots,X_n$ 独立且密度均为 $f_1$，
则 $S_n$ 的密度为 $f_n$。变量 $S_n/n$ 的密度为
\[
 y\longmapsto n f_n(ny)
 =\frac1{\pi(1+y^2)}=f_1(y).
\]
这证明 $S_n/n\overset{\mathrm d}=X_1$，也就完成了例~\ref{ex:5-4-2-2}所用的稳定性计算。
若另行证明 Fourier 恒等式 $\varphi_{f_a}(t)=e^{-a|t|}$，还可把卷积恒等式化成一次乘法；
本例不调用该恒等式，上面的推导完全建立在实变量积分上。
\end{example}

所有整数矩存在也未必能唯一识别分布；这不是一句抽象警告，而有一族可以逐项计算的密度。

\begin{remark}[所有整数矩仍可能不唯一决定分布]\label{rem:5-4-3-moment-indeterminate}
对 $x>0$ 令
\[
 p_0(x)=\frac1{x\sqrt{2\pi}}e^{-(\log x)^2/2},
 \qquad
 p_\theta(x)=p_0(x)\bigl[1+\theta\sin(2\pi\log x)\bigr],
 \quad |\theta|\leq1.
\]
每个 $p_\theta$ 非负。作换元 $y=\log x$。对整数 $n\geq0$ 和实数 $b$，配方并平移
$y=u+n$，再使用例~\ref{ex:5-4-3-1} 的高斯余弦积分，得到
\begin{align*}
 &\frac1{\sqrt{2\pi}}\int_{\R}
 e^{ny-y^2/2}\sin(by)\dd y\\
 &\qquad=e^{n^2/2}\frac1{\sqrt{2\pi}}
 \int_{\R}e^{-u^2/2}\bigl[\sin(bu)\cos(bn)+\cos(bu)\sin(bn)\bigr]\dd u\\
 &\qquad=e^{n^2/2}\left[0\cdot\cos(bn)+e^{-b^2/2}\sin(bn)\right]
 =e^{(n^2-b^2)/2}\sin(bn).
\end{align*}
其中第一项因被积函数为奇函数而为零。相同的配方还给出
\[
 \int_0^\infty x^np_0(x)\dd x
 =\frac1{\sqrt{2\pi}}\int_{\R}e^{ny-y^2/2}\dd y
 =e^{n^2/2}<\infty,
\]
而取 $n=0$ 时特别得到
\[
 \int_0^\infty p_0(x)\dd x
 =\frac1{\sqrt{2\pi}}\int_{\R}e^{-y^2/2}\dd y=1.
\]
取 $b=2\pi$，便有
\[
 \int_0^\infty x^np_0(x)\sin(2\pi\log x)\dd x
 =e^{(n^2-4\pi^2)/2}\sin(2\pi n)=0.
\]
结合上式，$n=0$ 的扰动积分为零说明每个 $p_\theta$ 都归一化；一般 $n$ 说明它们的
$n$ 阶矩都存在且等于 $e^{n^2/2}$。
然而当 $\theta\ne\theta'$ 时，两个密度在
$\sin(2\pi\log x)\ne0$ 的正测度集合上不同。因此“全部矩存在”本身不足以决定分布。
这不排除在额外矩增长条件下成立的矩唯一性定理。
\end{remark}


点态极限还剩一个危险：它可能不再是特征函数。最简单的失败发生在质量整体逃向无穷时。

\begin{example}[Bernoulli 与 Rademacher 的零点展开]\label{ex:5-4-3-2}
若 $R$ 为 Rademacher 变量，则 $\varphi_R(t)=\cos t$，所以
\[
 \varphi_R(t)=1-\frac{t^2}{2}+o(t^2)\qquad(t\to0).
\]
对独立副本 $R_i$，
\[
 \varphi_{n^{-1/2}\sum_{i=1}^nR_i}(t)
 =\cos(t/\sqrt n)^n\longrightarrow e^{-t^2/2}.
\]
更一般地，若 $B\sim\operatorname{Bernoulli}(p)$、$0<p<1$，则
$X=(B-p)/\sqrt{p(1-p)}$。记 $q=1-p$，则 $X$ 以概率 $p$ 取
$\sqrt{q/p}$，以概率 $q$ 取 $-\sqrt{p/q}$，所以
\begin{align*}
 \varphi_X(t)
 &=p e^{it\sqrt{q/p}}+q e^{-it\sqrt{p/q}}\\
 &=p\left(1+it\sqrt{\frac qp}-\frac{t^2q}{2p}+o(t^2)\right)
  +q\left(1-it\sqrt{\frac pq}-\frac{t^2p}{2q}+o(t^2)\right)\\
 &=1+it(\sqrt{pq}-\sqrt{pq})-\frac{t^2}{2}(q+p)+o(t^2)\\
 &=1-\frac{t^2}{2}+o(t^2).
\end{align*}
因而它与 Rademacher 变量具有同一个二阶主项；偏斜度只进入更高阶余项。
\end{example}


\begin{example}[零点处不连续的极限]\label{ex:5-4-3-3}
令 $\mu_n$ 为 $[-n,n]$ 上的均匀概率。直接积分给
\[
 \varphi_n(t)=
 \begin{cases}
  \dfrac{\sin(nt)}{nt},&t\ne0,\\[5pt]
  1,&t=0.
 \end{cases}
\]
对每个 $t\ne0$，$\varphi_n(t)\to0$，而在 $t=0$ 处始终等于一，所以点态极限在零点
不连续。另一方面，对固定 $R>0$ 与 $n\geq R$，
$\mu_n([-R,R])=R/n\to0$；任意紧集都包含在某个这样的区间中，因而没有一个紧集能统一抓住
这些概率的大部分质量。零点连续性将在下面的证明中恰好变成紧性。
\end{example}

\begin{theorem}[\Levy{} 连续性定理]\label{theorem:5-4-3-3}
设 $\mu_n,\mu$ 是 $\R^d$ 上的概率测度。
\begin{enumerate}
 \item 若 $\mu_n\Rightarrow\mu$，则
 $\varphi_{\mu_n}(t)\to\varphi_\mu(t)$ 对每个 $t\in\R^d$ 成立；
 \item 反之，若 $\varphi_{\mu_n}(t)\to\varphi(t)$ 对每个 $t$ 成立，且
 $\varphi$ 在 $0$ 处连续，则存在概率测度 $\mu$，使
 $\varphi=\varphi_\mu$ 且 $\mu_n\Rightarrow\mu$。
\end{enumerate}
\end{theorem}

\begin{proof}
若 $\mu_n\Rightarrow\mu$，对固定 $t$，函数
$x\mapsto e^{i\langle t,x\rangle}$ 有界连续，故弱收敛的定义直接给第一部分。

现在假设第二部分的条件。先证 $(\mu_n)$ 是紧族。先处理 $d=1$，并在规范空间上令
$X_n(x)=x$、$X_n\sim\mu_n$。对 $T>0$，Fubini 定理给
\begin{align*}
 \frac1{2T}\int_{-T}^T\Re\varphi_{\mu_n}(t)\dd t
 &=\E\left[\frac1{2T}\int_{-T}^T\cos(tX_n)\dd t\right]\\
 &=\E\left[\frac{\sin(TX_n)}{TX_n}\right],
\end{align*}
其中 $X_n=0$ 时把比值定义为一。若 $|X_n|\geq2/T$，则
$|\sin(TX_n)/(TX_n)|\leq1/2$，所以
\[
 \1_{\{|X_n|\geq2/T\}}
 \leq2\left(1-\frac{\sin(TX_n)}{TX_n}\right).
\]
取期望得到
\begin{equation}\label{eq:levy-tightness-estimate}
 \Prob(|X_n|\geq2/T)
 \leq2\left(1-\frac1{2T}\int_{-T}^T
                 \Re\varphi_{\mu_n}(t)\dd t\right).
\end{equation}
固定 $T$ 后，由 $|\varphi_{\mu_n}|\leq1$、点态收敛和支配收敛定理，右端收敛到
\[
 2\left(1-\frac1{2T}\int_{-T}^T\Re\varphi(t)\dd t\right).
\]
$\varphi(0)=\lim_n\varphi_{\mu_n}(0)=1$，且它在零点连续，所以这个量在
$T\downarrow0$ 时趋于零。给定
$\varepsilon>0$，先取 $T$ 足够小，再由上述收敛取 $N$，便能用区间
$[-2/T,2/T]$ 抓住所有 $n\geq N$ 的至少 $1-\varepsilon$ 质量。由
\cref{proposition:5-3-3-polish-radon}，剩余有限多个概率各自可由一个
有界闭区间抓住至少 $1-\varepsilon$ 的质量；取这些区间与前一区间的有限并，得到全列的紧性。

在 $\R^d$ 中，对每个坐标轴方向 $e_j$ 把同一估计应用于
$t\mapsto\varphi_{\mu_n}(te_j)$。给定 $\varepsilon>0$，利用 $\varphi$ 在零点的连续性，
可取同一个充分小的 $T>0$，使每个坐标的极限上界都小于 $\varepsilon/(2d)$；再取共同的
$N$，使对所有 $n\geq N$、所有 $1\leq j\leq d$ 都有
\[
 \mu_n\bigl(\{x:|x_j|\geq2/T\}\bigr)<\varepsilon/d.
\]
并集界说明立方体 $[-2/T,2/T]^d$ 抓住这些尾项至少 $1-\varepsilon$ 的质量。剩余有限多个
概率各自具有紧逼近，把相应有限多个紧集与该立方体作有限并，便得到对全列有效的紧集。因此一般维数的
$(\mu_n)$ 也是紧族。

由 Prokhorov 定理~\ref{theorem:5-3-3-2}，存在子列
$\mu_{n_k}\Rightarrow\mu$。第一部分说明
\[
 \varphi_\mu(t)=\lim_k\varphi_{\mu_{n_k}}(t)=\varphi(t),
\]
所以 $\varphi$ 确为一个概率的特征函数。任取原序列的任一子列；其紧性不变，故仍有进一步
弱收敛子列，而其极限的特征函数仍为 $\varphi$。唯一性定理
\ref{theorem:5-4-3-2} 说明所有这类子列极限都等于 $\mu$。

最后说明整列收敛。若 $\mu_n$ 不弱收敛到 $\mu$，由命题
\ref{proposition:5-3-3-weak-metric} 的有界 Lipschitz 度量化，存在 $\eta>0$ 和一个子列
满足 $d_{\mathrm{BL}}(\mu_{n_k},\mu)\geq\eta$。该子列又有进一步弱收敛子列，且其极限必须是
$\mu$，这与 $d_{\mathrm{BL}}$ 沿该进一步子列趋于零矛盾。因此 $\mu_n\Rightarrow\mu$。
\end{proof}

证明把两项工作严格分开：Prokhorov 从紧性产生候选极限，特征函数唯一性识别所有候选。
只有唯一性而没有存在性，不能排除质量逃逸；例~\ref{ex:5-4-3-3} 正是这个缺口的可计算模型。

现在回到独立和。二阶矩在零点给出特征函数的二阶展开，而 $\sqrt n$ 缩放恰好让
$n$ 个二次项累计成有限指数。

\begin{definition}[标准化和]\label{def:5-4-3-3}
若 $(X_i)$ i.i.d.，$\E X_1=\mu$ 且
$0<\Var(X_1)=\sigma^2<\infty$，定义
\[
 S_n=\sum_{i=1}^nX_i,
 \qquad
 Z_n=\frac{S_n-n\mu}{\sigma\sqrt n}.
\]
\end{definition}

\begin{theorem}[Lindeberg--\Levy{} 中心极限定理]\label{theorem:5-4-3-4}
若 $(X_i)$ i.i.d.，$\E X_1=\mu$ 且
$0<\Var(X_1)=\sigma^2<\infty$，则
\[
 \frac{S_n-n\mu}{\sigma\sqrt n}\Rightarrow N(0,1).
\]
\end{theorem}

\begin{proof}
用 $Y_i=(X_i-\mu)/\sigma$ 替换 $X_i$，只需处理
$\E Y_1=0$、$\E Y_1^2=1$ 的情形。记 $\varphi(u)=\E e^{iuY_1}$。
对实数 $v$，积分余项公式给
\[
 e^{iv}-1-iv=-\int_0^v(v-s)e^{is}\dd s,
 \qquad
 |e^{iv}-1-iv|\leq\frac{v^2}{2}.
\]
把
\[
 r(v)=
 \begin{cases}
 \dfrac{e^{iv}-1-iv}{v^2},&v\ne0,\\[5pt]
 -\dfrac12,&v=0
 \end{cases}
\]
看作连续函数。事实上，在余项积分中作 $s=uv$ 换元，对每个 $v\in\R$ 都有
\[
 r(v)=-\int_0^1(1-u)e^{iuv}\dd u.
\]
该表示式直接给出 $r(v)\to-1/2$ 当 $v\to0$，并且
$|r(v)|\le\int_0^1(1-u)\dd u=1/2$。于是
\[
 \frac{e^{iuY_1}-1-iuY_1}{u^2}=Y_1^2r(uY_1).
\]
右端被 $Y_1^2/2$ 支配，并在 $u\to0$ 时趋于 $-Y_1^2/2$。
由支配收敛定理及 $\E Y_1=0$、$\E Y_1^2=1$，
\[
 \varphi(u)=1-\frac{u^2}{2}+o(u^2).
\]

独立性给出，对每个固定 $t\in\R$，
\[
 \varphi_{n^{-1/2}\sum_{i=1}^nY_i}(t)
 =\varphi(t/\sqrt n)^n.
\]
令 $w_n=n(\varphi(t/\sqrt n)-1)$，则 $w_n\to-t^2/2$，所以右端为
$(1+w_n/n)^n$。这里不需要选择复对数分支。事实上，对任意复数列 $w_n\to w$，
二项式展开写成
\[
 \left(1+\frac{w_n}{n}\right)^n
 =\sum_{k=0}^{\infty}
 \1_{\{k\leq n\}}\frac{n(n-1)\cdots(n-k+1)}{n^k}
 \frac{w_n^k}{k!}.
\]
对每个固定 $k$，通项趋于 $w^k/k!$；当 $n$ 充分大时，绝对值被
$M^k/k!$ 支配，其中 $M>\sup_n|w_n|$。在计数测度上使用支配收敛定理，得到极限
$\sum_{k\geq0}w^k/k!=e^w$。因此
\[
 \varphi(t/\sqrt n)^n\longrightarrow e^{-t^2/2}.
\]
例~\ref{ex:5-4-3-1} 已证明右端是 $N(0,1)$ 的特征函数，并且它在零点连续。
\Levy{} 连续性定理~\ref{theorem:5-4-3-3} 于是给出所需弱收敛。
\end{proof}

\begin{center}
\begin{tikzpicture}[node distance=8mm and 8mm]
 \node[book node] (local) {单项零点展开\\$\varphi(u)=1-u^2/2+o(u^2)$};
 \node[book strong node,right=of local] (scale) {$\sqrt n$ 缩放\\$u=t/\sqrt n$};
 \node[book node,right=of scale] (power) {独立乘积\\$\varphi(t/\sqrt n)^n$};
 \node[book warm node,below=of power] (gauss) {高斯极限\\$e^{-t^2/2}$};
 \node[book strong node,left=of gauss] (levy) {\Levy{} 连续性\\$\Rightarrow$ 分布收敛};
 \draw[book accent arrow] (local)--(scale);
 \draw[book accent arrow] (scale)--(power);
 \draw[book accent arrow] (power)--(gauss);
 \draw[book accent arrow] (gauss)--(levy);
\end{tikzpicture}
\end{center}

这张图也标出定理的边界。有限二阶矩用于二阶展开；独立性用于乘法；零点连续性与
\Levy{} 定理用于把函数极限升级为概率极限。定理本身不给收敛速度；要得到定量误差率，还需要
三阶绝对矩等附加假设。多维中心极限定理可沿同一路线对每个 $t\in\R^d$ 展开
$\langle t,X\rangle$；其陈述还须明确向量矩条件、协方差矩阵以及退化协方差的情形。

本小节还澄清了矩母函数与特征函数的分工。前者在存在时能给强尾界，却可能对重尾分布失效；
后者始终存在并唯一决定分布，却不自动提供指数尾部。Fourier 反演、复解析延拓和定量局部极限定理
都需要额外工具；这里的唯一性、\Levy{} 定理与中心极限定理证明只使用实变量积分、紧性和
特征函数的基本性质。

\begin{exercise}[特征函数与极限]
\begin{enumerate}
 \item 设 $N$ 取值于 $\{0,1,2,\ldots\}$，且
 $\Prob(N=k)=e^{-\lambda}\lambda^k/k!$（$k\geq0$，$\lambda>0$）。计算 $N$ 的特征函数，
 并由独立乘法验证两个此类独立变量之和仍属于同一分布族，其参数为原参数之和；
 \item 证明若 $X_n\Rightarrow X$ 且 $\sup_n\E|X_n|^2<\infty$，则 $(\Law(X_n))$ 是紧族；
       再说明这个矩界本身不能推出二阶矩收敛；
 \item 对 $B\sim\operatorname{Bernoulli}(p)$ 直接展开标准化变量的特征函数到二阶，并与
       定理~\ref{theorem:5-4-3-4} 对照；
 \item 检查例~\ref{ex:5-4-3-3} 的点态极限为何不能是任何概率的特征函数。
\end{enumerate}
\end{exercise}

至此，独立和的三种尺度已经分开：不缩放时集中不等式给有限样本尾部，除以 $n$ 后大数律给
确定极限，除以 $\sqrt n$ 后中心极限定理保留高斯波动。下一章把单个随机变量换成时间索引族；
有限维分布仍是有限观测，而一致性条件将决定这些有限观测能否拼成路径空间上的分布。

